\newcommand{\PrefaceAncientVisionDiagram}{%
\begin{center}\begin{tikzpicture}[scale=.86,>=Latex]
  \begin{scope}[xshift=-3.8cm]
    \draw[very thick] (-2.1,-1.0)--(-1.45,-.35)--(-2.1,.3);
    \fill (-1.45,-.35) circle (.07);
    \draw[DeepBlue,->,very thick] (-1.4,-.35)--(.25,.45);
    \draw[DeepBlue,->,very thick] (-1.4,-.35)--(.4,-.35);
    \draw[DeepBlue,->,very thick] (-1.4,-.35)--(.25,-1.05);
    \draw[very thick,rounded corners=2pt] (.55,-1.1) rectangle (1.8,.65);
    \node[clearlabel,below,align=center] at (-.15,-1.45){古代“视觉触须”设想：\\由眼睛向物体发出微粒};
  \end{scope}
  \begin{scope}[xshift=3.6cm]
    \draw[very thick] (-2.2,-.9) rectangle (2.2,.9);
    \draw[DeepBlue,->,very thick] (-1.85,0)--(-.35,0);
    \fill (-.2,0) circle (.16);
    \foreach \a in {25,65,105,145,205,245,295,335}{
      \draw[Blue,->] (-.2,0)--++(\a:1.45);
    }
    \node[clearlabel,below,align=center] at (0,-1.35){量子态的传播不能简单理解为\\一颗经典台球沿唯一轨迹运动};
  \end{scope}
\end{tikzpicture}\figcap{图1.1—图1.2：早期视觉模型与量子粒子“弥散”的概念示意。}\end{center}}

\newcommand{\PrefaceParticleDiagram}{%
\begin{center}\begin{tikzpicture}[scale=.9,>=Latex]
  \shade[ball color=DeepBlue] (-2,0) circle (.45);
  \shade[ball color=Gold] (2,0) circle (.45);
  \draw[DeepBlue,very thick,->] (-1.45,0)--(-.35,0);
  \draw[Gold,very thick,->] (1.45,0)--(.35,0);
  \draw[Blue,very thick,->] (0,0)--++(55:1.75) node[clearlabel,above]{$\gamma$};
  \draw[Blue,very thick,->] (0,0)--++(235:1.75) node[clearlabel,below]{$\gamma$};
  \node[clearlabel,below] at (-2,-.65){电子 $e^-$};
  \node[clearlabel,below] at (2,-.65){正电子 $e^+$};
\end{tikzpicture}\figcap{图1.3：电子与正电子湮灭并产生光子的概念示意；能量和动量必须同时守恒。}\end{center}}

\newcommand{\EarthMoonScaleDiagram}{%
\begin{center}\begin{tikzpicture}[x=1cm,y=1cm,>=Latex]
  % 主图按同一线性比例绘制：中心距 10.5 cm，对应天体半径分别
  % 为 1.74 mm 和 0.48 mm。这样保留题图“真实尺度”的物理作用。
  \coordinate (E) at (-5.25,0);
  \coordinate (M) at (5.25,0);
  \shade[ball color=Blue!70] (E) circle (.174);
  \shade[ball color=TextGray!45] (M) circle (.048);
  \draw[<->,TextGray] (-5.25,-.45)--(5.25,-.45)
    node[midway,clearlabel,below=2pt]
      {平均中心距 $3.84\times10^5\,\mathrm{km}$};
  \draw[TextGray] (E)--++(-.42,.58);
  \node[clearlabel,above right,align=left] at (-5.72,.55)
    {地球 $R_E\approx6.37\times10^3\,\mathrm{km}$\\$M_E\approx5.97\times10^{24}\,\mathrm{kg}$};
  \draw[TextGray] (M)--++(.35,.58);
  \node[clearlabel,above left,align=right] at (5.67,.55)
    {月球 $R_M\approx1.74\times10^3\,\mathrm{km}$\\$M_M\approx7.35\times10^{22}\,\mathrm{kg}$};
  \node[clearlabel,below] at (0,-1.05)
    {半径与中心距采用同一比例；$M_M/M_E\approx0.0123$};
\end{tikzpicture}\figcap{图2.1：月球、地球及月地距离的真实比例。研究月球绕地球公转时可把月球近似为质点；研究着陆时则不能。}\end{center}}

\newcommand{\SprintGraph}{%
\begin{center}\begin{tikzpicture}
\begin{axis}[
  width=.78\linewidth,height=6.0cm,
  axis lines=left,xlabel={$t/\mathrm{s}$},ylabel={$x/\mathrm{m}$},
  xmin=0,xmax=12.5,ymin=0,ymax=105,
  xtick={0,...,12},ytick={0,20,...,100},
  grid=major,grid style={draw=RuleGray},
  every axis label/.style={fill=none},tick label style={font=\small}
]
\addplot[kcurve,mark=*,mark size=1.3pt] coordinates {
  (0,0)(1,1)(2,4)(3,9)(4,16)(5,25)(6,35)(7,45)(8,55)(9,65)(10,75)(11,85)(12,100)
};
\end{axis}
\end{tikzpicture}\figcap{图2.2：短跑选手位置随时间变化。曲线越陡，位置变化越快。}\end{center}}

\newcommand{\PositionSlopeDiagram}{%
\begin{center}\begin{tikzpicture}[scale=.94,>=Latex]
  \foreach \dx/\dy/\mode in {-3/0/1,3/0/2,0/-5.2/3}{
    \begin{scope}[shift={(\dx,\dy)}]
      \draw[kaxis] (-1.8,0)--(1.9,0) node[right,clearlabel]{$t$};
      \draw[kaxis] (-1.55,-1.35)--(-1.55,1.55) node[above,clearlabel]{$x$};
      \ifnum\mode=1
        \draw[kcurve] (-1.35,-.9)--(1.55,1.05);
        \node[clearlabel,below,align=center] at (0,-1.55){匀速运动\\直线斜率恒定};
      \else\ifnum\mode=2
        \draw[kcurve] plot[smooth] coordinates{(-1.35,.12)(-.7,.35)(0,.64)(.7,1.0)(1.55,1.27)};
        \coordinate (A) at (-.7,.35);\coordinate (B) at (.7,1.0);
        \fill (A) circle (.04);
        \fill (B) circle (.04);
        \draw[secant] (A)--(B);
        \draw[kguide] (-.7,0)--(A) (.7,0)--(B);
        \draw[kguide] (-1.55,.35)--(A) (-1.55,1.0)--(B);
        \draw[<->,TextGray] (-.7,-.45)--(.7,-.45)
          node[midway,clearlabel,below=1pt]{$\Delta t$};
        \draw[<->,TextGray] (1.18,.35)--(1.18,1.0)
          node[midway,clearlabel,right=1pt]{$\Delta x$};
        \draw[TextGray] (A)--(.38,.35);
        \draw[TextGray] (-.30,.35) arc[start angle=0,end angle=24.90,radius=.40];
        \node[clearlabel] at (.20,.58){$\alpha$};
        \node[clearlabel,above left=2pt] at (A) {$A$};
        \node[clearlabel,above right=2pt] at (B) {$B$};
        \node[clearlabel,below=1pt] at (-.7,0) {$t_1$};
        \node[clearlabel,below=1pt] at (.7,0) {$t_2$};
        \node[clearlabel,left=1pt] at (-1.55,.35) {$x$};
        \node[clearlabel,left=1pt] at (-1.55,1.0) {$x+\Delta x$};
        \node[clearlabel,below,align=center] at (0,-1.55){割线斜率：平均速度\\$\Delta t=t_2-t_1$};
      \else
        \draw[kcurve,domain=-1.35:1.5,samples=80] plot (\x,{.12+.65*\x-.1*\x*\x});
        \coordinate (A) at (0,.12);\fill (A) circle (.04);
        \draw[tangent] (-1.0,-.53)--(1.1,.835);
        \node[clearlabel,below right=3pt] at (A) {$A$};
        \node[clearlabel,below,align=center] at (0,-1.55){切线斜率\\瞬时速度};
      \fi\fi
    \end{scope}
  }
\end{tikzpicture}\figcap{图2.3：$x$-$t$ 图像中平均速度与瞬时速度的几何意义。}\end{center}}

\newcommand{\VelocityAreaDiagram}{%
\begin{center}\begin{tikzpicture}[scale=.82,>=Latex]
  \foreach \dx/\mode in {-4.7/1,0/2,4.7/3}{
    \begin{scope}[xshift=\dx cm]
      \draw[kaxis] (-1.8,0)--(1.9,0) node[right,clearlabel]{$t$};
      \draw[kaxis] (-1.55,-.15)--(-1.55,1.65) node[above,clearlabel]{$v$};
      \ifnum\mode=1
        \fill[Blue!12] (-1.1,0) rectangle (1.3,.9);
        \draw[kcurve] (-1.25,.9)--(1.5,.9);
        \draw[kguide] (-1.1,0)--(-1.1,.9) (1.3,0)--(1.3,.9);
        \node[clearlabel,below=2pt] at (-1.1,0){$A$};
        \node[clearlabel,below=2pt] at (1.3,0){$B$};
        \node[clearlabel,above left=1pt] at (-1.1,.9){$C$};
        \node[clearlabel,above right=1pt] at (1.3,.9){$D$};
        \node[clearlabel,below] at (0,-.55){匀速：矩形面积};
      \else
        \path[fill=Blue!12] (-1.1,0)--
          plot[domain=-1.1:1.3,samples=60]
            ({\x},{0.55+0.75*(\x+1.1)-0.22*(\x+1.1)^2})--(1.3,0)--cycle;
        \draw[kcurve]
          plot[domain=-1.1:1.3,samples=60]
            ({\x},{0.55+0.75*(\x+1.1)-0.22*(\x+1.1)^2});
        \draw[kguide] (-1.1,0)--(-1.1,.55) (1.3,0)--(1.3,1.08);
        \node[clearlabel,below=2pt] at (-1.1,0){$A$};
        \node[clearlabel,below=2pt] at (1.3,0){$B$};
        \node[clearlabel,above left=1pt] at (-1.1,.55){$C$};
        \node[clearlabel,above right=1pt] at (1.3,1.08){$D$};
        \ifnum\mode=2
          \foreach \x in {-1.05,-.85,...,1.25}{
            \draw[Blue!65] (\x,0)--
              (\x,{0.55+0.75*(\x+1.1)-0.22*(\x+1.1)^2});}
          \node[clearlabel,below] at (0,-.55){窄矩形求和};
        \else
          \node[clearlabel,below] at (0,-.55){极限：曲线下面积};
        \fi
      \fi
    \end{scope}
  }
\end{tikzpicture}\figcap{图2.4：$v$-$t$ 图像下的有向面积等于位移；若求路程，应对 $|v|$ 积分。}\end{center}}

\newcommand{\ExampleTwoSixGraphs}{%
\begin{center}\begin{tikzpicture}[scale=.78,>=Latex]
  \foreach \dx/\mode in {-4.6/1,0/2,4.6/3}{
    \begin{scope}[xshift=\dx cm]
      \draw[kaxis] (-1.65,0)--(1.75,0) node[right,clearlabel]{$t$};
      \draw[kaxis] (-1.45,-.55)--(-1.45,2.05) node[above,clearlabel]{$x$};
      \ifnum\mode=1
        \draw[kcurve] (-1.35,0)--(0,1.35)--(1.35,0);
        \draw[vector] (-.92,.43)--(-.42,.93);
        \draw[vector] (.42,.93)--(.92,.43);
        \node[clearlabel,below] at (0,-.62){先正向、后反向};
      \else\ifnum\mode=2
        \draw[kcurve,domain=-1.35:1.30,samples=80] plot (\x,{.25*(\x+1.35)^2});
        \draw[vector] (-.65,.12)--(-.15,.34);
        \draw[vector] (.35,.73)--(.82,1.25);
        \node[clearlabel,below] at (0,-.62){斜率为正且增大};
      \else
        \draw[kcurve,domain=-1.35:1.35,samples=100] plot (\x,{1.45*sqrt(max(0,1-(\x/1.35)^2))});
        \draw[vector] (-.85,.99)--(-.35,1.34);
        \draw[vector] (.35,1.34)--(.85,.99);
        \node[clearlabel,below] at (0,-.62){斜率由正变为负};
      \fi\fi
    \end{scope}}
\end{tikzpicture}\figcap{例2.6：三幅 $x$-$t$ 图的切线斜率分别给出速度的方向与变化趋势。}\end{center}}

\newcommand{\ExampleTwoSevenGraphs}{%
\begin{center}\begin{tikzpicture}[scale=.9,>=Latex]
  \begin{scope}[shift={(-4.9,-1.1)}]
    \draw[kaxis] (0,0)--(2.75,0) node[right,clearlabel]{$t$};
    \draw[kaxis] (0,-1.55)--(0,1.85) node[above,clearlabel]{$v$};
    \draw[kcurve] (0,1.3)--(2.4,-1.3);
    \node[clearlabel,left] at (0,1.3){$v_0$};
    \node[clearlabel,left] at (0,-1.3){$-v_0$};
    \draw[kguide] (0,-1.3)--(2.4,-1.3);
    \node[clearlabel,below left] at (0,0){$O$};
    \node[clearlabel,above right=2pt] at (1.2,0){$T/2$};
    \node[clearlabel,above] at (2.4,0){$T$};
    \node[clearlabel,below] at (1.2,-1.7){已知速度图};
  \end{scope}
  \foreach \dx/\dy/\mode/\tag in {-.7/0/1/A,3.3/0/2/B,-.7/-3.3/3/C,3.3/-3.3/4/D}{
    \begin{scope}[shift={(\dx,\dy)}]
      \draw[kaxis] (0,0)--(2.7,0) node[right,clearlabel]{$t$};
      \draw[kaxis] (0,-.25)--(0,1.65) node[above,clearlabel]{$x$};
      \ifnum\mode=1
        \draw[kcurve] (0,0)--(1.2,1.15)--(2.4,0);
      \else\ifnum\mode=2
        \draw[kcurve,domain=0:1,samples=70]
          plot ({2.4*\x},{4.6*\x*(1-\x)});
      \else\ifnum\mode=3
        \draw[kcurve,domain=0:1,samples=70] plot ({2.4*\x},{1.15*(2*\x-\x*\x)});
      \else
        \draw[kcurve] (0,1.15)--(1.2,0)--(2.4,.85);
      \fi\fi\fi
      \node[clearlabel,below left] at (0,0){$O$};
      \node[clearlabel,below=3pt] at (1.2,0){$T/2$};
      \node[clearlabel,below=3pt] at (2.4,0){$T$};
      \node[clearlabel,below] at (1.2,-.65){$\tag$};
    \end{scope}}
\end{tikzpicture}\figcap{例2.7：左图为给定的 $v$-$t$ 图，右侧四图为待判断的位移图；正确图必须连续且斜率等于给定速度。}\end{center}}

\newcommand{\ExampleTwoEightGraphs}{%
\begin{center}\begin{tikzpicture}[scale=.78,>=Latex]
  \foreach \dx/\mode in {-4.4/1,0/2,4.4/3}{
    \begin{scope}[xshift=\dx cm]
      \draw[kaxis] (-1.55,0)--(1.75,0) node[right,clearlabel]{时间};
      \draw[kaxis] (-1.35,-.25)--(-1.35,1.85) node[above,clearlabel]{$v$};
      \ifnum\mode=1
        \fill[Blue!12] (-1.25,0)--(1.25,0)--(1.25,1.35)--cycle;
        \draw[kcurve] (-1.25,0)--(1.25,1.35);
        \draw[kguide] (1.25,0)--(1.25,1.35);
        \draw[kguide] (-1.35,1.35)--(1.25,1.35);
        \node[clearlabel,left] at (-1.35,1.35){$v$};
      \else\ifnum\mode=2
        \fill[Blue!12] (-1.25,0)--(-1.25,.55)--(1.25,1.35)--(1.25,0)--cycle;
        \draw[kcurve] (-1.25,.55)--(1.25,1.35);
        \draw[kguide] (1.25,0)--(1.25,1.35);
        \draw[kguide] (-1.35,.55)--(-1.25,.55) (-1.35,1.35)--(1.25,1.35);
        \node[clearlabel,left] at (-1.35,.55){$v_0$};
        \node[clearlabel,left] at (-1.35,1.35){$v_1$};
      \else
        \fill[Blue!12] (-1.25,0)--(-.42,1.35)--(1.25,0)--cycle;
        \draw[kcurve] (-1.25,0)--(-.42,1.35)--(1.25,0);
        \draw[kguide] (-.42,0)--(-.42,1.35);
        \draw[kguide] (-1.35,1.35)--(-.42,1.35);
        \node[clearlabel,left] at (-1.35,1.35){$v_0$};
        \node[clearlabel,below] at (-.42,0){$t/3$};
      \fi\fi
      \node[clearlabel,below] at (1.25,0){$t$};
    \end{scope}}
\end{tikzpicture}\figcap{例2.8：总路程等于三幅 $v$-$t$ 图中阴影面积；梯形或三角形面积直接给出结果。}\end{center}}

\newcommand{\FreeFallMeetingDiagram}{%
\begin{center}\begin{tikzpicture}[scale=.9,>=Latex]
  \draw[very thick] (-1.4,0)--(1.4,0);
  \node[point,label={[clearlabel]below right=3pt:$B$}] (B) at (0,0) {};
  \node[point,label={[clearlabel]left:$A$}] (A) at (0,3.2) {};
  \draw[<->,TextGray] (.8,0)--(.8,3.2) node[midway,clearlabel,right]{$h$};
  \draw[DeepBlue,very thick,->] (B)--(0,1.15) node[midway,clearlabel,left]{$v_0$};
  \draw[Blue,very thick,->] (A)--(0,2.45) node[midway,clearlabel,left]{$g$};
\end{tikzpicture}\figcap{例2.17解图：两球具有相同的重力加速度，因此相对加速度为零，初始间距以相对速度 $v_0$ 关闭。}\end{center}}

\newcommand{\BounceDiagram}{%
\begin{center}\begin{tikzpicture}[scale=.85,>=Latex]
  \draw[very thick] (-1.7,0)--(1.7,0);
  \fill (0,2.7) circle (.12);
  \fill (0,.12) circle (.12);
  \draw[kguide] (0,.27)--(0,2.55);
  \draw[<->,TextGray] (.88,0)--(.88,2.7)
    node[midway,clearlabel,right=2pt]{$H$};
  \draw[DeepBlue,very thick,->] (-.35,1.0)--(-.35,.25) node[midway,clearlabel,left]{$u$};
  \draw[Blue,very thick,->] (.35,.25)--(.35,1.0) node[midway,clearlabel,right]{$eu$};
\end{tikzpicture}\figcap{例2.19解图：每次反弹速度按 $e$ 成等比衰减，飞行时间按 $e$ 衰减，路程按 $e^2$ 衰减。}\end{center}}

\newcommand{\CoordinateSystemsDiagram}{%
\begin{center}\begin{tikzpicture}[scale=.9,>=Latex]
  \begin{scope}[xshift=-3.2cm]
    \draw[kaxis] (-2,0)--(2.2,0) node[right,clearlabel]{$x$};
    \draw[kaxis] (0,-1.7)--(0,1.9) node[above,clearlabel]{$y$};
    \coordinate (P) at (1.35,1.15);\node[point,label={[clearlabel]above:$P(x,y)$}] at (P) {};
    \draw[kguide] (P)--(1.35,0) (P)--(0,1.15);
    \node[below left] at (0,0){$O$};
  \end{scope}
  \begin{scope}[xshift=3.2cm]
    \draw[kaxis] (0,0)--(2.3,0) node[right,clearlabel]{$N$};
    \draw[kaxis] (0,0)--(0,1.9);
    \coordinate (P) at (1.45,1.25);\draw[vector] (0,0)--(P) node[midway,clearlabel,above left]{$r$};
    \node[point,label={[clearlabel]above:$P(r,\theta)$}] at (P) {};
    \draw[->,TextGray] (.65,0) arc[start angle=0,end angle=41,radius=.65];
    \node[clearlabel] at (.82,.28){$\theta$};
    \node[below left] at (0,0){$O$};
  \end{scope}
\end{tikzpicture}\figcap{图2.5：平面直角坐标系与平面极坐标系。}\end{center}}

\newcommand{\PathChoicesDiagram}{%
\begin{center}\begin{minipage}{.96\linewidth}\centering
\begin{tikzpicture}[scale=.9,>=Latex]
 \coordinate (A) at (0,0);\coordinate (B) at (8,1.5);
 \draw[kcurve] (A)--(B) node[pos=.72,above=3pt]{$P_1$};
 \draw[kcurve] (A)..controls (1,2.9) and (5.9,3.0)..(B) node[pos=.46,above=3pt]{$P_2$};
 \draw[kcurve] (A)..controls (1.3,1.4) and (2.1,1.65)..(3.2,.25)
 ..controls (5,-1.0) and (7.7,0)..(B);
 \draw[kcurve] (A)..controls (-2.5,4) and (3.6,4.3)..(3.0,-.1)
 ..controls (2.8,-1.8) and (8.9,-1.1)..(B);
 \draw[kcurve] (A)..controls (3.6,-2.4) and (8.8,-2.0)..(B);
 \node at (5,.3){$P_3$};\node at (.8,3.35){$P_4$};\node at (5.5,-1.85){$P_5$};
 \node[point,label=below left:$A$] at (A){};\node[point,label=above right:$B$] at (B){};
\end{tikzpicture}
\figcap{图2.6：五条路径 $P_1,\ldots,P_5$ 均从 $A$ 到 $B$；路径不同，位移相同。}
\end{minipage}\end{center}}

\newcommand{\DisplacementFanDiagram}{%
\begin{center}\begin{minipage}{.98\linewidth}\centering
\begin{tikzpicture}[scale=.95,>=Latex]
 \coordinate (A) at (0,0);
 \foreach \i in {1,...,10}{
  \pgfmathsetmacro{\xx}{.95*\i}\pgfmathsetmacro{\yy}{.8*\xx-.06*\xx*\xx}
  \draw[DeepBlue!60,->,thin] (A)--(\xx,\yy);
 }
 \draw[kcurve,domain=0:9.75,samples=100] plot (\x,{.8*\x-.06*\x*\x});
 \foreach \i in {1,...,10}{
  \pgfmathsetmacro{\xx}{.95*\i}\pgfmathsetmacro{\yy}{.8*\xx-.06*\xx*\xx}
  \node[point] at (\xx,\yy){};\node[above=5pt,font=\small] at (\xx,\yy){$B_{\i}$};
 }
 \node[point,label=below left:$A$] at (A){};
\end{tikzpicture}
\figcap{图2.7：粗曲线是实际路径；从 $A$ 指向各 $B_i$ 的直箭头是对应位移。}
\end{minipage}\end{center}}

\newcommand{\VectorOperationsDiagram}{%
\begin{center}\begin{tikzpicture}[scale=.82,>=Latex]
  \begin{scope}[xshift=-4.4cm]
    \draw[vector] (0,0)--(1.5,.35) node[midway,clearlabel,below]{$\vect A$};
    \draw[vector] (1.5,.35)--(2.15,1.55) node[midway,clearlabel,right]{$\vect B$};
    \draw[vector] (0,0)--(2.15,1.55) node[midway,clearlabel,above left]{$\vect A+\vect B$};
  \end{scope}
  \begin{scope}
    \draw[vector] (0,0)--(1.55,.4) node[midway,clearlabel,below]{$\vect A$};
    \draw[vector] (0,0)--(.55,1.55) node[midway,clearlabel,left]{$\vect B$};
    \draw[vector] (.55,1.55)--(1.55,.4) node[midway,clearlabel,above right]{$\vect A-\vect B$};
  \end{scope}
  \begin{scope}[xshift=4.4cm]
    \draw[vector] (0,0)--(1.7,.35) node[midway,clearlabel,below]{$\vect A$};
    \draw[vector] (0,0)--(.75,1.45) node[midway,clearlabel,left]{$\vect B$};
    \draw[->,TextGray] (.65,.13) arc[start angle=12,end angle=63,radius=.67];
    \node[clearlabel] at (.9,.75){$\theta$};
    \node[clearlabel,below] at (.85,-.55){$\vect A\cdot\vect B=AB\cos\theta$};
  \end{scope}
\end{tikzpicture}\figcap{图2.8：矢量加法、减法与内积。}\end{center}}

\newcommand{\VectorComponentsDiagram}{%
\begin{center}\begin{tikzpicture}[scale=.9,>=Latex]
  \draw[kaxis] (-.4,0)--(4.4,0) node[right,clearlabel]{$x$};
  \draw[kaxis] (0,-.35)--(0,3.0) node[above,clearlabel]{$y$};
  \coordinate (P) at (3.4,2.2);\draw[vector] (0,0)--(P) node[midway,clearlabel,above left]{$\vect A$};
  \draw[kguide] (P)--(3.4,0) node[clearlabel,below]{$A_x$};
  \draw[kguide] (P)--(0,2.2) node[clearlabel,left]{$A_y$};
\end{tikzpicture}\figcap{图2.9：矢量平移到原点后，终点坐标就是矢量分量。}\end{center}}

\newcommand{\CurveVelocityDiagram}{%
\begin{center}\begin{minipage}{.95\linewidth}\centering
\begin{tikzpicture}[>=Latex]
 \draw[kcurve,domain=-3:3,samples=100] plot (\x,{.09*\x^3-.55*\x});
 \foreach \xx in {-2.3,-.9,.8,2.2}{
  \pgfmathsetmacro{\yy}{.09*\xx^3-.55*\xx}
  \pgfmathsetmacro{\slope}{.27*(\xx)*(\xx)-.55}\pgfmathsetmacro{\dx}{1/sqrt(1+(\slope)*(\slope))}
  \node[point] at (\xx,\yy){};\draw[vector] (\xx,\yy)--++(\dx,{\slope*\dx});
 }
\end{tikzpicture}
\figcap{图2.10：每个速度箭头都从所在位置出发，并沿该点的轨迹切线。}
\end{minipage}\end{center}}

\newcommand{\AccelerationVectorDiagram}{%
\begin{center}\begin{tikzpicture}[scale=.9,>=Latex]
  \coordinate (O) at (0,0);\coordinate (A) at (2.3,1.6);\coordinate (B) at (3.2,.75);
  \draw[vector] (O)--(A) node[midway,clearlabel,above left]{$\vect v(t)$};
  \draw[vector] (O)--(B) node[midway,clearlabel,below=7pt]{$\vect v(t+\Delta t)$};
  \draw[vector] (A)--(B) node[midway,clearlabel,above right]{$\Delta\vect v$};
  \node[clearlabel,right] at (3.65,1.2){$\displaystyle \vect a=\lim_{\Delta t\to0}\frac{\Delta\vect v}{\Delta t}$};
\end{tikzpicture}\figcap{图2.11：速度矢量首端重合，终端之差给出 $\Delta\vect v$。}\end{center}}

\newcommand{\RescueRouteDiagram}{%
\begin{center}\begin{minipage}{.96\linewidth}\centering
\begin{tikzpicture}[>=Latex]
 \coordinate (A) at (-2.8,2.3);\coordinate (C) at (.2,0);\coordinate (B) at (1.7,-1.9);
 \coordinate (U) at (.2,1.4);\coordinate (D) at (.2,-1.4);
 \draw[thick] (-3.5,0)--(3.6,0) node[right]{岸线};
 \draw[kguide] (A)--(-2.8,0) (B)--(1.7,0) (U)--(D);
 \draw[vector] (A)--(C) node[pos=.35,above right=2pt]{$v_1$};
 \draw[vector] (C)--(B) node[pos=.45,above right=2pt]{$v_2$};
 \draw pic[draw=TextGray,angle radius=7mm] {angle=U--C--A};
 \draw pic[draw=TextGray,angle radius=7mm] {angle=D--C--B};
 \node at (-.24,1.03){$\theta_1$};\node at (.55,-1.1){$\theta_2$};
 \draw[<->,TextGray] (-3.3,0)--(-3.3,2.3) node[midway,left]{$d_1$};
 \draw[<->,TextGray] (2.15,0)--(2.15,-1.9) node[midway,right]{$d_2$};
 \draw[<->,TextGray] (-2.8,-2.5)--(.2,-2.5) node[midway,below]{$x$};
 \draw[<->,TextGray] (.2,-2.5)--(1.7,-2.5) node[midway,below]{$L-x$};
 \node[point,label=above left:$A$] at (A){};\node[point,label=below right:$B$] at (B){};
 \node[point,label=above right:$C$] at (C){};
\end{tikzpicture}
\figcap{例2.31：入水点为 $C$；$\theta_1,\theta_2$ 从岸线法线量起，两垂足间距为 $L$。}
\end{minipage}\end{center}}

\newcommand{\FootballInterceptionDiagram}{%
\begin{center}\begin{minipage}{.96\linewidth}\centering
\begin{tikzpicture}[scale=1.1,>=Latex]
 \coordinate (A) at (0,0);\coordinate (B) at (4,2);\coordinate (C) at (5,0);\coordinate (X) at (6,0);
 \draw[kaxis] (A)--(X) node[right]{$X$};
 \draw[thick] (A)--(B) node[midway,above left]{$d$};\draw[kguide] (B)--(C);
 \draw[vector] (B)--($(B)!.43!(C)$) node[midway,right=5pt]{$\vect v_p$};
 \draw[vector] (A)--(1.25,0) node[midway,below=4pt]{$\vect v$};
 \draw pic[draw=TextGray,angle radius=8mm] {angle=X--A--B};\node at (1.03,.27){$\theta$};
 \draw pic[draw=TextGray,angle radius=4mm] {right angle=A--B--C};
 \foreach \p/\pos in {A/below left,B/above,C/below}\node[point,label=\pos:$\p$] at (\p){};
\end{tikzpicture}
\figcap{例2.29：最低奔跑速率对应 $BC\perp AB$；队员速度沿 $BC$，球的速度沿 $AC$。}
\end{minipage}\end{center}}

\newcommand{\TwoParticlesMinimumDistanceDiagram}{%
\begin{center}\begin{tikzpicture}[scale=.82,>=Latex]
  \coordinate (P) at (-3.7,-.75);\coordinate (Q) at (2.25,-.75);\coordinate (R) at (-.3,2.25);
  \draw[very thick] (P)--(Q);
  \draw[kguide] (Q)--(R);
  \draw[vector] (-3.25,-.75)--(-1.8,-.75) node[midway,clearlabel,above]{$\vect v_1$};
  \draw[vector] (1.93,-.38)--(.95,.78) node[midway,clearlabel,right]{$\vect v_2$};
  \draw pic[draw=TextGray,->,angle radius=8mm,"$\alpha$"{clearlabel}] {angle=R--Q--P};
  \draw[<->,TextGray] (P)--(Q) node[midway,clearlabel,below=3pt]{$l$};
  \node[point,label={[clearlabel]left:$A,P$}] at (P) {};
  \node[point,label={[clearlabel]right:$B,Q$}] at (Q) {};
  \node[clearlabel,above] at (R){$R$};
\end{tikzpicture}\figcap{例2.30：初始连线为 $PQ$；$A$ 沿 $PQ$，$B$ 沿 $QR$，最近距离由相对速度直线确定。}\end{center}}

\newcommand{\ProjectileVectorsDiagram}{%
\begin{center}\begin{minipage}{.98\linewidth}\centering
\begin{tikzpicture}[scale=.9,>=Latex]
 \begin{scope}[xshift=-3.6cm]
  \coordinate (O) at (0,0);\coordinate (A) at (2,3);\coordinate (B) at (2,-1);
  \draw[vector] (O)--(A) node[midway,above left]{$\vect v_0$};
  \draw[vector] (A)--(B) node[midway,right=3pt]{$\vect g t$};
  \draw[vector] (O)--(B) node[midway,below left]{$\vect v$};
  \node[below] at (1,-1.45){(a) 速度三角形};
 \end{scope}
 \begin{scope}[xshift=3cm]
  \coordinate (O) at (0,0);\coordinate (A) at (2,3);\coordinate (B) at (2,1);\coordinate (C) at (2,-1);
  \draw[vector] (O)--(A) node[midway,above left]{$\vect v_0t$};
  \draw[vector] (A)--(B) node[midway,right=4pt]{$\tfrac12\vect g t^2$};
  \draw[vector] (O)--(B) node[pos=.6,below right=3pt]{$\vect r$};
  \draw[TextGray,dashed,->] (B)--(C);
  \draw[TextGray,dashed,->] (O)--(C) node[midway,below left]{$t\vect v$};
  \node[right=4pt,align=left,font=\small] at (2,-.2){再接一段\\$\tfrac12\vect g t^2$};
  \node[below] at (1,-1.45){(b) 位移及辅助三角形};
 \end{scope}
\end{tikzpicture}
\figcap{图2.12：此刻速度向下、位移仍向上；右图加入虚线段后，外侧三角形才与左图相似。}
\end{minipage}\end{center}}

\newcommand{\TennisServeDiagram}{%
\begin{center}\begin{minipage}{.95\linewidth}\centering
\begin{tikzpicture}[>=Latex]
 \draw[thick] (-.15,0)--(6.7,0);\draw[thick] (4,0)--(4,1.2);
 \draw[<->,TextGray] (4.25,0)--(4.25,1.2) node[midway,right=2pt]{$h$};
 \draw[kguide] (0,0)--(0,3) (6,0)--(6,.3);
 \draw[kcurve,->,domain=0:5.6,samples=90] plot(\x,{3-3*\x*\x/31.36});
 \draw[vector] (0,3)--(1.3,3) node[midway,above]{$v_0$};
 \draw[<->,TextGray] (-.5,0)--(-.5,3) node[midway,left]{$H$};
 \draw[<->,TextGray] (0,-.75)--(4,-.75) node[midway,below]{$L$};
 \draw[<->,TextGray] (4,-.75)--(6,-.75) node[midway,below]{$l$};
 \node[above left] at (0,3){$A$};\node[below=3pt] at (4,0){球网};\node[below=3pt] at (6,0){发球线};
\end{tikzpicture}
\figcap{例2.32：合法轨迹从 $A$ 水平出发，从球网上方通过，并落在发球线之前。}
\end{minipage}\end{center}}

\newcommand{\HorizontalThrowImpactDiagram}{%
\begin{center}\begin{minipage}{.95\linewidth}\centering
\begin{tikzpicture}[>=Latex]
 \coordinate (A) at (0,2);\coordinate (B) at (4,0);\coordinate (V) at (5.1,-1.1);\coordinate (X) at (5.3,0);
 \draw[thick] (-.2,0)--(4.35,0);\draw[kguide] (A)--(0,0) (B)--(X);
 \draw[kcurve,dashed,domain=0:4,samples=80] plot(\x,{2-\x*\x/8});
 \draw[vector] (A)--(1.3,2) node[midway,above]{$\vect v_0$};
 \draw[vector] (B)--(V) node[right=4pt]{$\vect v$};
 \draw pic[draw=TextGray,angle radius=6mm] {angle=V--B--X};
 \node[right] at (4.7,-.21){$45^\circ$};
 \draw[<->,TextGray] (-.5,0)--(-.5,2) node[midway,left]{$h$};
 \node[point] at (A){};\node[point] at (B){};
\end{tikzpicture}
\figcap{例2.33：轨迹在抛出点水平，在落地点的切线与水平线成 $45^\circ$。}
\end{minipage}\end{center}}

\newcommand{\ProjectileRangeDiagram}{%
\begin{center}\begin{minipage}{.95\linewidth}\centering
\begin{tikzpicture}[scale=.95,>=Latex]
 \coordinate (A) at (0,0);\coordinate (B) at (7.2,0);\coordinate (V) at (52:1.5);
 \draw[thick] (-.4,0)--(7.6,0);
 \draw[kcurve,->,domain=0:7.2,samples=100] plot(\x,{tan(52)*\x*(1-\x/7.2)});
 \draw[vector] (A)--(V) node[above left]{$v_0$};
 \draw pic[draw=TextGray,angle radius=7mm] {angle=B--A--V};\node at (.97,.39){$\theta$};
 \draw[<->,TextGray] (0,-.55)--(7.2,-.55) node[midway,below]{$S$};
\end{tikzpicture}
\figcap{例2.34：斜抛轨迹先升后降，回到出发高度时的水平距离为射程 $S$。}
\end{minipage}\end{center}}

\newcommand{\ComplementaryProjectileDiagram}{%
\begin{center}\begin{minipage}{.96\linewidth}\centering
\begin{tikzpicture}[scale=.9,>=Latex]
 \coordinate (A) at (0,0);\coordinate (B) at (6.4,0);
 \draw[thick] (-.3,0)--(6.8,0);
 \foreach \ang in {22,68}{\draw[kcurve,dashed,domain=0:6.4,samples=100] plot(\x,{tan(\ang)*\x*(1-\x/6.4)});
 \draw[vector] (A)--(\ang:1.3);}
 \draw[TextGray] (.7,0) arc[start angle=0,end angle=22,radius=.7];
 \draw[TextGray] (1.1,0) arc[start angle=0,end angle=68,radius=1.1];
 \node at (1.39,.95){$\theta_1$};\node at (1.55,.22){$\theta_2$};
 \node[above] at (3.2,4){高弧线};\node[above] at (3.2,.7){低弧线};
 \node[below] at (3.2,-.4){$\theta_1+\theta_2=\pi/2$，两次初速率相同};
\end{tikzpicture}
\figcap{例2.35：互余抛射角对应同一落点；高弧线飞行时间更长。}
\end{minipage}\end{center}}

\newcommand{\ObstacleClearanceDiagram}{%
\begin{center}\begin{minipage}{.96\linewidth}\centering
\begin{tikzpicture}[>=Latex]
 \pgfmathsetmacro{\uu}{(sqrt(25+2.25)+1.5)/5}
 \pgfmathsetmacro{\kk}{(5*\uu-1.5)/25}
 \draw[thick] (-.3,0)--(5.9,0) (5,0)--(5,1.5);
 \draw[kcurve,dashed,->,domain=0:5,samples=90] plot(\x,{\uu*\x-\kk*\x*\x});
 \draw[vector] (0,0)--({1.25/sqrt(1+\uu^2)},{1.25*\uu/sqrt(1+\uu^2)}) node[above left]{$v_0$};
 \draw[TextGray] (.7,0) arc[start angle=0,end angle={atan(\uu)},radius=.7];
 \node at (.95,.42){$\theta$};\node[below left] at (0,0){$A$};
 \draw[<->,TextGray] (0,-.55)--(5,-.55) node[midway,below]{$L$};
 \draw[<->,TextGray] (5.5,0)--(5.5,1.5) node[midway,right]{$H$};
\end{tikzpicture}
\figcap{例2.36：临界轨迹经过障碍顶端；严格越过时须留出正的高度余量。}
\end{minipage}\end{center}}

\newcommand{\ProjectileReachabilityDiagram}{%
\begin{center}\begin{minipage}{.96\linewidth}\centering
\begin{tikzpicture}
\begin{axis}[width=.90\linewidth,height=7cm,axis lines=middle,xlabel={$x$},ylabel={$y$},
 xmin=0,xmax=7.2,ymin=-1,ymax=3.7,ticks=none,clip=true]
 \path[fill=Blue!7] (axis cs:0,-1)--plot[domain=0:6.6,samples=100]
 (axis cs:\x,{3-\x*\x/12})--(axis cs:6.6,-1)--cycle;
 \foreach \ang in {20,35,50,65,78}{
  \addplot[Blue!75,thick,domain=0:6.6,samples=140]{x*tan(\ang)-x*x/(12*cos(\ang)^2)};
 }
 \addplot[DeepBlue,very thick,domain=0:6.6,samples=100]{3-x*x/12};
 \node at (axis cs:5.7,2.8){不可达区};
 \node at (axis cs:3.9,-.5){可达区};
 \node at (axis cs:2.35,3.35){分界线};
 \draw[TextGray] (axis cs:2.35,3.08)--(axis cs:2.35,2.54);
\end{axis}
\end{tikzpicture}
\figcap{例2.37：同一初速率的各条轨迹均在包络线下方，且分别与它相切。}
\end{minipage}\end{center}}

\newcommand{\HillsideTargetDiagram}{%
\begin{center}\begin{minipage}{.96\linewidth}\centering
\begin{tikzpicture}[>=Latex]
 \coordinate (C) at (0,0);\coordinate (T) at (5,1.5);\coordinate (X) at (5.8,0);
 \draw[thick] (C)--(X) (C)--(5.6,1.68);
 \draw[kcurve,dashed,->,domain=0:5,samples=100] plot(\x,{tan(55)*\x-(5*tan(55)-1.5)*\x*\x/25});
 \draw[vector] (C)--(55:1.3) node[above left]{$v_0$};
 \draw[TextGray] (.75,0) arc[start angle=0,end angle={atan(.3)},radius=.75];
 \draw[TextGray] (1.3,0) arc[start angle=0,end angle=55,radius=1.3];
 \node at (.99,.10){$\alpha$};\node at (1.48,.75){$\beta$};
 \draw[kguide] (T)--(5,0);
 \draw[<->,TextGray] (0,-.55)--(5,-.55) node[midway,below]{$L$};
 \node[point,label=below left:$C$] at (C){};\node[point,label=above right:$T$] at (T){};
\end{tikzpicture}
\figcap{例2.39：目标坐标为 $(L,L\tan\alpha)$；炮弹初速度沿轨迹起点的切线。}
\end{minipage}\end{center}}

\newcommand{\HorizontalAtObstacleDiagram}{%
\begin{center}\begin{minipage}{.96\linewidth}\centering
\begin{tikzpicture}[>=Latex]
 \draw[thick] (-.3,0)--(6.5,0) (5,0)--(5,2);
 \draw[kcurve,dashed,domain=0:5,samples=100] plot(\x,{2-2*(\x-5)^2/25});
 \draw[vector] (0,0)--({1.3/sqrt(1.64)},{1.04/sqrt(1.64)}) node[above left]{$v_0$};
 \draw[vector] (5,2)--(6.3,2) node[midway,above]{$v_T$};
 \draw[TextGray] (.7,0) arc[start angle=0,end angle={atan(.8)},radius=.7];
 \node at (1.02,.31){$\theta$};
 \draw[<->,TextGray] (0,-.55)--(5,-.55) node[midway,below]{$L$};
 \draw[<->,TextGray] (5.5,0)--(5.5,2) node[midway,right]{$H$};
 \node[below left] at (0,0){$A$};\node[above left] at (5,2){$T$};
\end{tikzpicture}
\figcap{例2.40：轨迹在障碍顶端 $T$ 达到最高点，此处速度为水平方向。}
\end{minipage}\end{center}}

\newcommand{\InclinedPlaneNormalImpactDiagram}{%
\begin{center}\begin{minipage}{.96\linewidth}\centering
\begin{tikzpicture}[scale=.95,>=Latex]
 \pgfmathsetmacro{\thetaLaunch}{20+atan(1/(2*tan(20)))}
 \pgfmathsetmacro{\uu}{tan(\thetaLaunch)}
 \pgfmathsetmacro{\kk}{(\uu-tan(20))/4.5}
 \coordinate (A) at (0,0);\coordinate (P) at (4.5,{4.5*tan(20)});
 \coordinate (B) at (5.6,{5.6*tan(20)});\coordinate (V) at ($(P)+(-70:1.2)$);
 \draw[thick] (A)--(5.9,0) (A)--(B)--(5.6,0);
 \draw[kcurve,dashed,domain=0:4.5,samples=100] plot(\x,{\uu*\x-\kk*\x*\x});
 \draw[vector] (A)--(\thetaLaunch:1.3) node[above left]{$v_0$};
 \draw[vector] (P)--(V) node[right=3pt]{$v_P$};
 \draw pic[draw=TextGray,angle radius=4mm] {right angle=A--P--V};
 \draw[TextGray] (.7,0) arc[start angle=0,end angle=20,radius=.7];
 \draw[TextGray] (1.3,0) arc[start angle=0,end angle=\thetaLaunch,radius=1.3];
 \node at (.96,.14){$\alpha$};\node at (1.42,.9){$\theta$};
 \node[point,label=above right:$P$] at (P){};
\end{tikzpicture}
\figcap{例2.41：碰撞点 $P$ 在斜面上；轨迹末端切线与斜面垂直。}
\end{minipage}\end{center}}

\newcommand{\StairBounceDiagram}{%
\begin{center}\begin{minipage}{.98\linewidth}\centering
\begin{tikzpicture}[scale=1,>=Latex]
 \pgfmathsetmacro{\hh}{1/(1-.45^2)}\pgfmathsetmacro{\xa}{1.4-1.8/1.45}
 \draw[very thick] (-.1,3)--(2.3,3)--(2.3,2)--(4.1,2)--(4.1,1)--(5.9,1)--(5.9,0)--(7.8,0);
 \draw[kcurve,domain=0:1,samples=70] plot({\xa+(1.4-\xa)*\x},{3+\hh-\hh*\x*\x});
 \foreach \n in {0,1,2}{
  \draw[kcurve,domain=0:1,samples=90] plot({1.4+1.8*\n+1.8*\x},{3-\n+2*.45/(1-.45)*\x-(1+.45)/(1-.45)*\x*\x});
 }
 \foreach \n in {0,1,2,3}{\node[point,label=below:{$C_{\n}$}] at ({1.4+1.8*\n},{3-\n}){};}
 \draw[vector] (\xa,{3+\hh})--++(1.1,0) node[midway,above]{$v_0$};
 \node[above left] at (\xa,{3+\hh}){$A$};
 \draw[<->,TextGray] (-.45,3)--(-.45,{3+\hh}) node[midway,left]{$H$};
 \draw[<->,TextGray] (3.2,3.55)--(5,3.55) node[midway,above]{$d$};
 \draw[<->,TextGray] (7.4,0)--(7.4,1) node[midway,right]{$l$};
\end{tikzpicture}
\figcap{例2.43：$A$ 在首个碰撞台面上方；各 $C_i$ 均在水平踏面上，反弹段可平移重合。}
\end{minipage}\end{center}}

\newcommand{\SphereClearanceDiagram}{%
\begin{center}\begin{minipage}{.98\linewidth}\centering
\begin{tikzpicture}[scale=1.2,>=Latex]
 \foreach \q/\shift/\lab in {.577350269/-2.8/{临界：先与球面相切},.70/2.8/{可行：到球顶前不接触}}{
  \begin{scope}[xshift=\shift cm]
   \pgfmathsetmacro{\ddist}{(\q+sqrt(\q*\q+32/3))/(8/3)}
   \draw[thick] (-2.05,0)--(1.3,0);
   \draw[thick] (0,1) circle (1);
   \draw[kguide] (0,0)--(0,2);
   \draw[kcurve,domain=-\ddist:0,samples=120] plot(\x,{2-\q*\x-4*\x*\x/3});
   \pgfmathsetmacro{\sl}{-\q+8*\ddist/3}
   \draw[vector] (-\ddist,0)--++({.65/sqrt(1+\sl^2)},{.65*\sl/sqrt(1+\sl^2)});
   \draw[vector] (0,2)--++({.7/sqrt(1+\q^2)},{-.7*\q/sqrt(1+\q^2)});
   \node[above right] at (.58,2.02){$\vect v_T$};
   \node[point,label=above left:$T$] at (0,2){};
   \node[point,label=below:$O$] at (0,1){};
   \draw[<->,TextGray] (-\ddist,-.4)--(0,-.4) node[midway,below]{$d$};
   \node[font=\small,align=center] at (-.4,-1.05){\lab};
   \ifdim\q pt<.6pt
    \node[point,label=left:$P$] at (-.8660254,1.5){};
   \fi
  \end{scope}
 }
\end{tikzpicture}
\figcap{例2.45：两条轨迹均在下降时到达球顶 $T$；左图的提前切点 $P$ 决定速度下界。}
\end{minipage}\end{center}}

\newcommand{\CirclePositionDiagram}{%
\begin{center}\begin{tikzpicture}[scale=.9,>=Latex]
 \draw[very thick] (0,0) circle (2);
 \coordinate (O) at (0,0);\coordinate (N) at (2,0);\coordinate (P) at (52:2);
 \draw[very thick] (O)--(N) node[midway,clearlabel,below]{$R$};
 \draw[very thick] (O)--(P);
 \draw[->,TextGray] (.75,0) arc[start angle=0,end angle=52,radius=.75];
 \node[clearlabel] at (1,.43){$\theta$};
 \node[point,label={[clearlabel]below left:$O$}] at (O) {};
 \node[point,label={[clearlabel]right:$N$}] at (N) {};
 \node[point,label={[clearlabel]above right:$P$}] at (P) {};
\end{tikzpicture}\figcap{图2.13：角位置 $\theta$ 由固定射线 $ON$ 转到半径 $OP$ 的有向角确定；逆时针为正。}\end{center}}

\newcommand{\AngularSignDiagram}{%
\begin{center}\begin{tikzpicture}[scale=.75,>=Latex]
 \begin{scope}[xshift=-2.4cm]
  \draw[very thick] (0,0) circle (1.25);\draw[DeepBlue,very thick,->] (1.45,-.2) arc[start angle=-8,end angle=95,radius=1.45];
  \node[clearlabel,below] at (0,-1.65){$\omega>0$（逆时针）};
 \end{scope}
 \begin{scope}[xshift=2.4cm]
  \draw[very thick] (0,0) circle (1.25);\draw[Blue,very thick,->] (-1.45,-.2) arc[start angle=188,end angle=85,radius=1.45];
  \node[clearlabel,below] at (0,-1.65){$\omega<0$（顺时针）};
 \end{scope}
\end{tikzpicture}\figcap{图2.14：角速度的符号只表示转动方向，不表示快慢。}\end{center}}

\newcommand{\CentripetalDerivationDiagram}{%
\begin{center}\begin{minipage}{.98\linewidth}\centering
\begin{tikzpicture}[scale=1,>=Latex]
 \begin{scope}[xshift=-3.5cm]
  \draw[thick] (0,0) circle (1.7);
  \coordinate (A) at (110:1.7);\coordinate (B) at (70:1.7);
  \draw[TextGray] (0,0)--(A) (0,0)--(B);
  \draw[vector] (A)--++(20:1.2) node[pos=.85,above=3pt]{$\vect v_A$};
  \draw[vector] (B)--++(-20:1.2) node[right=3pt]{$\vect v_B$};
  \draw[TextGray] (70:.65) arc[start angle=70,end angle=110,radius=.65];
  \node at (0,1.22){$|\Delta\theta|$};
  \node[point,label=above left:$A$] at (A){};
  \node[point] at (B){};\node[above right] at ($(B)+(.12,.17)$){$B$};
  \node[below] at (0,-2){(a) 同一圆周上的两速度};
 \end{scope}
 \begin{scope}[xshift=2.4cm]
  \coordinate (S) at (0,0);\coordinate (P) at (20:2.5);\coordinate (Q) at (-20:2.5);
  \draw[vector] (S)--(P) node[midway,above left]{$\vect v_A$};
  \draw[vector] (S)--(Q) node[midway,below left]{$\vect v_B$};
  \draw[vector] (P)--(Q) node[midway,right=4pt]{$\Delta\vect v$};
  \draw[TextGray] (-20:.8) arc[start angle=-20,end angle=20,radius=.8];
  \node[right] at (.85,0){$|\Delta\theta|$};
  \node[left] at (S){$S$};\node[above] at (P){$P$};\node[below] at (Q){$Q$};
  \node[below] at (1.1,-2){(b) 平移后作矢量差};
 \end{scope}
\end{tikzpicture}
\figcap{图2.15：右图保持左图两速度的方向；$\Delta\vect v$ 由旧速度终点指向新速度终点。}
\end{minipage}\end{center}}

\newcommand{\RadialCollisionBase}[1]{%
\begin{tikzpicture}[scale=.75,>=Latex]
 \draw[very thick] (0,0) circle (2);\coordinate (P) at (0,2);\coordinate (A) at (25:2);\coordinate (B) at (0,.72);
 \node[point,label={[clearlabel]below left:$O$}] at (0,0) {};
 \node[point,label={[clearlabel]above:$P$}] at (P) {};
 \node[point,label={[clearlabel]below left:$A$}] at (A) {};
 \node[point,label={[clearlabel]right:$B$}] at (B) {};
 \draw[kguide] (0,0)--(P);
 \node[clearlabel,left,xshift=-2pt] at (0,.35){$R$};
 \draw[vector] (B)--(0,1.60) node[midway,clearlabel,right,xshift=3pt]{$#1$};
 \draw[DeepBlue,very thick,->] (4:2.25) arc[start angle=4,end angle=55,radius=2.25];
 \node[right] at (2.42,1.28){$\omega$};
\end{tikzpicture}}

\newcommand{\RadialUniformCollisionDiagram}{%
\begin{center}\RadialCollisionBase{\vect v_B}
\figcap{例2.47：$B$ 以恒定速度沿射线 $OP$ 运动，故与圆周质点 $A$ 的相撞点只能是 $P$。}\end{center}}

\newcommand{\RadialAcceleratedCollisionDiagram}{%
\begin{center}\RadialCollisionBase{\vect a_B}
\figcap{例2.48：$B$ 从 $O$ 由静止沿 $OP$ 匀加速；几何相撞点仍只能是 $P$，但到达时间条件改变。}\end{center}}

\newcommand{\SameCirclePursuitDiagram}{%
\begin{center}\begin{tikzpicture}[scale=.8,>=Latex]
  \draw[very thick] (0,0) circle (2.25);
  \coordinate (O) at (0,0);\coordinate (P) at (90:2.25);\coordinate (Q) at (135:2.25);
  \draw[very thick] (O)--(P);\draw[kguide] (O)--(Q);
  \draw[DeepBlue,very thick,->] (125:2.5) arc[start angle=125,end angle=175,radius=2.5]
    node[midway,clearlabel,left]{$\omega_1$};
  \draw[Blue,very thick,->] (-2.5,-.55) arc[start angle=192,end angle=230,radius=2.5]
    node[midway,clearlabel,left]{$\omega_2$};
  \draw pic[draw=TextGray,->,angle radius=8mm,"$\theta$"{clearlabel}] {angle=P--O--Q};
  \node[point] at (Q) {};\node[left] at ($(Q)+(-.7,.32)$){$Q$};
  \node[point,label={[clearlabel]above:$P$}] at (P) {};
  \node[point,label={[clearlabel]below right:$O$}] at (O) {};
\end{tikzpicture}\figcap{例2.49：两质点从同一点出发沿同一圆周逆时针运动，第二次相遇位置相对 $OP$ 转过角 $\theta$。}\end{center}}

\newcommand{\ConcentricCollinearityDiagram}{%
\begin{center}\begin{tikzpicture}[scale=.8,>=Latex]
  \draw[very thick] (0,0) circle (1.35);\draw[very thick] (0,0) circle (2.35);
  \coordinate (O) at (0,0);\coordinate (A) at (28:1.35);\coordinate (B) at (28:2.35);
  \draw[kguide] (O)--(B);
  \draw[vector] (A)--++(118:1.05) node[midway,clearlabel,right]{$\vect v_A$};
  \draw[vector] (B)--++(118:.82);
  \node at (2.12,2.17){$\vect v_B$};
  \draw[TextGray] (O)--(145:1.35) node[midway,clearlabel,above=3pt]{$R_1$};
  \draw[TextGray] (O)--(225:2.35);
  \node at (-1.75,-.65){$R_2$};
  \node[point] at (O) {};\node at (.25,-.30){$O$};
  \node[point,label={[clearlabel]right:$A$}] at (A) {};
  \node[point,label={[clearlabel]below right:$B$}] at (B) {};
\end{tikzpicture}\figcap{例2.50：$A,B$ 分别在半径 $R_1,R_2$ 的同心圆上；题目要求两点与圆心再次共线。}\end{center}}

\newcommand{\ProjectileCircleCollisionDiagram}{%
\begin{center}\begin{minipage}{.96\linewidth}\centering
\begin{tikzpicture}[scale=1,>=Latex]
 \draw[TextGray,dashed] (3,0) arc[start angle=0,end angle=180,x radius=3,y radius=.75];
 \draw[thick] (-3,0) arc[start angle=180,end angle=360,x radius=3,y radius=.75];
 \coordinate (O) at (0,0);\coordinate (A) at (3,0);\coordinate (S) at (0,3);
 \draw[kguide] (O)--(A) (O)--(S);
 \draw[vector] (S)--(1.25,3) node[midway,above]{$v_0$};
 \draw[kcurve,dashed,->,domain=0:3,samples=90] plot(\x,{3-\x*\x/3});
 \draw[Blue,->,thick] (3.3,0) arc[start angle=0,end angle=52,x radius=3.3,y radius=.83];
 \node[above right] at (2.6,.7){$\omega$};
 \draw[<->,TextGray] (-.45,0)--(-.45,3) node[midway,left]{$H$};
 \node[below] at (1.5,0){$R$};
 \node[point,label=below left:$O$] at (O){};
 \node[point,label=right:$A$] at (A){};
 \node[point,label=above left:{$S\;(B\text{的初始位置})$}] at (S){};
\end{tikzpicture}
\figcap{例2.51：$S$ 在圆心 $O$ 正上方；抛体从 $S$ 水平出发，在圆周上与整周返回的质点相遇。}
\end{minipage}\end{center}}

\newcommand{\AccelerationVelocityAngleDiagram}{%
\begin{center}\begin{minipage}{.96\linewidth}\centering
\begin{tikzpicture}[scale=.95,>=Latex]
 \draw[thick] (0,0) circle (2.0);\coordinate (P) at (35:2);
 \coordinate (V) at ($(P)+(125:1.4)$);\coordinate (A) at ($(P)+(170:1.4)$);
 \draw[kguide] (0,0)--(P) node[pos=.4,below right]{$R$};
 \draw[vector] (P)--(V) node[above right]{$\vect v$};
 \draw[vector] (P)--(A) node[above left]{$\vect a$};
 \draw pic[draw=TextGray,angle radius=7mm] {angle=V--P--A};
 % The circumference crosses this narrow wedge. Show its label separately,
 % retaining the physical vectors and the main angle arc without a mask.
 \begin{scope}[shift={(4.7,1.3)}]
 \draw[vector] (0,0)--(125:1.6) node[above right]{$\vect v$};
 \draw[vector] (0,0)--(170:1.6) node[above left]{$\vect a$};
 \draw[TextGray] (125:.6) arc(125:170:.6);
 \node at (147.5:.95){$\theta$};
 \node[below] at (-.5,-.4){夹角放大};
 \end{scope}
 \node[point,label=below left:$O$] at (0,0){};\node[point] at (P){};
\end{tikzpicture}
\figcap{例2.59：右侧单独放大速度与加速度的夹角 $\theta$；该角取值在 $0$ 与 $\pi$ 之间。}
\end{minipage}\end{center}}

\newcommand{\ConcentricSightDiagram}{%
\begin{center}\begin{minipage}{.96\linewidth}\centering
\begin{tikzpicture}[scale=.9,>=Latex]
 \pgfmathsetmacro{\ax}{1.5^2/2.8}\pgfmathsetmacro{\ay}{sqrt(1.5^2-\ax^2)}
 \coordinate (O) at (0,0);\coordinate (B) at (2.8,0);\coordinate (A) at (\ax,\ay);
 \draw[thick] (O) circle (1.5) (O) circle (2.8);
 \draw[thick] (O)--(B) (B)--(A);\draw[kguide] (O)--(A);
 \draw pic[draw=TextGray,angle radius=4mm] {right angle=O--A--B};
 \draw pic[draw=TextGray,angle radius=8mm] {angle=A--B--O};
 \node at (1.75,.27){$\theta$};
 \draw[TextGray] (O)--(145:1.5) node[midway,below left]{$R_1$};
 \draw[TextGray] (O)--(230:2.8) node[pos=.78,below right=4pt]{$R_2$};
 \node[point,label=below:$O$] at (O){};\node[point,label=right:$B$] at (B){};
 \node[point,label=above:$A$] at (A){};
\end{tikzpicture}
\figcap{例2.52：$A$ 为从 $B$ 引向内圆的切点，因而 $OA\perp AB$。}
\end{minipage}\end{center}}

\newcommand{\RacetrackDiagram}{%
\begin{center}\begin{tikzpicture}[scale=.72,>=Latex]
 \draw[very thick] (-2.4,-2.4)--(-2.4,0) arc[start angle=180,end angle=0,radius=2.4]--(2.4,-2.4);
 \draw[very thick] (-1.4,-2.4)--(-1.4,0) arc[start angle=180,end angle=0,radius=1.4]--(1.4,-2.4);
 \draw[->,DeepBlue,very thick] (-1.4,-1.9)--(-1.4,-.85);
 \draw[->,Blue,very thick] (-2.4,-1.9)--(-2.4,-.85);
 \draw[<->,TextGray] (0,0)--(0,1.4) node[midway,clearlabel,right]{$R_1$};
 \draw[<->,TextGray] (0,0)--(135:2.4) node[pos=.85,clearlabel,above right=3pt]{$R_2$};
 \node[clearlabel,below] at (0,-2.75){内、外路线均为半圆弯道};
\end{tikzpicture}\figcap{例2.57解图：外道允许更高速度，但半圆弧也更长。}\end{center}}

\newcommand{\FixedCoordinatesDiagram}{%
\begin{center}\begin{tikzpicture}[scale=.78,>=Latex]
 \draw[kaxis] (-.5,0)--(6.4,0) node[right,clearlabel]{$x$};
 \draw[kaxis] (0,-.4)--(0,4.2) node[above,clearlabel]{$y$};
 \draw[kcurve] plot[smooth] coordinates{(.5,.5)(1.4,2.8)(2.8,3.4)(4.1,2.0)(5.7,3.2)};
 \coordinate (A) at (1.4,2.8);\coordinate (B) at (4.1,2.0);
 \node[point,label={[clearlabel]above left:$A$}] at (A) {};
 \node[point,label={[clearlabel]below right:$B$}] at (B) {};
 \draw[kguide] (A)--(1.4,0) node[clearlabel,below]{$x(t_1)$};\draw[kguide] (A)--(0,2.8) node[clearlabel,left]{$y(t_1)$};
 \draw[kguide] (B)--(4.1,0) node[clearlabel,below]{$x(t_2)$};\draw[kguide] (B)--(0,2) node[clearlabel,left]{$y(t_2)$};
\end{tikzpicture}\figcap{图2.16：分量运动方程 $x(t),y(t)$ 在每个时刻唯一确定曲线上的位置。}\end{center}}

\newcommand{\ParametricEllipseDiagram}{%
\begin{center}\begin{minipage}{.96\linewidth}\centering
\begin{tikzpicture}[scale=.85,>=Latex]
 \draw[kaxis] (-3.8,0)--(3.9,0) node[right]{$x$};\draw[kaxis] (0,-2.3)--(0,2.5) node[above]{$y$};
 \draw[kcurve] (0,0) ellipse (3 and 1.65);
 \draw[vector,domain=10:40,samples=25] plot({3*cos(\x)},{1.65*sin(\x)});
 \node[point,label=below right:{$t=0$}] at (3,0){};
 \draw[<->,TextGray] (0,-.3)--(3,-.3) node[midway,below]{$a$};
 \draw[<->,TextGray] (.3,0)--(.3,1.65) node[midway,right]{$b$};
\end{tikzpicture}
\figcap{例2.61：闭合椭圆的半轴为 $a,b$；短箭头表示 $\omega>0$ 时的运动方向。}
\end{minipage}\end{center}}

\newcommand{\CycloidDiagram}{%
\begin{center}\begin{tikzpicture}[scale=.72,>=Latex]
 \draw[very thick] (-.4,0)--(10,0);
 \draw[kaxis] (0,0)--(10.2,0) node[right,clearlabel]{$x$};\draw[kaxis] (0,0)--(0,3.4) node[above,clearlabel]{$y$};
 \draw[kcurve,domain=0:6.283,samples=120] plot ({1.35*(\x-sin(\x r))},{1.35*(1-cos(\x r))});
 \draw[very thick] (0,1.35) circle (1.35);\node[point,label={[clearlabel]below left:$P$}] at (0,0) {};
 \draw[vector] (0,1.35)--(1.1,1.35) node[pos=.22,clearlabel,above=3pt]{$v$};
\end{tikzpicture}\figcap{例2.62：无滑动滚动圆周点的摆线。尖点处圆周点与地面瞬时相对静止。}\end{center}}

\newcommand{\HelixDiagram}{%
\begin{center}\begin{tikzpicture}[scale=.76,>=Latex,
  x={(1cm,0cm)},y={(-.50cm,-.28cm)},z={(0cm,.80cm)}]
 \draw[kaxis] (0,0,0)--(3.1,0,0) node[right,clearlabel]{$x$};
 \draw[kaxis] (0,0,0)--(0,2.7,0) node[below left,clearlabel]{$y$};
 \draw[kaxis] (0,0,0)--(0,0,5.2) node[above,clearlabel]{$z$};
 \draw[TextGray,dashed,domain=0:360,samples=90] plot ({1.4*cos(\x)},{1.4*sin(\x)},0);
 \draw[TextGray,dashed,domain=0:360,samples=90] plot ({1.4*cos(\x)},{1.4*sin(\x)},4.4);
 \foreach \a in {0,90,180,270}{\draw[TextGray!65] ({1.4*cos(\a)},{1.4*sin(\a)},0)--({1.4*cos(\a)},{1.4*sin(\a)},4.4);}
 \draw[kcurve,->,domain=0:720,samples=220]
   plot ({1.4*cos(\x)},{1.4*sin(\x)},{.0061*\x});
 \draw[<->,TextGray] (0,0,-.75)--(1.4,0,-.75) node[midway,clearlabel,below]{$R$};
 \draw[<->,TextGray] (1.72,0,0)--(1.72,0,2.2) node[midway,clearlabel,right]{$2\pi v/\omega$};
 \node[clearlabel,right,align=left] at (2.15,0,3.8){水平投影为圆\\竖直方向匀速上升};
\end{tikzpicture}\figcap{例2.63解图：$xy$ 投影是圆，$z$ 坐标匀速增加，合成轨迹为圆柱螺旋线。}\end{center}}

\newcommand{\CircleComponentsDiagram}{%
\begin{center}\begin{tikzpicture}[scale=.78,>=Latex]
  \draw[kaxis] (-3.1,0)--(3.25,0) node[right,clearlabel]{$x$};
  \draw[kaxis] (0,-3.0)--(0,3.15) node[above,clearlabel]{$y$};
  \draw[very thick] (0,0) circle (2.25);
  \coordinate (O) at (0,0);\coordinate (P) at (38:2.25);
  \draw[very thick] (O)--(P) node[midway,clearlabel,below right]{$R$};
  \draw[kguide] (P)--(P|-O) node[clearlabel,below]{$x$};
  \draw[kguide] (P)--(P-|O) node[clearlabel,left]{$y$};
  \draw[vector] (P)--++(128:1.25) node[midway,clearlabel,right]{$\vect v$};
  \draw[vector] (P)--++(218:1.15) node[midway,clearlabel,above left]{$\vect a$};
  \node[point,label={[clearlabel]below left:$O$}] at (O) {};
  \node[point,label={[clearlabel]above right:$P$}] at (P) {};
\end{tikzpicture}\figcap{例2.64：圆周运动的位置投影为 $x,y$；速度沿切线，加速度沿半径指向圆心。}\end{center}}

\newcommand{\RollingContactDiagram}{%
\begin{center}\begin{tikzpicture}[scale=.8,>=Latex]
  \draw[very thick] (-3.7,0)--(3.7,0);
  \draw[very thick] (0,1.55) circle (1.55);
  \coordinate (O) at (0,1.55);\coordinate (P) at (0,0);
  \draw[vector] (O)--++(1.45,0) node[midway,clearlabel,above]{$\vect v_O$};
  \draw[vector] (P)--++(-1.45,0) node[midway,clearlabel,below=2pt]{$\vect v_{P/O}$};
  \draw[vector] (P)--++(1.45,0);
  \node at (1.7,.3){$\vect v_O$};
  \draw[DeepBlue,very thick,->] ({1.85*cos(30)},{1.55+1.85*sin(30)}) arc[start angle=30,end angle=-20,radius=1.85];
  \node[clearlabel,right] at (2.05,1.15){$\omega=v/R$};
  \draw[very thick] (O)--(P) node[midway,clearlabel,right]{$R$};
  \node[point] at (P) {};\node[below=4pt] at (P){$P$};
  \node[point,label={[clearlabel]above left:$O$}] at (O) {};
  \node[clearlabel] at (0,-1.1){$\vect v_P=\vect0$};
\end{tikzpicture}\figcap{例2.65：接地点 $P$ 的平动速度与绕圆心转动速度大小相等、方向相反。}\end{center}}

\newcommand{\ArchimedeanSpiralDiagram}{%
\begin{center}\begin{tikzpicture}[scale=.64,>=Latex]
  \begin{scope}[xshift=-5.7cm]
    \coordinate (O) at (0,0);\coordinate (P) at (38:2.5);
    \node[point,label={[clearlabel]below left:$O$}] at (O) {};
    \draw[kaxis] (O)--(3,0) node[right]{$N$};
    \draw[vector] (O)--(P) node[midway,above left=3pt]{$r=vt$};
    \node[point,label=above:$P$] at (P){};
    \draw[TextGray] (.9,0) arc[start angle=0,end angle=38,radius=.9];
    \node at (1.35,.32){$\theta$};
    \node[clearlabel,below] at (1,-1.0){$\theta=\omega t$};
  \end{scope}
  \begin{scope}[xshift=2.7cm]
    \draw[kaxis] (-4.3,0)--(4.5,0) node[right,clearlabel]{$x$};
    \draw[kaxis] (0,-4.0)--(0,4.2) node[above,clearlabel]{$y$};
    \draw[kcurve,->,domain=0:18.3,samples=180] plot ({.21*\x*cos(\x r)},{.21*\x*sin(\x r)});
    \node[point] at (0,0){};\node at (.3,-.35){$O$};
    \node[clearlabel,below] at (0,-4.55){$r=(v/\omega)\theta$};
  \end{scope}
\end{tikzpicture}\figcap{例2.66：左图在一个时刻确定 $P$ 的极坐标；右图以 $O$ 为极点，连续描出 $r\propto\theta$ 的螺线。}\end{center}}

\newcommand{\NaturalCoordinatesDiagram}{%
\begin{center}\begin{minipage}{.98\linewidth}\centering
\begin{tikzpicture}[scale=1.25,>=Latex]
 \draw[kcurve,domain=-2.7:2.7,samples=120] plot(\x,{.15*\x^3-.6*\x});
 \coordinate (P) at (-1.4,.4284);\coordinate (Q) at (-.2,.1188);
 \pgfmathsetmacro{\ang}{atan(.282)}
 \draw[vector] (P)--++(\ang:1.5) node[above=5pt]{$\vect v(t),\ \vect e_t$};
 \draw[vector] (Q)--++({atan(-.582)}:1.25) node[below right=2pt]{$\vect v(t+\Delta t)$};
 \draw[vector] (P)--++({\ang-90}:1.15) node[below right]{$a_n,\ \vect e_n$};
 \draw[vector] (P)--++({\ang+180}:.8) node[above left=3pt]{$a_t<0$};
 \node[point,label=above left:$P$] at (P){};\node[point] at (Q){};\node[above=5pt] at (Q){$Q$};
\end{tikzpicture}
\figcap{切向沿速度，法向指向 $P$ 处曲线的凹侧；图中切向加速度取减速情形。}
\label{fig:k06-natural-frame}\end{minipage}\end{center}}

\newcommand{\VelocityIncrementNaturalDiagram}{%
\begin{center}\begin{tikzpicture}[scale=.82,>=Latex]
  \coordinate (A) at (-4.0,0);\coordinate (B) at (0,0);\coordinate (C) at (2.25,2.7);
  \coordinate (D) at (2.25,0);\coordinate (E) at (4.1,0);
  \draw[vector] (A)--(B) node[midway,clearlabel,below]{$\vect v(t)$};
  \draw[vector] (B)--(C);\node at (.2,1.2){$\vect a\,\Delta t$};
  \draw[vector] (A)--(C) node[midway,clearlabel,above left]{$\vect v(t+\Delta t)$};
  \draw[kguide] (C)--(D);
  \draw[kguide] (B)--(D);
  \draw[<->,TextGray] (B)--(D) node[midway,clearlabel,below=3pt]{$a\cos\alpha\,\Delta t$};
  \draw[<->,TextGray] (D)--(C) node[midway,clearlabel,right=3pt]{$a\sin\alpha\,\Delta t$};
  \draw pic[draw=TextGray,->,angle radius=7mm] {angle=E--B--C};
  \node[clearlabel] at ($(B)+(1.1,.45)$){$\alpha$};
  \draw pic[draw=TextGray,->,angle radius=10mm] {angle=B--A--C};
  \node at (-2.45,.34){$\Delta\theta$};
  \node[clearlabel,below left] at (A){$A$};\node[clearlabel,below left] at (B){$B$};
  \node[clearlabel,above] at (C){$C$};\node[clearlabel,below right] at (D){$D$};
\end{tikzpicture}\figcap{图2.17：速度三角形把 $\vect a\Delta t$ 分解为平行、垂直于原速度的两部分，分别决定速率变化和方向转角。}\end{center}}

\newcommand{\OsculatingCircleDiagram}{%
\begin{center}\begin{minipage}{.96\linewidth}\centering
\begin{tikzpicture}[scale=1,>=Latex]
 \draw[TextGray,thick] (0,{5/3}) circle ({5/3});
 \draw[kcurve,domain=-2.4:2.4,samples=100] plot(\x,{.3*\x*\x});
 \draw[thick] (-2.5,0)--(2.5,0) node[right]{$S$};
 \draw[TextGray,->] (0,0)--(0,{5/3}) node[midway,right=4pt]{$\rho$};
 \draw[TextGray] (.18,0)--(.18,.18)--(0,.18);
 \node[point,label=below:$P$] at (0,0){};
 \node[point,label=above left:$O$] at (0,{5/3}){};
\end{tikzpicture}
\figcap{图2.18：圆与曲线在 $P$ 处共切线 $S$、共曲率；半径 $PO$ 垂直于 $S$。}
\end{minipage}\end{center}}

% Figure contract: source example 2.69, physical p59; simple geometry, mode C.
% Reading order: broad valley P1 -> ascending P2 -> sharp turn P3 -> P4 left-below P3.
% Preserve the folded topology, not a single-valued y(x) hill. No artistic smoothing of labels.
% Layout: full curve spans x=-2.8..4.75, y=0..2.8, 10cm final width; labels >=9pt.
% Question: only curve and P1..P4. Answer: circles derived from same C2 spline.
% Blue physical curve, thin grey circle arcs, white paper, all labels transparent.
% Expected collision: tiny circle P3 / label; label offset above-right; P4 left.
% No decorative panels. Circle arcs may intersect remote curve segments legitimately.
% Sketch:  \___P1__/P2__/P3
%                            /P4\___

\newcommand{\FoldedCurvePath}{(-2.80000000,1.40000000)
 ..controls (-1.73540121,0.68166038) and (-0.67080242,-0.03667924)..(0.00000000,0.00000000)
 ..controls (0.67080242,0.03667924) and (0.94780848,0.82837734)..(1.45000000,1.55000000)
 ..controls (1.80153407,2.05513586) and (2.26340901,2.52593477)..(2.65000000,2.65000000)
 ..controls (2.81568185,2.70317081) and (2.96753604,2.69265585)..(3.00000000,2.65000000)
 ..controls (3.04328528,2.59312553) and (2.87432126,2.47911170)..(2.75000000,2.30000000)
 ..controls (2.56351810,2.03133245) and (2.47748239,1.61619468)..(2.52000000,1.30000000)
 ..controls (2.59086268,0.77300887) and (3.01881791,0.52085968)..(3.45000000,0.40000000)
 ..controls (3.88118209,0.27914032) and (4.31559105,0.28957016)..(4.75000000,0.30000000)}

\newcommand{\FoldedCurvePoints}{%
\node[point,label={[label distance=3pt]below left:{$P_1$}}] at (0.00000000,0.00000000) {};
\node[point,label={[label distance=3pt]above left:{$P_2$}}] at (1.45000000,1.55000000) {};
\node[point,label={[label distance=3pt]above right:{$P_3$}}] at (3.00000000,2.65000000) {};
\node[point,label={[label distance=3pt]left:{$P_4$}}] at (2.52000000,1.30000000) {};
}

\newcommand{\FoldedCurveCircles}{%
\begin{scope}
\clip (-3.0,-1.1) rectangle (5,3.9);
\draw[TextGray,thin] (-0.04766861,0.87177976) circle (0.87308204);
\draw[TextGray,thin] (5.68213522,-1.39522683) circle (5.15609635);
\draw[TextGray,thin] (2.97003756,2.62719657) circle (0.03765295);
\draw[TextGray,thin] (3.59629028,1.44472504) circle (1.08597703);
\end{scope}
}

\newcommand{\FoldedCurveTipZoom}{%
\begin{scope}[scale=10]
\clip (2.88,2.54) rectangle (3.08,2.78);
\draw[kcurve] \FoldedCurvePath;
\draw[TextGray,thin] (2.97003756,2.62719657) circle (0.03765295);
\end{scope}
}

\newcommand{\FourCurvaturesDiagram}{%
\begin{center}\begin{minipage}{.98\linewidth}\centering
\begin{tikzpicture}[scale=1.15,>=Latex]
 \FoldedCurveCircles
 \draw[kcurve] \FoldedCurvePath;
 \FoldedCurvePoints
\end{tikzpicture}
\figcap{例2.69：题图的回折曲线与对应密切圆；大圆仅显示局部圆弧。}
\par\medskip
\begin{tikzpicture}[>=Latex]
\FoldedCurveTipZoom
\end{tikzpicture}
\par\smallskip {\small $P_3$ 邻域放大十倍：灰色小圆为此处的密切圆。}
\end{minipage}\end{center}}

\newcommand{\ProjectileCurvatureDiagram}{%
\begin{center}\begin{tikzpicture}[scale=.78,>=Latex]
  \begin{scope}[xshift=-3.4cm]
    \draw[kcurve,->,domain=-2:2,samples=70] plot (\x,{1.75-.48*\x*\x});
    \coordinate (U) at (-.9,{1.75-.48*.9*.9});
    \draw[vector] (U)--++({atan(.864)}:1.15) node[midway,clearlabel,above left]{$\vect v$};
    \draw[vector] (U)--++(0,-1.2) node[midway,clearlabel,right]{$\vect g$};
    \node[clearlabel,below] at (0,-1.35){上升：$\vect g\cdot\vect v<0$};
  \end{scope}
  \begin{scope}[xshift=3.4cm]
    \draw[kcurve,->,domain=-2:2,samples=70] plot (\x,{1.75-.48*\x*\x});
    \coordinate (D) at (.9,{1.75-.48*.9*.9});
    \draw[vector] (D)--++({-atan(.864)}:1.15) node[midway,clearlabel,above right]{$\vect v$};
    \draw[vector] (D)--++(0,-1.2) node[midway,clearlabel,left]{$\vect g$};
    \node[clearlabel,below] at (0,-1.35){下降：$\vect g\cdot\vect v>0$};
  \end{scope}
  \begin{scope}[yshift=-3.8cm]
    \draw[kaxis] (-2.4,0)--(2.4,0) node[right,clearlabel]{$x$};
    \draw[kaxis] (0,-.3)--(0,1.7) node[above,clearlabel]{$y$};
    \draw[kcurve,domain=-1.7:1.7,samples=70] plot (\x,{.45*\x*\x});
    \draw[TextGray] (0,{1/.9}) circle ({1/.9});
    \node[point,label={[clearlabel]below right:$x=0$}] at (0,0) {};
    \draw[<->,TextGray] (0,0)--(0,{1/.9}) node[midway,clearlabel,right]{$\rho(0)$};
  \end{scope}
\end{tikzpicture}\figcap{例2.68：上升、下降阶段的速度—重力夹角决定速率增减；下图在抛物线顶点画出密切圆。}\end{center}}

\newcommand{\EllipseCurvatureDiagram}{%
\begin{center}\begin{tikzpicture}[scale=.74,>=Latex]
  \draw[kaxis] (-4.0,0)--(4.0,0) node[right,clearlabel]{$x$};
  \draw[kaxis] (0,-2.6)--(0,2.6) node[above,clearlabel]{$y$};
  \draw[very thick] (0,0) ellipse (3.25 and 1.65);
  \coordinate (A) at (3.25,0);\coordinate (B) at (0,1.65);
  \draw[TextGray] ({3.25-1.65*1.65/3.25},0) circle ({1.65*1.65/3.25});
  \begin{scope}
    \clip (-4.0,-2.55) rectangle (4.0,2.55);
    \draw[TextGray] (0,{1.65-3.25*3.25/1.65}) circle ({3.25*3.25/1.65});
  \end{scope}
  \draw[<->,TextGray] (0,-2.05)--(3.25,-2.05) node[midway,clearlabel,below]{$a$};
  \draw[<->,TextGray] (.25,0)--(.25,1.65) node[midway,clearlabel,right]{$b$};
  \node[point] at (A) {};\node[below right=4pt] at (A){$A$};
  \node[point,label={[clearlabel]above:$B$}] at (B) {};
  \draw[TextGray] (2.62,.78)--(3.28,1.45);
  \node[clearlabel,above] at (3.15,1.52){$\rho_A=b^2/a$};
  \draw[TextGray] (-1.45,1.86)--(-2.25,2.18);
  \node[clearlabel,left] at (-2.18,2.20){$\rho_B=a^2/b$};
\end{tikzpicture}\figcap{例2.70：椭圆长轴端 $A$ 的密切圆较小，短轴端 $B$ 的密切圆较大。}\end{center}}




\newcommand{\CrossRiverRoundTripDiagram}{%
\begin{center}\begin{tikzpicture}[scale=.78,>=Latex]
 \fill[Blue!7] (-3,-1.45) rectangle (3,1.45);
 \draw[very thick] (-3,-1.45)--(3,-1.45) (-3,1.45)--(3,1.45);
 \coordinate (A) at (0,-1.45);\coordinate (B) at (0,1.45);
 \node[point,label={[clearlabel]below:$A$}] at (A) {};
 \node[point,label={[clearlabel]above:$B$}] at (B) {};
 \draw[DeepBlue,very thick,->] (.10,-1.25)--(.10,1.20);
 \draw[DeepBlue,very thick,->] (-.10,1.20)--(-.10,-1.25);
 \draw[vector] (-2.35,.15)--(-1.05,.15) node[midway,clearlabel,above]{$\vect u$};
 \draw[vector] (A)--++(145:1.25) node[near end,clearlabel,below left]{$\vect v$};
 \draw[vector] (B)--++(215:1.25) node[near end,clearlabel,above left]{$\vect v$};
 \draw[<->,TextGray] (2.45,-1.45)--(2.45,1.45) node[midway,clearlabel,right]{$L$};
\end{tikzpicture}\figcap{例2.73解图：$A,B$ 位于两岸正对位置；往返时都用船对水速度的上游分量抵消水流。}\end{center}}

\newcommand{\EqualSpeedFrameDiagram}{%
\begin{center}\begin{tikzpicture}[scale=.82,>=Latex]
 \draw[very thick] (-3,0)--(3,0);
 \coordinate (A) at (-2.1,0);\coordinate (B) at (2.1,0);
 \node[point,label={[clearlabel]below:$A$}] at (A) {};
 \node[point,label={[clearlabel]below:$B$}] at (B) {};
 \draw[vector] (A)--++(1.55,0) node[midway,clearlabel,above]{$v$};
 \draw[vector] (-.75,1.0)--(.75,1.0) node[midway,clearlabel,above]{$S'\text{ 的速度 }u$};
 \draw[->,TextGray] (-1.55,-.58)--(-.2,-.58);
 \draw[->,TextGray] (1.55,-.58)--(.2,-.58);
 \node[below] at (0,-.65){在 $S'$ 中迎面接近};
\end{tikzpicture}\figcap{例2.74解图：在地面系中 $A$ 向右、$B$ 静止，参考系 $S'$ 也向右平动。}\end{center}}



\newcommand{\VelocityTriangleDiagram}{%
\begin{center}\begin{tikzpicture}[scale=.8,>=Latex]
 \coordinate (O) at (0,0);\coordinate (A) at (2.7,0);\coordinate (B) at (2.7,1.8);
 \draw[vector] (O)--(A) node[midway,clearlabel,below]{$\vect v_{P/A}$};
 \draw[vector] (A)--(B) node[midway,clearlabel,right]{$\vect v_{A/O}$};
 \draw[vector] (O)--(B) node[midway,clearlabel,above left]{$\vect v_{P/O}$};
\end{tikzpicture}\figcap{相对速度三角形：先走“物体相对载体”，再走“载体相对地面”，得到物体相对地面。}\end{center}}



























\newcommand{\SphereStreamlinesDiagram}{%
\begin{center}\begin{tikzpicture}[scale=.85,>=Latex]
 \shade[ball color=TextGray!35] (0,0) circle (1.2);
 \input{/Users/dongxingxi/Documents/物理竞赛一轮资料/sample/xes-original/restructured-complete/assets/ch13/sphere-paths.tex}
 \draw[Blue,thick,->] (-4,0)--(-1.25,0);\draw[Blue,thick,->] (1.25,0)--(4,0);
\end{tikzpicture}\figcap{图2.23：圆球绕流的子午面流线。不同流线绕过球面，不在球顶或球底合并。}\end{center}}

\newcommand{\StreamtubeDiagram}{%
\begin{center}\begin{tikzpicture}[scale=.95,>=Latex]
 % Center c(t)=(t,0,.1t^2); section basis e_y and (-.2t,0,1)/sqrt(1+.04t^2).
 \foreach \ang in {0,60,120,180,240,300}{
  \draw[Blue,thick,domain=0:5,samples=75,variable=\t]
   plot ({\t+(.70-.055*\t)*(.45*cos(\ang)-.2*\t*sin(\ang)/sqrt(1+.04*\t*\t))},
         {.1*\t*\t+(.70-.055*\t)*(.22*cos(\ang)+sin(\ang)/sqrt(1+.04*\t*\t))});}
 \foreach \t in {0,5}{
  \draw[Ink,thick,domain=0:360,samples=80,variable=\a]
   plot ({\t+(.70-.055*\t)*(.45*cos(\a)-.2*\t*sin(\a)/sqrt(1+.04*\t*\t))},
         {.1*\t*\t+(.70-.055*\t)*(.22*cos(\a)+sin(\a)/sqrt(1+.04*\t*\t))});}
 \foreach \t/\lab in {0.7/1,2.5/2,4.2/3}{
  \draw[->,DeepBlue,very thick] (\t,{.1*\t*\t})--++(.55,{.11*\t});
  \draw[TextGray,thin] (\t,{.1*\t*\t})--(\t,{.1*\t*\t-1.0});
  \node[below,font=\small] at (\t,{.1*\t*\t-1.07}) {$\vect v_{\lab}$};}
 \node[left] at (-.5,0) {$C_1$};\node[left] at (-.5,-.45) {$S_A$};
 \node[right] at (5.55,2.8) {$C_2$};\node[right] at (5.55,2.35) {$S_B$};
\end{tikzpicture}\figcap{图2.24：闭合截线 $C_1,C_2$ 围出横截面，侧壁由流线构成。箭头沿局部速度方向；图为弯曲流管的透视示意。}\end{center}}

\newcommand{\PipeContinuityDiagram}{%
\begin{center}\begin{tikzpicture}[scale=.74,>=Latex]
 \draw[very thick] (-4,.55)--(-1.5,.55)--(1.2,1.2)--(4,1.2);
 \draw[very thick] (-4,-.55)--(-1.5,-.55)--(1.2,-1.2)--(4,-1.2);
 \draw[vector] (-2.9,0)--(-1.65,0);
 \draw[TextGray,thin] (-2.3,-.12)--(-2.3,-.95);
 \node[below] at (-2.3,-1.0) {$v_A$};
 \draw[vector] (1.65,0)--(2.4,0) node[midway,below=4pt]{$v_B$};
 \draw[DeepBlue] (-3.25,-.55)--(-3.25,.55) node[above=5pt]{$S_A$};
 \draw[DeepBlue] (3.0,-1.2)--(3.0,1.2) node[above=5pt]{$S_B$};
\end{tikzpicture}\figcap{例2.97：管道的两个横截面及流速方向，$v_B$ 为待求量。}\end{center}}

\newcommand{\UniformPipeDiagram}{%
\begin{center}\begin{tikzpicture}[scale=.8,>=Latex]
 \draw[thick] (-4,.65)--(4,.65) (-4,-.65)--(4,-.65);
 \foreach \xx/\lab in {-2.8/1,2.0/2}{
  \draw[dashed,TextGray] (\xx,-.65)--(\xx,.65) node[above=5pt]{$S$};
  \draw[->,DeepBlue,thick] (\xx+.2,0)--(\xx+1.25,0) node[midway,below=5pt]{$v_{\lab}$};}
\end{tikzpicture}\figcap{例2.98：同一圆柱管的两个等截面；在一维稳态模型中比较 $Sv_1$ 与 $Sv_2$。}\end{center}}

\newcommand{\CompressiblePipeDiagram}{%
\begin{center}\begin{tikzpicture}[scale=.8,>=Latex]
 \draw[thick] (-4,.6)--(-1.5,.6)--(1.3,1.3)--(4,1.3);
 \draw[thick] (-4,-.6)--(-1.5,-.6)--(1.3,-1.3)--(4,-1.3);
 \foreach \xx/\rr/\lab in {-3/.6/A,3/1.3/B}{
  \draw[DeepBlue,dashed] (\xx,-\rr)--(\xx,\rr) node[above=6pt]{$S_{\lab}$};
  \node[below=6pt] at (\xx,-\rr) {$\rho_{\lab}$};}
 \draw[->,DeepBlue,thick] (-2.65,0)--(-1.65,0);
 \draw[TextGray,thin] (-1.9,.1)--(-1.9,1.05);
 \node[above] at (-1.9,1.1) {$v_A$};
 \draw[->,DeepBlue,thick] (1.55,0)--(2.5,0) node[midway,below=5pt]{$v_B$};
\end{tikzpicture}\figcap{例2.100：两截面的面积、密度及流速符号，$v_B$ 为待求量。}\end{center}}

\newcommand{\FallingJetDiagram}{%
\begin{center}\begin{tikzpicture}[scale=.85,>=Latex]
 \draw[thick] (-1.8,1)--(-1.8,0)--(-.50,0) (.50,0)--(1.8,0)--(1.8,1);
 \foreach \sgn in {-1,1}{\draw[DeepBlue,thick,domain=0:3.7,samples=90,variable=\s]
  plot ({\sgn*.5/(1+1.2*\s)^.25},{-\s});}
 \draw[TextGray,dashed] (-.75,0)--(1.3,0) (-.75,-3)--(1.3,-3);
 \draw[<->,TextGray] (-1.1,0)--(-1.1,-3) node[midway,left]{$h$};
 \draw[DeepBlue,dashed] (-.5,0)--(.5,0);
 \draw[DeepBlue,thick] ({-.5/4.6^.25},-3)--({.5/4.6^.25},-3);
 \node[above] at (0,.12) {$S$};\node[left] at (-.55,-3.25) {$S(h)$};
 \draw[->,DeepBlue,thick] (1.3,0)--(1.3,-.65) node[right]{$v$};
 \draw[->,DeepBlue,thick] (1.3,-3)--(1.3,-4.1) node[right]{$v(h)$};
\end{tikzpicture}\figcap{例2.99：高度差 $h$ 的末端就是截面 $S(h)$；两速度箭头分别从对应截面起画。}\end{center}}

% Reserve space with discardable glue so a natural page break is not followed
% by a second forced break.  Local to the chapter-four maintained units.
\newcommand{\KinematicKeepSpace}[1]{%
\par\begingroup
\dimen0=#1\relax
\dimen2=\lastskip
\advance\dimen0 by \dimen2
\vskip-\dimen2
\vskip0pt plus\dimen0
\penalty-100
\vskip0pt plus-\dimen0
\vskip\dimen0
\penalty9999
\vskip-\dimen0
\vskip\dimen2
\endgroup}
% Maintained chapter-four figures. Transparent labels; all panels use explicit geometry.
% Generated by generate_ch04_geometry.py.
\newcommand{\MarsCoplanarPath}{\draw[kcurve] plot coordinates {(-3.3333333,-0.5510407) (-3.3055556,-0.5222965) (-3.2777778,-0.4937709) (-3.2500000,-0.4654674) (-3.2222222,-0.4373898) (-3.1944444,-0.4095421) (-3.1666667,-0.3819279) (-3.1388889,-0.3545514) (-3.1111111,-0.3274165) (-3.0833333,-0.3005273) (-3.0555556,-0.2738881) (-3.0277778,-0.2475030) (-3.0000000,-0.2213764) (-2.9722222,-0.1955128) (-2.9444444,-0.1699166) (-2.9166667,-0.1445923) (-2.8888889,-0.1195447) (-2.8611111,-0.0947784) (-2.8333333,-0.0702983) (-2.8055556,-0.0461093) (-2.7777778,-0.0222163) (-2.7500000,0.0013755) (-2.7222222,0.0246611) (-2.6944444,0.0476352) (-2.6666667,0.0702925) (-2.6388889,0.0926277) (-2.6111111,0.1146351) (-2.5833333,0.1363092) (-2.5555556,0.1576444) (-2.5277778,0.1786350) (-2.5000000,0.1992751) (-2.4722222,0.2195588) (-2.4444444,0.2394801) (-2.4166667,0.2590329) (-2.3888889,0.2782111) (-2.3611111,0.2970085) (-2.3333333,0.3154188) (-2.3055556,0.3334355) (-2.2777778,0.3510522) (-2.2500000,0.3682624) (-2.2222222,0.3850595) (-2.1944444,0.4014369) (-2.1666667,0.4173877) (-2.1388889,0.4329052) (-2.1111111,0.4479827) (-2.0833333,0.4626132) (-2.0555556,0.4767897) (-2.0277778,0.4905054) (-2.0000000,0.5037533) (-1.9722222,0.5165262) (-1.9444444,0.5288173) (-1.9166667,0.5406194) (-1.8888889,0.5519254) (-1.8611111,0.5627284) (-1.8333333,0.5730214) (-1.8055556,0.5827972) (-1.7777778,0.5920490) (-1.7500000,0.6007698) (-1.7222222,0.6089529) (-1.6944444,0.6165913) (-1.6666667,0.6236784) (-1.6388889,0.6302076) (-1.6111111,0.6361725) (-1.5833333,0.6415667) (-1.5555556,0.6463840) (-1.5277778,0.6506184) (-1.5000000,0.6542641) (-1.4722222,0.6573154) (-1.4444444,0.6597670) (-1.4166667,0.6616138) (-1.3888889,0.6628508) (-1.3611111,0.6634734) (-1.3333333,0.6634776) (-1.3055556,0.6628593) (-1.2777778,0.6616150) (-1.2500000,0.6597415) (-1.2222222,0.6572363) (-1.1944444,0.6540970) (-1.1666667,0.6503219) (-1.1388889,0.6459097) (-1.1111111,0.6408597) (-1.0833333,0.6351718) (-1.0555556,0.6288464) (-1.0277778,0.6218845) (-1.0000000,0.6142879) (-0.9722222,0.6060591) (-0.9444444,0.5972011) (-0.9166667,0.5877179) (-0.8888889,0.5776141) (-0.8611111,0.5668950) (-0.8333333,0.5555671) (-0.8055556,0.5436373) (-0.7777778,0.5311136) (-0.7500000,0.5180049) (-0.7222222,0.5043209) (-0.6944444,0.4900722) (-0.6666667,0.4752704) (-0.6388889,0.4599281) (-0.6111111,0.4440586) (-0.5833333,0.4276764) (-0.5555556,0.4107969) (-0.5277778,0.3934361) (-0.5000000,0.3756115) (-0.4722222,0.3573410) (-0.4444444,0.3386438) (-0.4166667,0.3195397) (-0.3888889,0.3000495) (-0.3611111,0.2801948) (-0.3333333,0.2599979) (-0.3055556,0.2394821) (-0.2777778,0.2186712) (-0.2500000,0.1975897) (-0.2222222,0.1762629) (-0.1944444,0.1547165) (-0.1666667,0.1329767) (-0.1388889,0.1110704) (-0.1111111,0.0890246) (-0.0833333,0.0668668) (-0.0555556,0.0446247) (-0.0277778,0.0223265) (0.0000000,0.0000000) (0.0277778,-0.0223265) (0.0555556,-0.0446247) (0.0833333,-0.0668668) (0.1111111,-0.0890246) (0.1388889,-0.1110704) (0.1666667,-0.1329767) (0.1944444,-0.1547165) (0.2222222,-0.1762629) (0.2500000,-0.1975897) (0.2777778,-0.2186712) (0.3055556,-0.2394821) (0.3333333,-0.2599979) (0.3611111,-0.2801948) (0.3888889,-0.3000495) (0.4166667,-0.3195397) (0.4444444,-0.3386438) (0.4722222,-0.3573410) (0.5000000,-0.3756115) (0.5277778,-0.3934361) (0.5555556,-0.4107969) (0.5833333,-0.4276764) (0.6111111,-0.4440586) (0.6388889,-0.4599281) (0.6666667,-0.4752704) (0.6944444,-0.4900722) (0.7222222,-0.5043209) (0.7500000,-0.5180049) (0.7777778,-0.5311136) (0.8055556,-0.5436373) (0.8333333,-0.5555671) (0.8611111,-0.5668950) (0.8888889,-0.5776141) (0.9166667,-0.5877179) (0.9444444,-0.5972011) (0.9722222,-0.6060591) (1.0000000,-0.6142879) (1.0277778,-0.6218845) (1.0555556,-0.6288464) (1.0833333,-0.6351718) (1.1111111,-0.6408597) (1.1388889,-0.6459097) (1.1666667,-0.6503219) (1.1944444,-0.6540970) (1.2222222,-0.6572363) (1.2500000,-0.6597415) (1.2777778,-0.6616150) (1.3055556,-0.6628593) (1.3333333,-0.6634776) (1.3611111,-0.6634734) (1.3888889,-0.6628508) (1.4166667,-0.6616138) (1.4444444,-0.6597670) (1.4722222,-0.6573154) (1.5000000,-0.6542641) (1.5277778,-0.6506184) (1.5555556,-0.6463840) (1.5833333,-0.6415667) (1.6111111,-0.6361725) (1.6388889,-0.6302076) (1.6666667,-0.6236784) (1.6944444,-0.6165913) (1.7222222,-0.6089529) (1.7500000,-0.6007698) (1.7777778,-0.5920490) (1.8055556,-0.5827972) (1.8333333,-0.5730214) (1.8611111,-0.5627284) (1.8888889,-0.5519254) (1.9166667,-0.5406194) (1.9444444,-0.5288173) (1.9722222,-0.5165262) (2.0000000,-0.5037533) (2.0277778,-0.4905054) (2.0555556,-0.4767897) (2.0833333,-0.4626132) (2.1111111,-0.4479827) (2.1388889,-0.4329052) (2.1666667,-0.4173877) (2.1944444,-0.4014369) (2.2222222,-0.3850595) (2.2500000,-0.3682624) (2.2777778,-0.3510522) (2.3055556,-0.3334355) (2.3333333,-0.3154188) (2.3611111,-0.2970085) (2.3888889,-0.2782111) (2.4166667,-0.2590329) (2.4444444,-0.2394801) (2.4722222,-0.2195588) (2.5000000,-0.1992751) (2.5277778,-0.1786350) (2.5555556,-0.1576444) (2.5833333,-0.1363092) (2.6111111,-0.1146351) (2.6388889,-0.0926277) (2.6666667,-0.0702925) (2.6944444,-0.0476352) (2.7222222,-0.0246611) (2.7500000,-0.0013755) (2.7777778,0.0222163) (2.8055556,0.0461093) (2.8333333,0.0702983) (2.8611111,0.0947784) (2.8888889,0.1195447) (2.9166667,0.1445923) (2.9444444,0.1699166) (2.9722222,0.1955128) (3.0000000,0.2213764) (3.0277778,0.2475030) (3.0555556,0.2738881) (3.0833333,0.3005273) (3.1111111,0.3274165) (3.1388889,0.3545514) (3.1666667,0.3819279) (3.1944444,0.4095421) (3.2222222,0.4373898) (3.2500000,0.4654674) (3.2777778,0.4937709) (3.3055556,0.5222965) (3.3333333,0.5510407)};}
\newcommand{\MarsInclinedPath}{\draw[TextGray,thick] plot coordinates {(-3.3295465,-2.0955284) (-3.2802682,-2.0811303) (-3.2311171,-2.0666671) (-3.1820946,-2.0521385) (-3.1332019,-2.0375442) (-3.0844402,-2.0228836) (-3.0358108,-2.0081564) (-2.9873151,-1.9933622) (-2.9389542,-1.9785006) (-2.8907295,-1.9635712) (-2.8426423,-1.9485736) (-2.7946940,-1.9335072) (-2.7468859,-1.9183719) (-2.6992194,-1.9031670) (-2.6516959,-1.8878922) (-2.6043168,-1.8725471) (-2.5570835,-1.8571312) (-2.5099975,-1.8416442) (-2.4630601,-1.8260855) (-2.4162729,-1.8104548) (-2.3696374,-1.7947516) (-2.3231551,-1.7789754) (-2.2768275,-1.7631259) (-2.2306561,-1.7472026) (-2.1846425,-1.7312051) (-2.1387883,-1.7151328) (-2.0930950,-1.6989855) (-2.0475643,-1.6827625) (-2.0021979,-1.6664635) (-1.9569973,-1.6500881) (-1.9119642,-1.6336357) (-1.8671004,-1.6171059) (-1.8224075,-1.6004983) (-1.7778873,-1.5838123) (-1.7335415,-1.5670477) (-1.6893719,-1.5502038) (-1.6453803,-1.5332802) (-1.6015685,-1.5162765) (-1.5579383,-1.4991923) (-1.5144916,-1.4820270) (-1.4712303,-1.4647802) (-1.4281562,-1.4474514) (-1.3852712,-1.4300402) (-1.3425774,-1.4125461) (-1.3000767,-1.3949687) (-1.2577710,-1.3773075) (-1.2156624,-1.3595620) (-1.1737529,-1.3417318) (-1.1320446,-1.3238164) (-1.0905395,-1.3058154) (-1.0492397,-1.2877284) (-1.0081475,-1.2695548) (-0.9672648,-1.2512942) (-0.9265939,-1.2329463) (-0.8861370,-1.2145104) (-0.8458963,-1.1959863) (-0.8058741,-1.1773734) (-0.7660726,-1.1586713) (-0.7264942,-1.1398796) (-0.6871411,-1.1209979) (-0.6480158,-1.1020257) (-0.6091205,-1.0829626) (-0.5704577,-1.0638082) (-0.5320299,-1.0445622) (-0.4938394,-1.0252240) (-0.4558888,-1.0057933) (-0.4181806,-0.9862697) (-0.3807174,-0.9666528) (-0.3435015,-0.9469423) (-0.3065358,-0.9271377) (-0.2698228,-0.9072387) (-0.2333651,-0.8872450) (-0.1971654,-0.8671562) (-0.1612265,-0.8469720) (-0.1255510,-0.8266920) (-0.0901418,-0.8063159) (-0.0550015,-0.7858435) (-0.0201331,-0.7652744) (0.0144606,-0.7446083) (0.0487768,-0.7238450) (0.0828125,-0.7029842) (0.1165648,-0.6820258) (0.1500307,-0.6609693) (0.1832073,-0.6398147) (0.2160915,-0.6185618) (0.2486803,-0.5972103) (0.2809706,-0.5757601) (0.3129592,-0.5542111) (0.3446432,-0.5325631) (0.3760191,-0.5108160) (0.4070840,-0.4889698) (0.4378346,-0.4670243) (0.4682675,-0.4449796) (0.4983796,-0.4228356) (0.5281674,-0.4005923) (0.5576278,-0.3782497) (0.5867573,-0.3558080) (0.6155525,-0.3332672) (0.6440100,-0.3106273) (0.6721263,-0.2878886) (0.6998980,-0.2650512) (0.7273216,-0.2421154) (0.7543935,-0.2190813) (0.7811103,-0.1959492) (0.8074682,-0.1727194) (0.8334638,-0.1493923) (0.8590933,-0.1259683) (0.8843532,-0.1024477) (0.9092398,-0.0788311) (0.9337493,-0.0551189) (0.9578781,-0.0313116) (0.9816224,-0.0074100) (1.0049784,0.0165855) (1.0279425,0.0406742) (1.0505107,0.0648552) (1.0726793,0.0891278) (1.0944445,0.1134912) (1.1158024,0.1379444) (1.1367491,0.1624866) (1.1572809,0.1871168) (1.1773937,0.2118338) (1.1970837,0.2366365) (1.2163470,0.2615238) (1.2351798,0.2864944) (1.2535780,0.3115470) (1.2715378,0.3366802) (1.2890553,0.3618925) (1.3061265,0.3871823) (1.3227475,0.4125482) (1.3389144,0.4379884) (1.3546232,0.4635011) (1.3698701,0.4890844) (1.3846513,0.5147366) (1.3989627,0.5404554) (1.4128005,0.5662389) (1.4261608,0.5920848) (1.4390399,0.6179909) (1.4514339,0.6439547) (1.4633391,0.6699738) (1.4747515,0.6960456) (1.4856676,0.7221674) (1.4960836,0.7483365) (1.5059959,0.7745500) (1.5154007,0.8008047) (1.5242946,0.8270978) (1.5326740,0.8534258) (1.5405354,0.8797856) (1.5478754,0.9061736) (1.5546905,0.9325863) (1.5609776,0.9590199) (1.5667332,0.9854707) (1.5719543,1.0119348) (1.5766377,1.0384079) (1.5807804,1.0648861) (1.5843795,1.0913649) (1.5874320,1.1178398) (1.5899353,1.1443064) (1.5918866,1.1707597) (1.5932834,1.1971951) (1.5941232,1.2236074) (1.5944036,1.2499916)};
\draw[DeepBlue,thick] plot coordinates {(1.5944036,1.2499916) (1.5941294,1.2760148) (1.5933054,1.3020001) (1.5919299,1.3279421) (1.5900008,1.3538355) (1.5875167,1.3796745) (1.5844759,1.4054533) (1.5808772,1.4311661) (1.5767192,1.4568067) (1.5720010,1.4823691) (1.5667215,1.5078468) (1.5608801,1.5332334) (1.5544761,1.5585223) (1.5475092,1.5837068) (1.5399791,1.6087800) (1.5318859,1.6337349) (1.5232296,1.6585643) (1.5140107,1.6832610) (1.5042297,1.7078176) (1.4938874,1.7322266) (1.4829847,1.7564803) (1.4715230,1.7805711) (1.4595036,1.8044909) (1.4469283,1.8282319) (1.4337989,1.8517860) (1.4201176,1.8751449) (1.4058869,1.8983004) (1.3911094,1.9212441) (1.3757880,1.9439675) (1.3599260,1.9664621) (1.3435267,1.9887192) (1.3265940,2.0107302) (1.3091318,2.0324862) (1.2911445,2.0539784) (1.2726366,2.0751981) (1.2536130,2.0961362) (1.2340789,2.1167838) (1.2140397,2.1371319) (1.1935012,2.1571716) (1.1724694,2.1768938) (1.1509508,2.1962896) (1.1289519,2.2153500) (1.1064797,2.2340659) (1.0835415,2.2524286) (1.0601450,2.2704290) (1.0362979,2.2880583) (1.0120084,2.3053078) (0.9872852,2.3221687) (0.9621369,2.3386324) (0.9365727,2.3546905) (0.9106020,2.3703343) (0.8842345,2.3855558) (0.8574802,2.4003466) (0.8303493,2.4146987) (0.8028525,2.4286044) (0.7750006,2.4420557) (0.7468048,2.4550453) (0.7182764,2.4675657) (0.6894271,2.4796097) (0.6602689,2.4911705) (0.6308139,2.5022413) (0.6010746,2.5128157) (0.5710636,2.5228872) (0.5407939,2.5324501) (0.5102784,2.5414984) (0.4795306,2.5500267) (0.4485639,2.5580300) (0.4173921,2.5655031) (0.3860290,2.5724416) (0.3544888,2.5788412) (0.3227855,2.5846979) (0.2909335,2.5900080) (0.2589474,2.5947682) (0.2268418,2.5989754) (0.1946313,2.6026272) (0.1623307,2.6057210) (0.1299549,2.6082551) (0.0975190,2.6102277) (0.0650378,2.6116375) (0.0325264,2.6124838) (0.0000000,2.6127660) (-0.0325264,2.6124838) (-0.0650378,2.6116375) (-0.0975190,2.6102277) (-0.1299549,2.6082551) (-0.1623307,2.6057210) (-0.1946313,2.6026272) (-0.2268418,2.5989754) (-0.2589474,2.5947682) (-0.2909335,2.5900080) (-0.3227855,2.5846979) (-0.3544888,2.5788412) (-0.3860290,2.5724416) (-0.4173921,2.5655031) (-0.4485639,2.5580300) (-0.4795306,2.5500267) (-0.5102784,2.5414984) (-0.5407939,2.5324501) (-0.5710636,2.5228872) (-0.6010746,2.5128157) (-0.6308139,2.5022413) (-0.6602689,2.4911705) (-0.6894271,2.4796097) (-0.7182764,2.4675657) (-0.7468048,2.4550453) (-0.7750006,2.4420557) (-0.8028525,2.4286044) (-0.8303493,2.4146987) (-0.8574802,2.4003466) (-0.8842345,2.3855558) (-0.9106020,2.3703343) (-0.9365727,2.3546905) (-0.9621369,2.3386324) (-0.9872852,2.3221687) (-1.0120084,2.3053078) (-1.0362979,2.2880583) (-1.0601450,2.2704290) (-1.0835415,2.2524286) (-1.1064797,2.2340659) (-1.1289519,2.2153500) (-1.1509508,2.1962896) (-1.1724694,2.1768938) (-1.1935012,2.1571716) (-1.2140397,2.1371319) (-1.2340789,2.1167838) (-1.2536130,2.0961362) (-1.2726366,2.0751981) (-1.2911445,2.0539784) (-1.3091318,2.0324862) (-1.3265940,2.0107302) (-1.3435267,1.9887192) (-1.3599260,1.9664621) (-1.3757880,1.9439675) (-1.3911094,1.9212441) (-1.4058869,1.8983004) (-1.4201176,1.8751449) (-1.4337989,1.8517860) (-1.4469283,1.8282319) (-1.4595036,1.8044909) (-1.4715230,1.7805711) (-1.4829847,1.7564803) (-1.4938874,1.7322266) (-1.5042297,1.7078176) (-1.5140107,1.6832610) (-1.5232296,1.6585643) (-1.5318859,1.6337349) (-1.5399791,1.6087800) (-1.5475092,1.5837068) (-1.5544761,1.5585223) (-1.5608801,1.5332334) (-1.5667215,1.5078468) (-1.5720010,1.4823691) (-1.5767192,1.4568067) (-1.5808772,1.4311661) (-1.5844759,1.4054533) (-1.5875167,1.3796745) (-1.5900008,1.3538355) (-1.5919299,1.3279421) (-1.5933054,1.3020001) (-1.5941294,1.2760148) (-1.5944036,1.2499916)};
\draw[TextGray,thick] plot coordinates {(-1.5944036,1.2499916) (-1.5941232,1.2236074) (-1.5932834,1.1971951) (-1.5918866,1.1707597) (-1.5899353,1.1443064) (-1.5874320,1.1178398) (-1.5843795,1.0913649) (-1.5807804,1.0648861) (-1.5766377,1.0384079) (-1.5719543,1.0119348) (-1.5667332,0.9854707) (-1.5609776,0.9590199) (-1.5546905,0.9325863) (-1.5478754,0.9061736) (-1.5405354,0.8797856) (-1.5326740,0.8534258) (-1.5242946,0.8270978) (-1.5154007,0.8008047) (-1.5059959,0.7745500) (-1.4960836,0.7483365) (-1.4856676,0.7221674) (-1.4747515,0.6960456) (-1.4633391,0.6699738) (-1.4514339,0.6439547) (-1.4390399,0.6179909) (-1.4261608,0.5920848) (-1.4128005,0.5662389) (-1.3989627,0.5404554) (-1.3846513,0.5147366) (-1.3698701,0.4890844) (-1.3546232,0.4635011) (-1.3389144,0.4379884) (-1.3227475,0.4125482) (-1.3061265,0.3871823) (-1.2890553,0.3618925) (-1.2715378,0.3366802) (-1.2535780,0.3115470) (-1.2351798,0.2864944) (-1.2163470,0.2615238) (-1.1970837,0.2366365) (-1.1773937,0.2118338) (-1.1572809,0.1871168) (-1.1367491,0.1624866) (-1.1158024,0.1379444) (-1.0944445,0.1134912) (-1.0726793,0.0891278) (-1.0505107,0.0648552) (-1.0279425,0.0406742) (-1.0049784,0.0165855) (-0.9816224,-0.0074100) (-0.9578781,-0.0313116) (-0.9337493,-0.0551189) (-0.9092398,-0.0788311) (-0.8843532,-0.1024477) (-0.8590933,-0.1259683) (-0.8334638,-0.1493923) (-0.8074682,-0.1727194) (-0.7811103,-0.1959492) (-0.7543935,-0.2190813) (-0.7273216,-0.2421154) (-0.6998980,-0.2650512) (-0.6721263,-0.2878886) (-0.6440100,-0.3106273) (-0.6155525,-0.3332672) (-0.5867573,-0.3558080) (-0.5576278,-0.3782497) (-0.5281674,-0.4005923) (-0.4983796,-0.4228356) (-0.4682675,-0.4449796) (-0.4378346,-0.4670243) (-0.4070840,-0.4889698) (-0.3760191,-0.5108160) (-0.3446432,-0.5325631) (-0.3129592,-0.5542111) (-0.2809706,-0.5757601) (-0.2486803,-0.5972103) (-0.2160915,-0.6185618) (-0.1832073,-0.6398147) (-0.1500307,-0.6609693) (-0.1165648,-0.6820258) (-0.0828125,-0.7029842) (-0.0487768,-0.7238450) (-0.0144606,-0.7446083) (0.0201331,-0.7652744) (0.0550015,-0.7858435) (0.0901418,-0.8063159) (0.1255510,-0.8266920) (0.1612265,-0.8469720) (0.1971654,-0.8671562) (0.2333651,-0.8872450) (0.2698228,-0.9072387) (0.3065358,-0.9271377) (0.3435015,-0.9469423) (0.3807174,-0.9666528) (0.4181806,-0.9862697) (0.4558888,-1.0057933) (0.4938394,-1.0252240) (0.5320299,-1.0445622) (0.5704577,-1.0638082) (0.6091205,-1.0829626) (0.6480158,-1.1020257) (0.6871411,-1.1209979) (0.7264942,-1.1398796) (0.7660726,-1.1586713) (0.8058741,-1.1773734) (0.8458963,-1.1959863) (0.8861370,-1.2145104) (0.9265939,-1.2329463) (0.9672648,-1.2512942) (1.0081475,-1.2695548) (1.0492397,-1.2877284) (1.0905395,-1.3058154) (1.1320446,-1.3238164) (1.1737529,-1.3417318) (1.2156624,-1.3595620) (1.2577710,-1.3773075) (1.3000767,-1.3949687) (1.3425774,-1.4125461) (1.3852712,-1.4300402) (1.4281562,-1.4474514) (1.4712303,-1.4647802) (1.5144916,-1.4820270) (1.5579383,-1.4991923) (1.6015685,-1.5162765) (1.6453803,-1.5332802) (1.6893719,-1.5502038) (1.7335415,-1.5670477) (1.7778873,-1.5838123) (1.8224075,-1.6004983) (1.8671004,-1.6171059) (1.9119642,-1.6336357) (1.9569973,-1.6500881) (2.0021979,-1.6664635) (2.0475643,-1.6827625) (2.0930950,-1.6989855) (2.1387883,-1.7151328) (2.1846425,-1.7312051) (2.2306561,-1.7472026) (2.2768275,-1.7631259) (2.3231551,-1.7789754) (2.3696374,-1.7947516) (2.4162729,-1.8104548) (2.4630601,-1.8260855) (2.5099975,-1.8416442) (2.5570835,-1.8571312) (2.6043168,-1.8725471) (2.6516959,-1.8878922) (2.6992194,-1.9031670) (2.7468859,-1.9183719) (2.7946940,-1.9335072) (2.8426423,-1.9485736) (2.8907295,-1.9635712) (2.9389542,-1.9785006) (2.9873151,-1.9933622) (3.0358108,-2.0081564) (3.0844402,-2.0228836) (3.1332019,-2.0375442) (3.1820946,-2.0521385) (3.2311171,-2.0666671) (3.2802682,-2.0811303) (3.3295465,-2.0955284)};
\draw[->,DeepBlue,thick] plot coordinates {(-0.8732643,-1.2086051) (-0.8296331,-1.1884459) (-0.7862609,-1.1681816) (-0.7431505,-1.1478119) (-0.7003050,-1.1273361) (-0.6577273,-1.1067537) (-0.6154204,-1.0860642) (-0.5733875,-1.0652670) (-0.5316316,-1.0443616)};
\draw[->,DeepBlue,thick] plot coordinates {(0.6894671,2.4795935) (0.6573750,2.4922850) (0.6249252,2.5043828) (0.5921344,2.5158783) (0.5590194,2.5267631) (0.5255976,2.5370295) (0.4918864,2.5466700) (0.4579039,2.5556774) (0.4236680,2.5640451)};
\draw[->,DeepBlue,thick] plot coordinates {(0.8732643,-1.2086051) (0.9171514,-1.2286600) (0.9612918,-1.2486111) (1.0056825,-1.2684589) (1.0503207,-1.2882040) (1.0952037,-1.3078471) (1.1403288,-1.3273886) (1.1856933,-1.3468291) (1.2312945,-1.3661693)};}

\newcommand{\ReferenceFramesDiagram}{%
\begin{center}\begin{minipage}{\linewidth}\centering\begin{tikzpicture}[>=Latex,scale=.88]
\draw[kaxis] (-.3,0)--(7,0) node[right]{$S:x$};
\draw[kaxis] (2,1.6)--(7,1.6) node[right]{$S':x'$};
\coordinate (P) at (5.7,2.4);
\draw[kguide] (2,-.6)--(2,1.6) (P)--(5.7,-.6);
\fill (0,0) circle(1.5pt) (2,1.6) circle(1.5pt) (P) circle(1.5pt);
\node[below left] at (0,0){$O$};\node[above left] at (2,1.6){$O'$};
\node[above] at (P){$P$};
\draw[<->] (0,-.65)--(2,-.65) node[midway,below]{$x_{O'}$};
\draw[<->,DeepBlue] (0,-1.25)--(5.7,-1.25) node[midway,below]{$x$};
\draw[<->,DeepBlue] (2,.8)--(5.7,.8) node[midway,above]{$x'$};
\draw[vector] (2.6,2.35)--(3.9,2.35) node[midway,above]{$u$};
\draw (5.7,0)--(5.7,.12) (5.7,1.6)--(5.7,1.72);
\end{tikzpicture}\figcap{图2.19：同一点的两个横向投影坐标满足 $x=x_{O'}+x'$；轴的上下错开不表示额外横向距离。}\end{minipage}\end{center}}

\newcommand{\AircraftWindDiagram}{%
\begin{center}\begin{minipage}{\linewidth}\centering\begin{tikzpicture}[>=Latex,scale=.88]
\begin{scope}
\draw[kaxis] (0,0)--(4,0) node[right]{东};\draw[kaxis] (0,0)--(0,2.4) node[above]{北};
\draw[vector] (0,0)--(3,0) node[midway,below]{$\vect V$};
\draw[vector] (3,0)--(3,1.5) node[midway,right]{$\vect u$};
\draw[vector] (0,0)--(3,1.5) node[midway,above left]{$\vect v_{\text{地}}$};
\node at (1.7,-.85){(a) 机头正东};
\end{scope}
\begin{scope}[xshift=7cm]
\draw[kaxis] (0,0)--(3.8,0) node[right]{东};\draw[kaxis] (0,0)--(0,2.4) node[above]{北};
\draw[vector] (0,0)--(2.598,-1.5) node[midway,below left]{$\vect V$};
\draw[vector] (2.598,-1.5)--(2.598,0) node[midway,right]{$\vect u$};
\draw[vector] (0,0)--(2.598,0) node[midway,above]{$\vect v_{\text{地}}$};
\draw (.8,0) arc(0:-30:.8);\node at (1.15,-.3){$\theta$};
\node at (1.5,-2.1){(b) 航迹正东};
\end{scope}
\end{tikzpicture}\figcap{例2.81解图：风速向北；两图均按 $\vect v_{\text{地}}=\vect V+\vect u$ 首尾相接。}\end{minipage}\end{center}}

\newcommand{\RiverCrossingDiagram}{%
\begin{center}\begin{minipage}{\linewidth}\centering\begin{tikzpicture}[>=Latex,scale=.9]
\begin{scope}
\fill[Blue!5] (-.5,0) rectangle (3.9,3);
\draw[thick] (-.5,0)--(3.9,0) (-.5,3)--(3.9,3);
\coordinate (A) at (0,0);\draw[kguide] (A)--(2.25,3);
\fill (A) circle(1.5pt);\node[below] at (A){$A$};
\draw[vector] (A)--(0,1.6) node[midway,left]{$\vect v_1$};
\draw[vector] (0,1.6)--(1.2,1.6) node[midway,above]{$\vect v_2$};
\draw[vector] (A)--(1.2,1.6);
\node[right] at (1.55,.65){$\vect v_{\text{岸}}$};\draw[thin] (1.55,.78)--(.82,1.02);
\draw[<->] (3.5,0)--(3.5,3) node[midway,left]{$d$};
\node at (1.7,-.8){(a) 最短时间的实际航迹};
\end{scope}
\begin{scope}[xshift=7cm,yshift=.35cm]
\coordinate (W) at (2.5,0);\coordinate (T) at (1.6,1.2);
\draw[kaxis] (-.3,0)--(4.4,0) node[right]{$v_x$};
\draw[kaxis] (0,-1.75)--(0,2.3) node[above]{$v_y$};
\draw (W) circle(1.5);\draw[kguide] (W)--(T);
\draw[vector] (0,0)--(T);\fill (W) circle(1.5pt) (T) circle(1.5pt);
\node[below left] at (0,0){$O$};\node[below] at (W){$W$};\node[above] at (T){$T$};
\node[above right] at (2.2,.65){$v_1$};
\draw[<->] (0,-1.9)--(2.5,-1.9) node[midway,below]{$v_2$};
\draw (.4,0) arc(0:36.87:.4);\node at (.67,.2){$\psi$};
\node at (1.8,-2.65){(b) $v_1<v_2$ 的速度圆};
\end{scope}
\end{tikzpicture}\figcap{例2.82解图：左图出发点在岸上，航迹平行于合速度；右图 $WT\perp OT$，切线给出最小位移方向。}\end{minipage}\end{center}}

\newcommand{\AircraftEqualSpeedDiagram}{%
\begin{center}\begin{minipage}{\linewidth}\centering\begin{tikzpicture}[>=Latex,scale=1]
\foreach \s/\shift/\txt in {1/0/东,-1/6/西}{
\begin{scope}[xshift=\shift cm]
\coordinate (O) at (0,0);\coordinate (G) at (0,3.6);
\coordinate (A) at ({\s*1.74284},3.15);\coordinate (F) at ($(A)!.5!(G)$);
\draw[vector] (O)--(A);\draw[vector] (A)--(G);\draw[vector] (O)--(G);
\draw[kguide] (O)--(F);
\node[below] at (O){$O$};\node[above] at (G){$G$};
\ifnum\s>0 \node[right] at (A){$A$};\node[left] at (0,1.45){$60$};
\else \node[left] at (A){$A$};\node[right] at (0,1.45){$60$};\fi
\node[above] at (F){$F$};\node at ({\s*1.35},1.6){$60$};
\draw (90:1.4) arc[start angle=90,end angle={90-\s*28.955},radius=1.4];
\node at ({\s*.61},1.77){$\theta$};
\node at ({\s*.9},3.9){$30$};
\node at (0,-.65){机头北偏\txt};
\end{scope}}
\node at (3,2.1){正北 $\uparrow$};
\end{tikzpicture}\figcap{例2.83解图：$OA$ 为空速，$AG$ 为风速，$OG$ 为地速；$F$ 为风速边中点，速度单位均为 $\mathrm{m/s}$。}\end{minipage}\end{center}}

\newcommand{\SmokeWindDiagram}{%
\begin{center}\begin{minipage}{\linewidth}\centering\begin{tikzpicture}[>=Latex,scale=.95]
\foreach \vv/\shift/\txt in {1/0/第一次,3/6/第二次}{
\begin{scope}[xshift=\shift cm]
\draw[kaxis] (-.3,0)--(3.3,0) node[right]{东};\draw[kaxis] (0,-1.6)--(0,2.6) node[above]{北};
\draw[vector] (0,0)--(2,-1) node[midway,below left]{$\vect w$};
\draw[vector] (2,-1)--(2,{\vv-1});
\node[right] at (2.1,{.5*\vv-1}){$-\vect v_{\text{船}}$};
\draw[vector] (0,0)--(2,{\vv-1});
\node[above left] at (1,{.5*(\vv-1)+.2}){$\vect v_{\text{烟/船}}$};
\node at (1.4,-2.1){\txt};
\end{scope}}
\end{tikzpicture}\figcap{例2.84解图：同一风速加上两种反向船速，分别得到向东、向东北的烟对船速度。}\end{minipage}\end{center}}

\newcommand{\InterceptTriangleDiagram}{%
\begin{center}\begin{minipage}{\linewidth}\centering\begin{tikzpicture}[>=Latex,x=.045cm,y=.045cm]
\coordinate (A) at (0,0);\coordinate (B) at (-93.9693,34.2020);
\coordinate (C1) at (-17.3065,34.2020);\coordinate (C2) at (79.9527,34.2020);
\draw[thick] (A)--(B);\draw[vector] (B)--(C2);
\draw[vector] (A)--(C1);\draw[vector] (A)--(C2);
\foreach \p/\lab/\pos in {A/A/below,B/B/above left,C1/C_1/above,C2/C_2/above right}{\fill (\p) circle(1.6pt);\node[\pos] at (\p){$\lab$};}
\node[left] at (-50,12){$L$};\node[above] at (30,38){敌船沿 $+x$};
\node[left] at (-8,20){$v_2t_1$};\node[below right] at (42,17){$v_2t_2$};
\draw (B)++(0:22) arc(0:-20:22);\node at (-64,29){$\alpha$};
\draw (23.16018:14) arc(23.16018:160:14);\node[below] at (0,-12){$\beta_1=\angle BAC_1,\quad\beta_2=\angle BAC_2$};
\end{tikzpicture}\figcap{例2.85解图：按文中数值例得到的两个未来交点；发射方向由各自的 $AC_i$ 决定，而不只由正弦值决定。}\end{minipage}\end{center}}

\newcommand{\BombInterceptDiagram}{%
\begin{center}\begin{minipage}{\linewidth}\centering\begin{tikzpicture}[>=Latex,scale=.92]
\draw[kaxis] (-5,0)--(3.8,0) node[right]{$x$};
\coordinate (B) at (-4,3);\coordinate (Q) at (-2,2.25);
\draw[kcurve,domain=0:1,samples=60,variable=\t] plot({-4+4*\t},{3*(1-\t*\t)});
\draw[kguide] (B)--(-4,0);\draw[vector] (B)--(-2.8,3) node[midway,above]{$v_1$};
\draw[thick,dashed] (-2,0)--(Q) (2,0)--(Q);
\foreach \x/\lab/\pos in {0/0/below,-2/-D/below,2/+D/below}{\fill (\x,0) circle(1.6pt);\node[\pos] at (\x,0){$\lab$};}
\node[above left] at (B){$(-L,H)$};\fill (Q) circle(1.6pt);\node[above right] at (Q){$Q$};
\draw[<->] (-4.65,0)--(-4.65,3) node[midway,left]{$H$};
\draw[<->] (-4,-.8)--(0,-.8) node[midway,below]{$L$};
\node at (1.6,3){两炮位分别讨论};
\end{tikzpicture}\figcap{例2.86解图：蓝线是炸弹轨迹；两条虚线仅连接炮位与相遇点，不是炮弹的抛物线轨迹。}\end{minipage}\end{center}}

\newcommand{\ThreeBoatsDiagram}{%
\begin{center}\begin{minipage}{\linewidth}\centering\begin{tikzpicture}[>=Latex,scale=.95]
\begin{scope}
\draw[kguide] (0,0)--(4.8,0);
\foreach \x/\lab in {0/A,1.8/B,3.6/C}{\fill (\x,0) circle(1.5pt);\node[below] at (\x,0){$\lab$};}
\draw[vector] (0,0)--(.6,0);\node[above] at (.3,0){$u$};
\draw[vector] (1.8,0)--(1.8,1.2) node[above]{$2u$};
\draw[vector] (3.6,0)--(4.8,1.2) node[above]{$2\sqrt2u$};
\draw[<->] (0,-.65)--(1.8,-.65) node[midway,below]{$a$};
\draw[<->] (1.8,-.65)--(3.6,-.65) node[midway,below]{$a$};
\node at (2,2.1){(a) $t=0$，速度同比例};
\end{scope}
\begin{scope}[xshift=7cm]
\coordinate (M) at (2.7,2.2);\coordinate (A) at (0,0);
\draw[kguide] (A)--(2.7,0)--(M);\draw[thick,->] (M)--(A);
\fill (M) circle(1.5pt) (A) circle(1.5pt);
\node[above] at (M){$M(t)$};\node[left] at (A){$A(t)$};
\node[below] at (1.35,0){$3a/2$};\node[right] at (2.7,1){$2ut$};
\draw[vector] (A)--(0,-1.2) node[below]{$\vect v_{A/M}=(0,-2u)$};
\node[above left] at (1.5,1.3){$\vect r_{A/M}$};
\node at (1.5,3.05){(b) 随 $M$ 平移的坐标};
\end{scope}
\end{tikzpicture}\figcap{例2.87解图：$M$ 与 $A$ 的横向速度相同，$MA$ 的水平分量固定、竖直分量随时间增加。}\end{minipage}\end{center}}

\newcommand{\PursuitDiagram}{%
\begin{center}\begin{minipage}{\linewidth}\centering\begin{tikzpicture}[>=Latex,scale=1.25]
\draw[thick] (-.3,0)--(3.5,0);
\draw[kguide] (0,3)--(3.1,3);
\draw[kcurve,domain=1:0,samples=100,variable=\u] plot({3*(2/3-\u+\u^3/3)},{3*(1-\u^2)});
\coordinate (D) at (.625,2.25);\coordinate (P) at (1.1875,3);
\draw[thick] (D)--(P);\draw[kguide] (D)--++(1.3,0);
\draw[vector] (D)--++(.27,.36);\draw[vector] (P)--++(.5,0) node[midway,above]{$v_0$};
\fill (D) circle(1.5pt) (P) circle(1.5pt);
\node[left] at (D){$D$};\node[above left] at (P){$P$};
\node[left] at (.86,2.7){$r$};
\draw (D)++(.38,0) arc(0:53.1301:.38);\node at (1.18,2.52){$\theta$};
\draw[vector] (0,0)--(0,.7) node[left]{$v_1$};
\draw[vector] (0,0)--(.65,0) node[below]{$a_0$};
\draw[<->] (-.65,0)--(-.65,3) node[midway,left]{$H$};
\node[below left] at (0,0){$A$};\node[above right] at (2,3){命中};
\end{tikzpicture}\figcap{例2.88解图：取 $v_1=2v_0$ 展示追击曲线；中途导弹 $D$ 的切线沿同一时刻飞机 $P$ 的视线，初始竖直、末端水平。}\end{minipage}\end{center}}

\newcommand{\SlidingRodDiagram}{%
\begin{center}\begin{minipage}{\linewidth}\centering\begin{tikzpicture}[>=Latex,scale=.88]
\begin{scope}
\coordinate (O) at (0,0);\coordinate (A) at (0,3.6);\coordinate (B) at (4,0);\coordinate (P) at (1.52,2.232);
\draw[thick] (0,4.1)--(O)--(4.9,0);\draw[kcurve] (A)--(B);
\foreach \p/\pos in {A/left,B/below,P/above right}{\fill (\p) circle(1.5pt);\node[\pos] at (\p){$\p$};}
\draw[vector] (A)--++(0,-.8) node[midway,left]{$v_A$};\draw[vector] (B)--++(.8,0) node[above]{$v_B$};
\draw (0,2.9) arc(-90:-41.987:.7);\node at (.42,2.7){$\theta$};
\node[above right] at (.75,3.25){$aL$};\node[above right] at (2.7,1.4){$(1-a)L$};
\end{scope}
\begin{scope}[xshift=7cm]
\draw[kaxis] (0,0)--(3.3,0) node[right]{$x$};\draw[kaxis] (0,0)--(0,4.1) node[above]{$y$};
\draw[kcurve,domain=0:90,samples=60] plot({.38*sqrt(28.96)*sin(\x)},{.62*sqrt(28.96)*cos(\x)});
\coordinate (P) at (1.52,2.232);\fill (P) circle(1.5pt);
\draw[vector] (P)--++(.342,-.62);\node[above right] at (P){$P$};
\node[right] at (2,1.55){$\vect v_P$};
\node[above left] at (0,3.337){$(1-a)L$};\node[below] at (2.045,0){$aL$};
\end{scope}
\end{tikzpicture}\figcap{例2.89解图：左图杆长固定、端点接触墙与地面；右图使用同一 $P$，速度沿四分之一椭圆的切线。}\end{minipage}\end{center}}

\newcommand{\BoatTowDiagram}{%
\begin{center}\begin{minipage}{\linewidth}\centering\begin{tikzpicture}[>=Latex,scale=.93]
\coordinate (O) at (0,2.5);\coordinate (B) at (4,0);
\draw[thick] (-3.8,2.3)--(-.15,2.3)--(-.15,-.4);
\fill[Blue!5] (0,-.7) rectangle (5.3,0);\draw[thin,Blue] (0,0)--(5.3,0);
\draw[thick] (-3.2,2.5) rectangle (-2.2,3);\draw (-3,2.4) circle(.1) (-2.4,2.4) circle(.1);
\draw[kcurve] (-2.2,2.5)--(O)--(B);
\fill (O) circle(1.8pt);\node[above right] at (O){$O$};
\draw[thick] (B)--(5.1,0)--(4.9,-.4)--(4.2,-.4)--cycle;
\fill (B) circle(1.5pt);\node[above right] at (B){$B$};
\draw[vector] (-2.3,3.45)--(-3.4,3.45) node[midway,above]{$v$};
\draw[vector] (4.9,-1)--(3.9,-1) node[midway,below]{$u$};
\node[above] at (-1.2,2.5){$l_c$};\node[above right] at (1.9,1.32){$l$};
\draw[<->] (-4.25,0)--(-4.25,2.5) node[midway,left]{$h$};
\draw[kguide] (-4.25,0)--(0,0) (-4.25,2.5)--(-3.5,2.5);
\draw (B)++(148:.9) arc(148:180:.9);\node at (2.82,.31){$\theta$};
\draw[<->] (0,-.85)--(4,-.85) node[midway,below]{$x$};
\end{tikzpicture}\figcap{例2.90解图：$O$ 表示半径忽略的定滑轮转向点；汽车侧增长 $\dot l_c=v$，船侧缩短 $\dot l=-v$。}\end{minipage}\end{center}}

\newcommand{\MovingConstraintDiagram}{%
\begin{center}\begin{minipage}{\linewidth}\centering\begin{tikzpicture}[>=Latex,scale=.9]
\begin{scope}
\draw[thick] (0,0)--(4,0)--(0,2)--cycle;
\coordinate (P) at (2.2,.9);\draw[very thick] (P)--(2.2,3.1);
\foreach \x in {2.06,2.34}{\draw[thick] (\x,2)--(\x,2.6);}
\draw (1.6,2)--(2.06,2) (2.34,2)--(2.8,2);
\fill (P) circle(1.5pt);\node[above right=3pt] at (P){$P$};
\draw[vector] (.6,-.5)--(1.7,-.5) node[midway,below]{$v_0$};
\draw[vector] (2.7,2.65)--(2.7,3.5) node[right]{$v$};
\draw (3.2,0) arc(180:153.435:.8);\node at (2.95,.24){$\theta$};
\node at (1.8,3.9){(a) 斜面与固定导向};
\end{scope}
\begin{scope}[xshift=7cm]
\coordinate (C) at (0,0);\coordinate (P) at (1.2,1.6);
\draw[thick] (-2,0)--(2,0) arc(0:180:2);
\draw[kguide] (C)--(0,2.2) (C)--(P) (P)--(1.2,0);
\draw[very thick] (P)--(1.2,3.2);
\foreach \x in {1.06,1.34}{\draw[thick] (\x,2.4)--(\x,3);}
\draw (.6,2.4)--(1.06,2.4) (1.34,2.4)--(1.8,2.4);
\fill (P) circle(1.5pt) (C) circle(1.5pt);
\node[below left] at (C){$C=(X,0)$};\node[above right=3pt] at (P){$P$};
\node at (.7,.46){$R$};\node[below] at (1.2,0){$x_0$};
\draw (0,.75) arc(90:53.1301:.75);\node at (.29,1.13){$\theta$};
\draw[vector] (1.8,2.7)--(1.8,3.5) node[right]{$v_y$};
\draw[vector] (.5,-.8)--(1.6,-.8) node[midway,below]{$v_0$};
\node at (0,3.9){(b) 半圆柱与固定导向};
\end{scope}
\end{tikzpicture}\figcap{例2.91—2.92解图：两侧固定导向只允许杆竖直平移；接触点由同一斜面方程或圆方程确定。}\end{minipage}\end{center}}

\newcommand{\RodPulleyDiagram}{%
\begin{center}\begin{minipage}{\linewidth}\centering\begin{tikzpicture}[>=Latex,scale=1]
\coordinate (O) at (0,0);\coordinate (B) at (0,3.4);\coordinate (A) at (2.3,1.2);\coordinate (C) at (-2.2,3.4);
\draw[very thick] (O)--(A) node[midway,below=5pt]{$R$};
\draw[kcurve] (A)--(B)--(C)--(-2.2,.65);
\draw[thick] (-2.55,0) rectangle (-1.85,.65);\node at (-2.2,.32){$M$};
\draw[kguide] (O)--(B);
\foreach \p/\pos in {O/below,B/above,C/above,A/right}{\fill (\p) circle(1.7pt);\node[\pos] at (\p){$\p$};}
\draw[<->] (-.65,0)--(-.65,3.4) node[midway,left]{$H$};
\draw (B)++(-90:.8) arc(-90:-43.727:.8);\node at (.42,2.4){$\alpha$};
\draw[vector] (A)--++(-.48,.92) node[above]{$\vect v_A$};
\draw[->,DeepBlue] (35:.75) arc(35:120:.75);\node[above] at (.3,1.2){$\omega$};
\draw[vector] (-2.85,.6)--(-2.85,-.3) node[below]{$v_M$};
\node[left] at (.83,2.0){$l$};
\end{tikzpicture}\figcap{例2.93解图：零半径转向点 $B,C$ 与绳折点重合；正角速度下 $AB$ 缩短而物块下降。}\end{minipage}\end{center}}

\newcommand{\EulerStepsDiagram}{%
\begin{center}\begin{minipage}{\linewidth}\centering\begin{tikzpicture}[>=Latex,scale=.87]
\begin{scope}
\draw[kaxis] (-.2,0)--(2,0) node[right]{$x$};\draw[kaxis] (0,-.2)--(0,2) node[above]{$y$};
\draw[vector] (0,0)--(35:2) node[above right]{$N$};
\draw[->] (.85,0) arc(0:35:.85);\node at (17.5:1.16){$\alpha$};
\node[below left] at (0,0){$O$};\node at (.7,-.85){(1) 绕 $z$};\node at (.7,-1.35){从 $+z$ 端看};
\end{scope}
\begin{scope}[xshift=5cm]
\draw[TextGray,thick] (-1.6,0)--(1.7,0) node[right]{$xy$};
\draw[DeepBlue,thick] (215:1.7)--(35:1.7) node[right]{$XY$};
\draw[kaxis] (0,0)--(90:2.1) node[above]{$z$};
\draw[vector] (0,0)--(125:2.1) node[above left]{$Z$};
\draw[->] (90:1.05) arc(90:125:1.05);\node at (107.5:1.45){$\beta$};
\draw (0,0) circle(.085);\fill (0,0) circle(.02);
\node[right] at (.12,-.4){$N$ 向外};
\node at (0,-1.4){(2) 绕 $N$};\node at (0,-1.9){从 $+N$ 端看};
\end{scope}
\begin{scope}[xshift=9.8cm]
\draw[kguide] (-.2,0)--(2,0) node[right]{$N$};
\draw[vector] (0,0)--(35:2) node[above right]{$X$};
\draw[vector] (0,0)--(125:1.9) node[above left]{$Y$};
\draw[->] (.85,0) arc(0:35:.85);\node at (17.5:1.16){$\gamma$};
\node[below left] at (0,0){$O$};\node at (.5,-.85){(3) 绕 $Z$};\node at (.5,-1.35){从 $+Z$ 端看};
\end{scope}
\end{tikzpicture}\figcap{图2.20：右手正转的三步欧拉角；第二图中 $XY$ 截线转过 $\beta$，其法线 $Z$ 同转过 $\beta$，二者始终垂直。}\end{minipage}\end{center}}

\newcommand{\RigidMotionTypesDiagram}{%
\begin{center}\begin{minipage}{\linewidth}\centering\begin{tikzpicture}[>=Latex,scale=.9]
\begin{scope}
\foreach \dy in {0,1}{\draw[kguide,domain=0:1,samples=45,variable=\t] plot({3*\t},{\dy+.9*sin(180*\t)});}
\draw[thick] (-.3,-.25) rectangle (.3,1.25);
\draw[thick] (2.7,-.25) rectangle (3.3,1.25);
\foreach \x in {0,3}{\foreach \y in {0,1}{\fill (\x,\y) circle(1.5pt);}}
\draw[vector] (0,0)--(.6,.565);\draw[vector] (0,1)--(.6,1.565);
\node[left] at (-.45,0){$P_1$};\node[left] at (-.45,1){$Q_1$};
\node[right] at (3.45,0){$P_2$};\node[right] at (3.45,1){$Q_2$};
\node at (1.5,2.4){(a) 曲线平动};
\node at (1.5,-1){两姿态的对应边平行};
\end{scope}
\begin{scope}[xshift=8cm]
\draw[kguide] (0,-1.2)--(0,1.2);
\draw[thick] (0,-2)--(0,-1.2) node[pos=0,below]{$B$};
\draw[kaxis] (0,1.2)--(0,2.4) node[above]{$A$};
\draw[thick,dashed] (2,0) arc(0:180:2 and 1.2);
\draw[thick] (-2,0) arc(180:360:2 and 1.2);
\coordinate (P) at ({2*cos(315)},{1.2*sin(315)});
\draw[thick] (0,0)--(P);\draw[kguide] (0,0)--(2,0);
\draw[vector] (P)--++(.78,.468);
\draw (.7,0) arc(0:-45:.7 and .42);
\node at (1,-.19){$\theta$};\node at (.46,-.62){$R$};
\fill (P) circle(1.5pt);\node[below] at (P){$P$};\node[left] at (0,0){$O$};
\node[above right] at (2.2,-.38){$\vect v_P$};
\node at (0,-2.7){(b) 定轴转动的圆平面投影};
\end{scope}
\end{tikzpicture}\figcap{图2.21—图2.22：平动可以沿曲线；定轴转动中 $OP=R$，点 $P$ 的切向速度为 $|\omega|R$。}\end{minipage}\end{center}}

\newcommand{\RigidBarCasesDiagram}{%
\begin{center}\begin{minipage}{\linewidth}\centering\begin{tikzpicture}[>=Latex,scale=.95]
\begin{scope}[xscale=1.5]
\draw[thick] (0,0)--(2.4,0);\draw[kguide] (2.4,0)--(4.8,0);
\draw[vector] (0,0)--(0,1.2) node[above]{$v_0$};\draw[vector] (2.4,0)--(2.4,.6) node[above]{$v_0/2$};
\foreach \x/\lab in {0/A,2.4/B,4.8/I}{\fill (\x,0) circle(1.5pt);\node[below] at (\x,0){$\lab$};}
\fill (1.2,0) circle(1.5pt);\node[above] at (1.2,0){$G$};
\draw[<->] (0,-.7)--(2.4,-.7) node[midway,below]{$l$};\draw[<->] (2.4,-.7)--(4.8,-.7) node[midway,below]{$l$};
\node at (2.4,2){(1) $x_G=l/2,\ x_I=2l$};
\end{scope}
\begin{scope}[yshift=-4.1cm,xscale=1.5]
\draw[thick] (0,0)--(2.4,0);\draw[vector] (0,0)--(0,1.2) node[above]{$v_0$};\draw[vector] (2.4,0)--(2.4,-1.2) node[below]{$-v_0$};
\node[left] at (0,0){$A$};\node[right] at (2.4,0){$B$};\fill (1.2,0) circle(1.6pt);\node[above] at (1.2,0){$G=I$};
\node at (1.2,2){(2) 等质量};
\end{scope}
\begin{scope}[xshift=7cm,yshift=-4.1cm]
\draw[thick] (0,0)--(3.6,0);\draw[vector] (0,0)--(0,1.2) node[above]{$v_0$};\draw[vector] (3.6,0)--(3.6,-1.2) node[below]{$-v_0$};
\node[left] at (0,0){$A$};\node[right] at (3.6,0){$B$};
\fill (1.2,0) circle(1.6pt) (1.8,0) circle(1.6pt);\node[above] at (1.2,0){$G$};\node[below] at (1.8,0){$I$};
\node at (1.8,2){(3) $m_A=2m_B$};
\draw[<->] (0,-.65)--(1.2,-.65) node[midway,below]{$l/3$};
\end{scope}
\end{tikzpicture}\figcap{例2.94解图：三问分别标出实际速度、质心 $G$ 与瞬心 $I$；质量改变不改变给定速度场的瞬心。}\end{minipage}\end{center}}

\newcommand{\RollingBallDiagram}{%
\begin{center}\begin{minipage}{\linewidth}\centering\begin{tikzpicture}[>=Latex,scale=1]
\draw[thick] (-2,0)--(4,0);\draw[thick] (0,1) circle(1);
\fill (0,1) circle(1.5pt) (0,0) circle(1.5pt);\node[below left] at (0,1){$C$};\node[below left] at (0,0){$P$};
\draw (0,1)--(-.707,1.707);\node[left] at (-.58,1.28){$R$};
\draw[vector] (0,1)--(2.6,1);\node[above] at (2,1){$v_0$};
\draw[->,DeepBlue,thick] (40:1.5)++(0,1) arc(40:145:1.5);
\node[above] at (0,2.65){$\omega_0>0$（逆时针）};
\draw[vector] (0,-.65)--(2.5,-.65) node[midway,below]{$v_P=v_0+R\omega_0$};
\end{tikzpicture}\figcap{例2.95解图：接触点的平动速度和转动速度初始都向右；纯滚动要求它们相消。}\end{minipage}\end{center}}

% Knife local coordinates in metres: handle end -0.15, COM 0, tip +0.10.
\newcommand{\KnifeOutline}{%
\draw[thick,fill=TextGray!16] (-.15,-.012) rectangle (-.045,.012);
\draw[thick,fill=Blue!8] (-.045,-.012)--(.10,0)--(-.045,.012)--cycle;}

\newcommand{\KnifeFlightDiagram}{%
\begin{center}\begin{minipage}{\linewidth}\centering\begin{tikzpicture}[>=Latex]
\begin{scope}[x=2.1cm,y=2.1cm]
\draw[thick] (5,-.1)--(5,1.1);\node[right] at (5,.7){木靶};
\draw[kguide,domain=0:1,samples=100,variable=\t] plot({4.9*\t},{4.9*tan(30)*\t*(1-\t)});
\foreach \t/\n in {0/1,.25/2,.5/3,.75/4,1/5}{
\begin{scope}[shift={({4.9*\t},{4.9*tan(30)*\t*(1-\t)})},rotate={60+300*\t}]
\KnifeOutline\fill (0,0) circle(.7pt);
\end{scope}
\node[above] at ({4.9*\t},{4.9*tan(30)*\t*(1-\t)+.17}){\n};}
\draw[<->] (0,-.35)--(5,-.35) node[midway,below]{$s=5.0\,\mathrm m$};
\node at (2.5,-.8){(a) 质心抛物线与五个等时姿态};
\end{scope}
\begin{scope}[xshift=2.1cm,yshift=-3.8cm]
\begin{scope}[scale=15]\KnifeOutline\end{scope}
\fill (0,0) circle(1.5pt);\node[below] at (0,-.25){$C$};\node[right] at (1.5,0){刀尖};
\draw[<->] (-2.25,-.65)--(1.5,-.65) node[midway,below]{$L=25\,\mathrm{cm}$};
\draw[<->] (0,.65)--(1.5,.65) node[midway,above]{$l=10\,\mathrm{cm}$};
\node at (0,-1.45){(b) 长度定义（局部放大）};
\end{scope}
\begin{scope}[xshift=8cm,yshift=-3.8cm]
\draw[kaxis] (-.3,0)--(2.5,0) node[right]{$x$};
\begin{scope}[rotate=60,scale=10]\KnifeOutline\end{scope}
\draw[vector] (0,0)--(30:2.1) node[above right]{$\vect v_C$};
\draw (.65,0) arc(0:30:.65);\node at (15:.95){$30^\circ$};
\draw (60:1.85) arc(60:360:1.85);\draw[->] (335:1.85) arc(335:360:1.85);
\node[left] at (-1.85,.45){$300^\circ$};
\node[right] at (.25,-.45){$C$};\draw[thin] (.25,-.45)--(0,0);
\node[above] at (60:2.05){$\theta=60^\circ$};
\node at (.3,-2.45){(c) 初姿态与最小正转角};
\end{scope}
\end{tikzpicture}\figcap{例2.96解图：尖端在刀轴正向；第5姿态刀尖恰达木靶，质心仍距靶 $l$，不是整段刀身穿靶。}\end{minipage}\end{center}}

\newcommand{\KnifeImpactDiagram}{%
\begin{center}\begin{minipage}{\linewidth}\centering\begin{tikzpicture}[>=Latex]
\begin{scope}[x=2cm,y=2cm]
\draw[kguide] (0,0)--(5,0);\draw[thick] (5,-.15)--(5,.5);
\begin{scope}[shift={(4.9,0)}]\KnifeOutline\end{scope}
\fill (0,0) circle(1.5pt);\node[above] at (0,0){$C_0$};\node[above] at (4.9,.15){$C_1$};
\draw[<->] (0,-.35)--(4.9,-.35) node[midway,below]{$d_1=s-l$};
\node at (2.5,.65){(a) 第一把刀击中靶点 $Q$};
\end{scope}
\begin{scope}[xshift=5cm,yshift=-3.9cm,x=15cm,y=15cm]
\begin{scope}[shift={(.10,0)}]\KnifeOutline\end{scope}
\begin{scope}[shift={(-.15,0)}]\KnifeOutline\end{scope}
\draw[thick] (.20,-.08)--(.20,.09);
\fill (-.05,0) circle(1.4pt) (-.15,0) circle(1.4pt) (.10,0) circle(1.4pt);
\node[above] at (-.05,.025){$P$};\node[below] at (-.15,-.015){$C_2$};\node[below] at (.10,-.015){$C_1$};\node[right] at (.20,0){$Q$};
\draw[<->] (-.05,.065)--(.20,.065) node[midway,above]{$L$};
\draw[<->] (-.15,-.085)--(-.05,-.085) node[midway,below]{$l$};
\node at (-.06,-.15){(b) 接触局部放大：$C_2P=l,\ PQ=L$};
\end{scope}
\end{tikzpicture}\figcap{例2.96碰撞几何：两刀仅在 $P$ 接触；第一把木柄远端就是第二把刀尖，故 $C_0C_2=s-L-l$。}\end{minipage}\end{center}}

\newcommand{\LightTimingDiagram}{%
\begin{center}\begin{minipage}{\linewidth}\centering\begin{tikzpicture}[>=Latex,scale=.95]
\coordinate (A) at (-3,0);\coordinate (B) at (3,0);\coordinate (O) at (-.7,2.5);
\draw[thick] (A)--(B);\draw[kguide] (A)--(O)--(B);
\foreach \p/\pos in {A/below left,B/below right,O/above}{\fill (\p) circle(1.5pt);\node[\pos] at (\p){$\p$};}
\node[above left] at (-1.95,1.2){$r_A$};\node[above right] at (1.4,1.2){$r_B$};
\draw[vector] (-.8,-.7)--(.8,-.7) node[midway,below]{运动员};
\node[above] at (0,0){$L=100\,\mathrm m$};
\end{tikzpicture}\figcap{例2.102解图：观察者 $O$ 接收起点 $A$、终点 $B$ 的光信号；传播路程一般不相等。}\end{minipage}\end{center}}

\newcommand{\ApparentTransverseDiagram}{%
\begin{center}\begin{minipage}{\linewidth}\centering\begin{tikzpicture}[>=Latex,scale=1]
\coordinate (E) at (0,0);\coordinate (P1) at (0,4);\coordinate (P2) at (2,2.6);\coordinate (F) at (0,2.6);
\draw[thick] (E)--(P1);\draw[thick] (E)--(P2);
\draw[vector] (P1)--(P2);\draw[kguide] (F)--(P2);
\fill (E) circle(1.6pt) (P1) circle(1.6pt) (P2) circle(1.6pt);
\node[below] at (E){地球};\node[above] at (P1){$P_1$};\node[right] at (P2){$P_2$};
\draw[<->] (-.9,0)--(-.9,4) node[midway,left]{$L$};
\node[above right] at (.85,3.55){$v\Delta t$};\node[below right] at (1.15,1.4){$D_2$};
\draw[kguide] (0,4.35)--(0,4.85) (P2)--(2,4.85);
\draw[<->] (0,4.85)--(2,4.85) node[midway,above]{$v\sin\theta\,\Delta t$};
\draw[<->] (-.4,2.6)--(-.4,4);\node[left,align=right] at (-1.3,3.3){接近量\\$v\cos\theta\,\Delta t$};\draw[thin] (-1.25,3.3)--(-.42,3.3);
\draw (0,3.35) arc(-90:-34.992:.65);\node at (.3,3.1){$\theta$};
\end{tikzpicture}\figcap{例2.103解图：横向与视线方向的位移分开标注；三角形给出 $D_2$，再对 $v\Delta t/L$ 作一阶展开。}\end{minipage}\end{center}}

\newcommand{\MarsRetrogradeDiagram}{%
\begin{center}\begin{minipage}{\linewidth}\centering\begin{tikzpicture}[>=Latex,scale=.9]
\foreach \t/\n/\dx/\dy in {-90/1/0/0,-60/2/3.5/0,-36.37/3/7/0,0/4/10.5/0,36.37/5/1.75/-4,60/6/5.25/-4,90/7/8.75/-4}{
\begin{scope}[shift={(\dx,\dy)}]
\draw[TextGray] (0,0) circle(.8);\draw[thick] (0,0) circle({.8*1.524});
\fill (0,0) circle(1.3pt);\node[above] at (0,0){$S$};
\pgfmathsetmacro{\ea}{.985609*\t-90}
\pgfmathsetmacro{\ma}{.985609*\t*pow(1.524,-1.5)-90}
\pgfmathsetmacro{\ex}{.8*cos(\ea)}\pgfmathsetmacro{\ey}{.8*sin(\ea)}
\pgfmathsetmacro{\mx}{.8*1.524*cos(\ma)}\pgfmathsetmacro{\my}{.8*1.524*sin(\ma)}
\draw[DeepBlue,dashed,thin,->] (\ex,\ey)--({\ex+1.7*(\mx-\ex)},{\ey+1.7*(\my-\ey)});
\fill (\ex,\ey) circle(1.5pt) (\mx,\my) circle(1.5pt);
\node[font=\scriptsize] at (\ea:.5){$E_{\n}$};
\ifdim\t pt<0pt \node[font=\scriptsize,left] at ({\mx-.13},{\my-.16}){$M_{\n}$};
\else \node[font=\scriptsize,right] at ({\mx+.13},{\my-.16}){$M_{\n}$};\fi
\node[font=\footnotesize] at (0,1.7){\n：$t=\t\,\mathrm d$};
\end{scope}}
\node at (5.25,-6){虚线由地球经过火星，指向背景星空；各小图比例相同};
\begin{scope}[xshift=5.25cm,yshift=-8.6cm]
\draw[kaxis] (-3.7,0)--(3.8,0) node[right]{$t/\mathrm d$};
\draw[kaxis] (0,-1.9)--(0,1.45) node[above]{$\lambda$};
\MarsCoplanarPath
\foreach \t/\n/\pos in {-90/1/above left,-60/2/above left,-36.37/3/above,0/4/below left,36.37/5/below,60/6/below right,90/7/below right}{
\pgfmathsetmacro{\ll}{atan2(1.524*sin(.985609*\t*pow(1.524,-1.5))-sin(.985609*\t),1.524*cos(.985609*\t*pow(1.524,-1.5))-cos(.985609*\t))/12}
\fill (\t/27,\ll) circle(1.5pt);\node[\pos] at (\t/27,\ll){\n};}
\draw[<->] (-1.347,-1.75)--(1.347,-1.75) node[midway,below]{逆行};
\node at (0,2.15){同平面模型：七个时刻的视经度};
\node at (0,-2.65){第3、5点：$\dot\lambda=0$};
\end{scope}
\end{tikzpicture}\figcap{例2.105解图：七组位置与七个视方向一一对应；冲日取 $t=0$，第3至第5点之间视经度随时间减小。}\end{minipage}\end{center}}

\newcommand{\MarsInclinedDiagram}{%
\begin{center}\begin{minipage}{\linewidth}\centering\begin{tikzpicture}[>=Latex]
\draw[kaxis] (-4.4,-2.1)--(4.6,-2.1) node[right]{视经度增加};
\draw[kaxis] (-4.4,-2.1)--(-4.4,3.05) node[above]{视纬度增加};
\MarsInclinedPath
\node[above] at (0,3.25){三维圆轨道示例：轨道倾角 $1.85^\circ$};
\node at (0,-2.8){箭头表示时间先后；蓝色区间才是逆行};
\end{tikzpicture}\figcap{回环与计算模型的区别：加入非零倾角后视纬度也会变化；蓝段按视经度减小的条件确定，不把整个回环都称为逆行。}\end{minipage}\end{center}}

\newcommand{\RomerDiagram}{%
\begin{center}\begin{minipage}{\linewidth}\centering\begin{tikzpicture}[>=Latex,scale=1]
\coordinate (S) at (0,0);\coordinate (J) at (5.2,0);
\pgfmathsetmacro{\ang}{acos(1/5.2)}
\draw[thick] (S) circle(1);\draw[TextGray] ({5.2*cos(-10)},{5.2*sin(-10)}) arc(-10:10:5.2);
\foreach \a/\lab/\pos in {\ang/E_+/above left,-\ang/E_-/below left}{
\coordinate (E) at (\a:1);\draw[kguide] (S)--(E);\draw[thick] (E)--(J);\fill (E) circle(1.6pt);\node[\pos] at (E){$\lab$};
\draw[vector] (E)--++({-.65*sin(\a)},{.65*cos(\a)});}
\fill (S) circle(1.6pt) (J) circle(1.6pt);\node[left] at (S){$S$};\node[right] at (J){$J$};
\draw[kguide] (S)--(J);
\draw[<->] (0,-2.25)--(5.2,-2.25) node[midway,below]{$R_J=5.2R_E$};
\node at (2,1.65){$E_+$：距离增加，周期最长};\node at (2,-1.65){$E_-$：距离减小，周期最短};
\begin{scope}[xshift=9cm]
\draw (0,0) circle(.65);\fill (0,0) circle(2pt);\node[below] at (0,0){$J$};
\fill (45:.65) circle(1.6pt);\node[above right] at (45:.65){木卫I};
\draw[->] (50:.9) arc(50:170:.9);\node at (0,-1.2){卫星轨道放大};
\end{scope}
\end{tikzpicture}\figcap{例2.106解图：$SE_\pm\perp E_\pm J$，两切点为周期极值位置；木卫I单独放大，不与日心轨道混用比例。}\end{minipage}\end{center}}

\newcommand{\AimLocusDiagram}{%
\begin{center}\begin{minipage}{\linewidth}\centering\begin{tikzpicture}[>=Latex,scale=.9]
\begin{scope}
\draw[thick] (-.4,0)--(4.1,0);\draw[thick] (3.6,0)--(3.6,4.4);
\coordinate (O) at (0,.6);\coordinate (Q) at (3.6,3.7);\coordinate (P) at (3.6,1.6);
\draw[kguide] (O)--(Q);\draw[kcurve,domain=0:1,samples=70,variable=\t] plot({3.6*\t},{.6+3.1*\t-2.1*\t^2});
\foreach \p in {O,Q,P}{\fill (\p) circle(1.5pt);}
\node[above left=3pt] at (O){$O$};\node[right] at (Q){$Q$};\node[below left] at (3.5,1.45){$P\ (z=h)$};
\draw[<->] (0,-.5)--(3.6,-.5) node[midway,below]{$r$};
\draw[<->] (4.65,.6)--(4.65,3.7) node[midway,right]{$Z$};
\draw[kguide] (0,.6)--(4.65,.6);
\node[left] at (0,.25){$z_0$};
\node at (1.8,5){(a) 包含 $OQ$ 的竖直截面};
\end{scope}
\begin{scope}[xshift=9cm]
\draw[kaxis] (-2.4,0)--(2.6,0) node[right]{$x$};\draw[kaxis] (0,0)--(0,4.6) node[above]{$z$};
\draw[DeepBlue,thick] (0,2.7) circle(1.5);\draw[thick] (-2.4,.7)--(2.4,.7);
\coordinate (Q) at (1.2,3.6);\coordinate (P) at (1.2,.7);
\draw[kguide] (Q)--(P);\fill (Q) circle(1.5pt) (P) circle(1.5pt) (0,2.7) circle(1.5pt);
\node[above right] at (Q){$Q$};\node[below right] at (P){$P$};
\node[right] at (2.45,.7){$h$};
\node[left] at (-.1,2.7){$K$};
\node at (0,5.25){(b) 正视墙面};
\end{scope}
\end{tikzpicture}\figcap{例2.109解图：瞄准点 $Q$ 位于以 $K$ 为圆心的墙面圆上，实际命中点 $P$ 位于水平线；两图分别说明投掷平面与墙面轨迹。}\end{minipage}\end{center}}

% Question figure contract — source physical pages in question-figure-map.json.
% Mode C; simple physics geometry, restrained blue/grey lines on white paper.
% Each question can be understood while hiding its solution. No solved labels.
% Reference geometry: preserve contacts, axis zero, directions, source point identities.
% Main viewport <=14cm; ordinary labels 9–11pt; all labels transparent.
% Read axes/objects first, givens second; no decoration or filled label rectangles.
% Cloned givens retain exact reviewed coordinates; answer overlays are excluded explicitly.
% Captions + source-specific labels bind each view to its own question, never previous answer.
% Anticipated risks: angle labels near rays, multiple dimensions, clipping tall multi-panels.
% Sketch: [question text] [given geometry or blank axes] -> [answer text + construction].

% Source example 2.1; question-only view.
\newcommand{\KQuestionAB}{%
\begin{center}\begin{minipage}{.97\linewidth}\centering
\begin{tikzpicture}[>=Latex,scale=.9]

\draw[kaxis] (-.25,0)--(6.6,0) node[right]{$t/\mathrm s$};
\draw[kaxis] (0,-2.7)--(0,1.2) node[above]{$x/\mathrm m$};
\foreach \t in {1,...,6}{\draw (\t,0)--(\t,.08);\node[above=2pt] at (\t,0){\t};}
\foreach \y/\v in {-2.4/-60,-1.6/-40,-.8/-20,.8/20}{\draw (0,\y)--(.08,\y);\node[left] at (0,\y){$\v$};}
\node[below left] at (0,0){$O$};

\end{tikzpicture}
\figcap{题图：按表中数据在空坐标系中作图。}\label{question-figure-2-1}
\end{minipage}\end{center}
}

% Source example 2.2; question-only view.
\newcommand{\KQuestionAC}{%
\begin{center}\begin{minipage}{.97\linewidth}\centering
\begin{tikzpicture}[>=Latex,scale=.85]

\foreach \dx/\i in {0/1,4.3/2,8.6/3}{\begin{scope}[xshift=\dx cm]
\draw[kaxis] (-.2,0)--(3.1,0) node[right]{$t$};
\draw[kaxis] (0,-.6)--(0,1.8) node[above]{$x$};
\node[below left] at (0,0){$O$};\node at (1.4,-.85){(\i)};
\end{scope}}

\end{tikzpicture}
\figcap{题图：分别画出三个运动函数的图像。}\label{question-figure-2-2}
\end{minipage}\end{center}
}

% Source example 2.3; question-only view.
\newcommand{\KQuestionAD}{%
\begin{center}\begin{minipage}{.97\linewidth}\centering
\begin{tikzpicture}[>=Latex,scale=.9]

\foreach \dx/\mode in {0/1,6.2/2}{\begin{scope}[xshift=\dx cm]
\draw[kaxis] (-.2,0)--(4.6,0) node[right]{$t$};
\draw[kaxis] (0,-1.25)--(0,1.3) node[above]{$x$};
\ifnum\mode=1
\draw[kcurve,domain=0:4.25,samples=200] plot(\x,{.8*sin(280*\x)});
\else
\draw[kcurve,domain=0:4.25,samples=240] plot(\x,{1.05*exp(-.85*\x)*cos(450*\x)});
\fi
\node[below left] at (0,0){$O$};\node at (2,-1.6){(\mode)};
\end{scope}}

\end{tikzpicture}
\figcap{题图：两个质点的一维位置随时间变化，坐标轴未给定数值比例。}\label{question-figure-2-3}
\end{minipage}\end{center}
}

% Source example 2.6; question-only view.
\newcommand{\KQuestionAG}{%
%
\begin{center}\begin{tikzpicture}[scale=.78,>=Latex]
  \foreach \dx/\mode in {-4.6/1,0/2,4.6/3}{
    \begin{scope}[xshift=\dx cm]
      \draw[kaxis] (-1.65,0)--(1.75,0) node[right,clearlabel]{$t$};
      \draw[kaxis] (-1.45,-.55)--(-1.45,2.05) node[above,clearlabel]{$x$};
      \ifnum\mode=1
        \draw[kcurve] (-1.35,0)--(0,1.35)--(1.35,0);
      \else\ifnum\mode=2
        \draw[kcurve,domain=-1.35:1.30,samples=80] plot (\x,{.25*(\x+1.35)^2});
      \else
        \draw[kcurve,domain=-1.35:1.35,samples=100] plot (\x,{1.45*sqrt(max(0,1-(\x/1.35)^2))});
      \fi\fi
    \end{scope}}
\end{tikzpicture}\figcap{题图：从左到右为三幅给定的位置—时间图像。}\label{question-figure-2-6}\end{center}
}

% Source example 2.7; question-only view.
\newcommand{\KQuestionAH}{%
%
\begin{center}\begin{tikzpicture}[scale=.9,>=Latex]
  \begin{scope}[shift={(-4.9,-1.1)}]
    \draw[kaxis] (0,0)--(2.75,0) node[right,clearlabel]{$t$};
    \draw[kaxis] (0,-1.55)--(0,1.85) node[above,clearlabel]{$v$};
    \draw[kcurve] (0,1.3)--(2.4,-1.3);
    \node[clearlabel,left] at (0,1.3){$v_0$};
    \node[clearlabel,left] at (0,-1.3){$-v_0$};
    \draw[kguide] (0,-1.3)--(2.4,-1.3);
    \node[clearlabel,below left] at (0,0){$O$};
    \node[clearlabel,above right=2pt] at (1.2,0){$T/2$};
    \node[clearlabel,above] at (2.4,0){$T$};
    \node[clearlabel,below] at (1.2,-1.7){已知速度图};
  \end{scope}
  \foreach \dx/\dy/\mode/\tag in {-.7/0/1/A,3.3/0/2/B,-.7/-3.3/3/C,3.3/-3.3/4/D}{
    \begin{scope}[shift={(\dx,\dy)}]
      \draw[kaxis] (0,0)--(2.7,0) node[right,clearlabel]{$t$};
      \draw[kaxis] (0,-.25)--(0,1.65) node[above,clearlabel]{$x$};
      \ifnum\mode=1
        \draw[kcurve] (0,0)--(1.2,1.15)--(2.4,0);
      \else\ifnum\mode=2
        \draw[kcurve,domain=0:1,samples=70]
          plot ({2.4*\x},{4.6*\x*(1-\x)});
      \else\ifnum\mode=3
        \draw[kcurve,domain=0:1,samples=70] plot ({2.4*\x},{1.15*(2*\x-\x*\x)});
      \else
        \draw[kcurve] (0,1.15)--(1.2,0)--(2.4,.85);
      \fi\fi\fi
      \node[clearlabel,below left] at (0,0){$O$};
      \node[clearlabel,below=3pt] at (1.2,0){$T/2$};
      \node[clearlabel,below=3pt] at (2.4,0){$T$};
      \node[clearlabel,below] at (1.2,-.65){$\tag$};
    \end{scope}}
\end{tikzpicture}\figcap{题图：左图给出速度随时间的变化，右侧 A—D 为待选的位置图像。}\label{question-figure-2-7}\end{center}
}

% Source example 2.8; question-only view.
\newcommand{\KQuestionAI}{%
%
\begin{center}\begin{tikzpicture}[scale=.78,>=Latex]
  \foreach \dx/\mode in {-4.4/1,0/2,4.4/3}{
    \begin{scope}[xshift=\dx cm]
      \draw[kaxis] (-1.55,0)--(1.75,0) node[right,clearlabel]{时间};
      \draw[kaxis] (-1.35,-.25)--(-1.35,1.85) node[above,clearlabel]{$v$};
      \ifnum\mode=1
        \draw[kcurve] (-1.25,0)--(1.25,1.35);
        \draw[kguide] (1.25,0)--(1.25,1.35);
        \draw[kguide] (-1.35,1.35)--(1.25,1.35);
        \node[clearlabel,left] at (-1.35,1.35){$v$};
      \else\ifnum\mode=2
        \draw[kcurve] (-1.25,.55)--(1.25,1.35);
        \draw[kguide] (1.25,0)--(1.25,1.35);
        \draw[kguide] (-1.35,.55)--(-1.25,.55) (-1.35,1.35)--(1.25,1.35);
        \node[clearlabel,left] at (-1.35,.55){$v_0$};
        \node[clearlabel,left] at (-1.35,1.35){$v_1$};
      \else
        \draw[kcurve] (-1.25,0)--(-.42,1.35)--(1.25,0);
        \draw[kguide] (-.42,0)--(-.42,1.35);
        \draw[kguide] (-1.35,1.35)--(-.42,1.35);
        \node[clearlabel,left] at (-1.35,1.35){$v_0$};
        \node[clearlabel,below] at (-.42,0){$t/3$};
      \fi\fi
      \node[clearlabel,below] at (1.25,0){$t$};
    \end{scope}}
\end{tikzpicture}\figcap{题图：三个运动过程的速度—时间图像，求相应时间内的路程。}\label{question-figure-2-8}\end{center}
}

% Source example 2.17; question-only view.
\newcommand{\KQuestionAR}{%
%
\begin{center}\begin{tikzpicture}[scale=.9,>=Latex]
  \draw[very thick] (-1.4,0)--(1.4,0);
  \node[point,label={[clearlabel]below right=3pt:$B$}] (B) at (0,0) {};
  \node[point,label={[clearlabel]left:$A$}] (A) at (0,3.2) {};
  \draw[<->,TextGray] (.8,0)--(.8,3.2) node[midway,clearlabel,right]{$h$};
  \draw[DeepBlue,very thick,->] (B)--(0,1.15) node[midway,clearlabel,left]{$v_0$};
  \draw[Blue,very thick,->] (A)--(0,2.45) node[midway,clearlabel,left]{$g$};
\end{tikzpicture}\figcap{题图：A 球从上方释放，B 球从下方以初速度 $v_0$ 竖直抛出，初始间距为 $h$。}\label{question-figure-2-17}\end{center}
}

% Source example 2.19; question-only view.
\newcommand{\KQuestionAT}{%
%
\begin{center}\begin{tikzpicture}[scale=.85,>=Latex]
  \draw[very thick] (-1.7,0)--(1.7,0);
  \fill (0,2.7) circle (.12);
  \fill (0,.12) circle (.12);
  \draw[kguide] (0,.27)--(0,2.55);
  \draw[<->,TextGray] (.88,0)--(.88,2.7)
    node[midway,clearlabel,right=2pt]{$H$};
  \draw[DeepBlue,very thick,->] (-.35,1.0)--(-.35,.25) node[midway,clearlabel,left]{$u$};
  \draw[Blue,very thick,->] (.35,.25)--(.35,1.0) node[midway,clearlabel,right]{$eu$};
\end{tikzpicture}\figcap{题图：初始高度为 $H$；一次碰撞前、后的速率分别用 $u$、$eu$ 表示。}\label{question-figure-2-19}\end{center}
}

% Source example 2.22; question-only view.
\newcommand{\KQuestionAW}{%
\begin{center}\begin{minipage}{.97\linewidth}\centering
\begin{tikzpicture}[>=Latex,scale=.70]

\foreach \x in {-4,-3,-2,-1,1,2,3,4,5}{\draw[RuleGray,thin] (\x,-5.5)--(\x,-.8) (\x,.1)--(\x,5);}
\foreach \y in {-5,-4,-3,-2,-1,1,2,3,4,5}{\draw[RuleGray,thin] (-4,\y)--(-1.05,\y) (.12,\y)--(5.5,\y);}
\draw[kaxis] (-4.2,0)--(5.8,0) node[right]{$x$};
\draw[kaxis] (0,-5.7)--(0,5.3) node[above]{$y$};
\foreach \x in {-3,-2,-1,1,2,3,4,5}\node[below=2pt] at (\x,0){$\x$};
\foreach \y in {-5,-4,-3,-2,-1,1,2,3,4}\node[left=2pt] at (0,\y){$\y$};
\node[above left=3pt] at (0,0){$O$};

\end{tikzpicture}
\figcap{题图：在直角坐标网格中标出题设各点。}\label{question-figure-2-22}
\end{minipage}\end{center}
}

% Source example 2.23; question-only view.
\newcommand{\KQuestionAX}{%
\begin{center}\begin{minipage}{.97\linewidth}\centering
\begin{tikzpicture}[>=Latex,scale=.9]

\draw[kaxis] (0,0)--(6,0) node[right]{$r$};
\foreach \r in {1,...,5}{\draw (\r,0)--(\r,.1);\node[below] at (\r,0){$\r$};}
\node[left] at (0,0){$O$};
\path[use as bounding box] (-.5,-1.4) rectangle (6.4,1.4);

\end{tikzpicture}
\figcap{题图：以向右射线为极轴，逆时针为正角方向。}\label{question-figure-2-23}
\end{minipage}\end{center}
}

% Source example 2.28; question-only view.
\newcommand{\KQuestionBC}{%
\begin{center}\begin{minipage}{.97\linewidth}\centering
\begin{tikzpicture}[>=Latex,scale=.9]

\coordinate (A) at (3,0);\coordinate (M) at (0,3);\coordinate (N) at (6,3);
\coordinate (B) at (1.125,3);\coordinate (H) at (3,3);
\draw[kguide] (M)--(N) (A)--(H) (A)--(B);
\node[point,label=below:$A$] at (A){};
\node[left] at (M){$M$};\node[right] at (N){$N$};
\draw[vector] (B)--++(1.1,0) node[midway,above]{$v$};
\draw[vector] (A)--++(0,1.7) node[midway,right]{$u$};
\node[below left] at (B){$B$};\node[above right] at (H){$H$};
\draw[TextGray] (A)++(90:.8) arc(90:122:.8);\node at (2.64,1.13){$\theta$};
\draw[<->] (3.7,0)--(3.7,3) node[midway,right]{$L$};

\end{tikzpicture}
\figcap{题图：潜艇 A 与航线 MN 的垂直距离为 $L$，$\theta$ 从航线法线量起。}\label{question-figure-2-28}
\end{minipage}\end{center}
}

% Source example 2.29; question-only view.
\newcommand{\KQuestionBD}{%
\begin{center}\begin{minipage}{.97\linewidth}\centering
\begin{tikzpicture}[>=Latex,scale=.9]

\coordinate (A) at (0,0);\coordinate (B) at (4,2);\coordinate (X) at (6,0);
\draw[kaxis] (A)--(X) node[right]{$X$};\draw[thick] (A)--(B) node[midway,above left]{$d$};
\draw[vector] (A)--(1.3,0) node[midway,below]{$v$};
\draw[vector] (B)--++(.5,-1.1) node[right]{$v_p$};
\draw (0.8,0) arc(0:26.565:.8);\node at (1.12,.3){$\theta$};
\node[point,label=below left:$A$] at (A){};\node[point,label=above:$B$] at (B){};

\end{tikzpicture}
\figcap{题图：球沿 AX 运动，队员从 B 开始奔跑；奔跑方向待选。}\label{question-figure-2-29}
\end{minipage}\end{center}
}

% Source example 2.30; question-only view.
\newcommand{\KQuestionBE}{%
%
\begin{center}\begin{tikzpicture}[scale=.82,>=Latex]
  \coordinate (P) at (-3.7,-.75);\coordinate (Q) at (2.25,-.75);\coordinate (R) at (-.3,2.25);
  \draw[very thick] (P)--(Q);
  \draw[kguide] (Q)--(R);
  \draw[vector] (-3.25,-.75)--(-1.8,-.75) node[midway,clearlabel,above]{$\vect v_1$};
  \draw[vector] (1.93,-.38)--(.95,.78) node[midway,clearlabel,right]{$\vect v_2$};
  \draw pic[draw=TextGray,->,angle radius=8mm,"$\alpha$"{clearlabel}] {angle=R--Q--P};
  \draw[<->,TextGray] (P)--(Q) node[midway,clearlabel,below=3pt]{$l$};
  \node[point,label={[clearlabel]left:$A,P$}] at (P) {};
  \node[point,label={[clearlabel]right:$B,Q$}] at (Q) {};
  \node[clearlabel,above] at (R){$R$};
\end{tikzpicture}\figcap{题图：A、B 初始位于 P、Q，分别沿 PQ、QR 匀速运动。}\label{question-figure-2-30}\end{center}
}

% Source example 2.31; question-only view.
\newcommand{\KQuestionBF}{%
%
\begin{center}\begin{minipage}{.96\linewidth}\centering
\begin{tikzpicture}[>=Latex]
 \coordinate (A) at (-2.8,2.3);\coordinate (C) at (.2,0);\coordinate (B) at (1.7,-1.9);
 \coordinate (U) at (.2,1.4);\coordinate (D) at (.2,-1.4);
 \draw[thick] (-3.5,0)--(3.6,0) node[right]{岸线};
 \draw[kguide] (A)--(-2.8,0) (B)--(1.7,0) (U)--(D);
 \draw[vector] (A)--(C) node[pos=.35,above right=2pt]{$v_1$};
 \draw[vector] (C)--(B) node[pos=.45,above right=2pt]{$v_2$};
 \draw pic[draw=TextGray,angle radius=7mm] {angle=U--C--A};
 \draw pic[draw=TextGray,angle radius=7mm] {angle=D--C--B};
 \node at (-.24,1.03){$\theta_1$};\node at (.55,-1.1){$\theta_2$};
 \draw[<->,TextGray] (-3.3,0)--(-3.3,2.3) node[midway,left]{$d_1$};
 \draw[<->,TextGray] (2.15,0)--(2.15,-1.9) node[midway,right]{$d_2$};
 \draw[<->,TextGray] (-2.8,-2.5)--(.2,-2.5) node[midway,below]{$x$};
 \draw[<->,TextGray] (.2,-2.5)--(1.7,-2.5) node[midway,below]{$L-x$};
 \node[point,label=above left:$A$] at (A){};\node[point,label=below right:$B$] at (B){};
 \node[point,label=above right:$C$] at (C){};
\node[below left=3pt] at (-3.5,0){$M$};\node[above left=6pt] at (3.6,0){$N$};
\end{tikzpicture}
\figcap{题图：A 在岸上、B 在水中，C 为待选择的入水点；两垂足间距为 $L$。}\label{question-figure-2-31}
\end{minipage}\end{center}
}

% Source example 2.32; question-only view.
\newcommand{\KQuestionBG}{%
%
\begin{center}\begin{minipage}{.95\linewidth}\centering
\begin{tikzpicture}[>=Latex]
 \draw[thick] (-.15,0)--(8.2,0);\draw[thick] (4,0)--(4,1.2);
 \draw[<->,TextGray] (4.25,0)--(4.25,1.2) node[midway,right=2pt]{$h$};
 \draw[kguide] (0,0)--(0,3) (6,0)--(6,.3);
 \draw[kcurve,->,domain=0:5.6,samples=90] plot(\x,{3-3*\x*\x/31.36});
 \draw[vector] (0,3)--(1.3,3) node[midway,above]{$v_0$};
 \draw[<->,TextGray] (-.5,0)--(-.5,3) node[midway,left]{$H$};
 \draw[<->,TextGray] (0,-.75)--(4,-.75) node[midway,below]{$L$};
 \draw[<->,TextGray] (4,-.75)--(6,-.75) node[midway,below]{$l$};
 \node[above left] at (0,3){$A$};\node[below=3pt] at (4,0){球网};\node[below=3pt] at (6,0){发球线};
\draw[thick] (8,-.12)--(8,.25);
\draw[<->,TextGray] (0,-1.4)--(8,-1.4) node[midway,below]{$2L$};
\node[above] at (8,.25){底线};
\end{tikzpicture}
\figcap{题图：A 为水平击球点；全场长 $2L$，球网到对方发球线的距离为 $l$。}\label{question-figure-2-32}
\end{minipage}\end{center}
}

% Source example 2.33; question-only view.
\newcommand{\KQuestionBH}{%
%
\begin{center}\begin{minipage}{.95\linewidth}\centering
\begin{tikzpicture}[>=Latex]
 \coordinate (A) at (0,2);\coordinate (B) at (4,0);\coordinate (V) at (5.1,-1.1);\coordinate (X) at (5.3,0);
 \draw[thick] (-.2,0)--(4.35,0);\draw[kguide] (A)--(0,0) (B)--(X);
 \draw[kcurve,dashed,domain=0:4,samples=80] plot(\x,{2-\x*\x/8});
 \draw[vector] (A)--(1.3,2) node[midway,above]{$\vect v_0$};
 \draw[vector] (B)--(V) node[right=4pt]{$\vect v$};
 \draw pic[draw=TextGray,angle radius=6mm] {angle=V--B--X};
 \node[right] at (4.7,-.21){$45^\circ$};
 \draw[<->,TextGray] (-.5,0)--(-.5,2) node[midway,left]{$h$};
 \node[point] at (A){};\node[point] at (B){};
\end{tikzpicture}
\figcap{题图：从高度 $h$ 水平抛出，落地速度与水平线成 $45^\circ$。}\label{question-figure-2-33}
\end{minipage}\end{center}
}

% Source example 2.34; question-only view.
\newcommand{\KQuestionBI}{%
%
\begin{center}\begin{minipage}{.95\linewidth}\centering
\begin{tikzpicture}[scale=.95,>=Latex]
 \coordinate (A) at (0,0);\coordinate (B) at (7.2,0);\coordinate (V) at (52:1.5);
 \draw[thick] (-.4,0)--(7.6,0);
 \draw[kcurve,->,domain=0:7.2,samples=100] plot(\x,{tan(52)*\x*(1-\x/7.2)});
 \draw[vector] (A)--(V) node[above left]{$v_0$};
 \draw pic[draw=TextGray,angle radius=7mm] {angle=B--A--V};\node at (.97,.39){$\theta$};
 \draw[<->,TextGray] (0,-.55)--(7.2,-.55) node[midway,below]{$S$};
\end{tikzpicture}
\figcap{题图：初速度为 $v_0$、仰角为 $\theta$，抛出点与落点等高。}\label{question-figure-2-34}
\end{minipage}\end{center}
}

% Source example 2.35; question-only view.
\newcommand{\KQuestionBJ}{%
%
\begin{center}\begin{minipage}{.96\linewidth}\centering
\begin{tikzpicture}[scale=.9,>=Latex]
 \coordinate (A) at (0,0);\coordinate (B) at (6.4,0);
 \draw[thick] (-.3,0)--(6.8,0);
 \foreach \ang in {22,68}{\draw[kcurve,dashed,domain=0:6.4,samples=100] plot(\x,{tan(\ang)*\x*(1-\x/6.4)});
 \draw[vector] (A)--(\ang:1.3);}
 \draw[TextGray] (.7,0) arc[start angle=0,end angle=22,radius=.7];
 \draw[TextGray] (1.1,0) arc[start angle=0,end angle=68,radius=1.1];
 \node at (1.39,.95){$\theta_1$};\node at (1.55,.22){$\theta_2$};
 \node[above] at (3.2,4){高弧线};\node[above] at (3.2,.7){低弧线};
 \node[below] at (3.2,-.4){$\theta_1+\theta_2=\pi/2$，两次初速率相同};
\end{tikzpicture}
\figcap{题图：两次初速率相同，两个抛射角互余。}\label{question-figure-2-35}
\end{minipage}\end{center}
}

% Source example 2.36; question-only view.
\newcommand{\KQuestionBK}{%
\begin{center}\begin{minipage}{.97\linewidth}\centering
\begin{tikzpicture}[>=Latex,scale=.9]

\draw[thick] (-.3,0)--(6,0) (5,0)--(5,1.7);
\draw[vector] (0,0)--(52:1.5) node[above left]{$v_0$};
\draw (.7,0) arc(0:52:.7);\node at (1.03,.4){$\theta$};
\node[below left] at (0,0){$A$};
\draw[<->] (0,-.55)--(5,-.55) node[midway,below]{$L$};
\draw[<->] (5.5,0)--(5.5,1.7) node[midway,right]{$H$};

\end{tikzpicture}
\figcap{题图：A 到障碍物的水平距离为 $L$，障碍物高 $H$。}\label{question-figure-2-36}
\end{minipage}\end{center}
}

% Source example 2.37; question-only view.
\newcommand{\KQuestionBL}{%
%
\begin{center}\begin{minipage}{.96\linewidth}\centering
\begin{tikzpicture}
\begin{axis}[width=.90\linewidth,height=7cm,axis lines=middle,xlabel={$x$},ylabel={$y$},
 xmin=0,xmax=7.2,ymin=-1,ymax=3.7,ticks=none,clip=true]
 
 \foreach \ang in {20,35,50,65,78}{
  \addplot[Blue!75,thick,domain=0:6.6,samples=140]{x*tan(\ang)-x*x/(12*cos(\ang)^2)};
 }
\end{axis}
\end{tikzpicture}
\figcap{题图：从同一点以同一初速率、不同仰角抛出的轨迹族。}\label{question-figure-2-37}
\end{minipage}\end{center}
}

% Source example 2.39; question-only view.
\newcommand{\KQuestionBN}{%
%
\begin{center}\begin{minipage}{.96\linewidth}\centering
\begin{tikzpicture}[>=Latex]
 \coordinate (C) at (0,0);\coordinate (T) at (5,1.5);\coordinate (X) at (5.8,0);
 \draw[thick] (C)--(X) (C)--(5.6,1.68);
 \draw[kcurve,dashed,->,domain=0:5,samples=100] plot(\x,{tan(55)*\x-(5*tan(55)-1.5)*\x*\x/25});
 \draw[vector] (C)--(55:1.3) node[above left]{$v_0$};
 \draw[TextGray] (.75,0) arc[start angle=0,end angle={atan(.3)},radius=.75];
 \draw[TextGray] (1.3,0) arc[start angle=0,end angle=55,radius=1.3];
 \node at (.99,.10){$\alpha$};\node at (1.48,.75){$\beta$};
 \draw[kguide] (T)--(5,0);
 \draw[<->,TextGray] (0,-.55)--(5,-.55) node[midway,below]{$L$};
 \node[point,label=below left:$C$] at (C){};\node[point,label=above right:$T$] at (T){};
\end{tikzpicture}
\figcap{题图：坡脚 C、坡上目标 T；$\alpha$ 为坡角，$\beta$ 为发射仰角。}\label{question-figure-2-39}
\end{minipage}\end{center}
}

% Source example 2.40; question-only view.
\newcommand{\KQuestionBO}{%
%
\begin{center}\begin{minipage}{.96\linewidth}\centering
\begin{tikzpicture}[>=Latex]
 \draw[thick] (-.3,0)--(6.5,0) (5,0)--(5,2);
 \draw[kcurve,dashed,domain=0:5,samples=100] plot(\x,{2-2*(\x-5)^2/25});
 \draw[vector] (0,0)--({1.3/sqrt(1.64)},{1.04/sqrt(1.64)}) node[above left]{$v_0$};
 \draw[vector] (5,2)--(6.3,2) node[midway,above]{$v_T$};
 \draw[TextGray] (.7,0) arc[start angle=0,end angle={atan(.8)},radius=.7];
 \node at (1.02,.31){$\theta$};
 \draw[<->,TextGray] (0,-.55)--(5,-.55) node[midway,below]{$L$};
 \draw[<->,TextGray] (5.5,0)--(5.5,2) node[midway,right]{$H$};
 \node[below left] at (0,0){$A$};\node[above left] at (5,2){$T$};
\end{tikzpicture}
\figcap{题图：要求抛体在障碍顶端 T 处的速度平行于水平地面。}\label{question-figure-2-40}
\end{minipage}\end{center}
}

% Source example 2.41; question-only view.
\newcommand{\KQuestionBP}{%
%
\begin{center}\begin{minipage}{.96\linewidth}\centering
\begin{tikzpicture}[scale=.95,>=Latex]
 \pgfmathsetmacro{\thetaLaunch}{20+atan(1/(2*tan(20)))}
 \pgfmathsetmacro{\uu}{tan(\thetaLaunch)}
 \pgfmathsetmacro{\kk}{(\uu-tan(20))/4.5}
 \coordinate (A) at (0,0);\coordinate (P) at (4.5,{4.5*tan(20)});
 \coordinate (B) at (5.6,{5.6*tan(20)});\coordinate (V) at ($(P)+(-70:1.2)$);
 \draw[thick] (A)--(5.9,0) (A)--(B)--(5.6,0);
 \draw[kcurve,dashed,domain=0:4.5,samples=100] plot(\x,{\uu*\x-\kk*\x*\x});
 \draw[vector] (A)--(\thetaLaunch:1.3) node[above left]{$v_0$};
 \draw[vector] (P)--(V) node[right=3pt]{$v_P$};
 \draw pic[draw=TextGray,angle radius=4mm] {right angle=A--P--V};
 \draw[TextGray] (.7,0) arc[start angle=0,end angle=20,radius=.7];
 \draw[TextGray] (1.3,0) arc[start angle=0,end angle=\thetaLaunch,radius=1.3];
 \node at (.96,.14){$\alpha$};\node at (1.42,.9){$\theta$};
 \node[point,label=above right:$P$] at (P){};
\end{tikzpicture}
\figcap{题图：从斜面底端抛出，要求再次碰到斜面时速度垂直于斜面。}\label{question-figure-2-41}
\end{minipage}\end{center}
}

% Source example 2.43; question-only view.
\newcommand{\KQuestionBR}{%
%
\begin{center}\begin{minipage}{.98\linewidth}\centering
\begin{tikzpicture}[scale=1,>=Latex]
 \pgfmathsetmacro{\hh}{1/(1-.45^2)}\pgfmathsetmacro{\xa}{1.4-1.8/1.45}
 \draw[very thick] (-.1,3)--(2.3,3)--(2.3,2)--(4.1,2)--(4.1,1)--(5.9,1)--(5.9,0)--(7.8,0);
 \draw[kcurve,domain=0:1,samples=70] plot({\xa+(1.4-\xa)*\x},{3+\hh-\hh*\x*\x});
 \foreach \n in {0,1,2}{
  \draw[kcurve,domain=0:1,samples=90] plot({1.4+1.8*\n+1.8*\x},{3-\n+2*.45/(1-.45)*\x-(1+.45)/(1-.45)*\x*\x});
 }
 \foreach \n in {0,1,2,3}{\node[point,label=below:{$C_{\n}$}] at ({1.4+1.8*\n},{3-\n}){};}
 \draw[vector] (\xa,{3+\hh})--++(1.1,0) node[midway,above]{$v_0$};
 \node[above left] at (\xa,{3+\hh}){$A$};
 \draw[<->,TextGray] (-.45,3)--(-.45,{3+\hh}) node[midway,left]{$H$};
 \draw[<->,TextGray] (3.2,3.55)--(5,3.55) node[midway,above]{$d$};
 \draw[<->,TextGray] (7.4,0)--(7.4,1) node[midway,right]{$l$};
\end{tikzpicture}
\figcap{题图：A 在首个踏面上方 $H$；台阶宽 $d$、高 $l$，小球逐次反弹。}\label{question-figure-2-43}
\end{minipage}\end{center}
}

% Source example 2.45; question-only view.
\newcommand{\KQuestionBT}{%
\begin{center}\begin{minipage}{.97\linewidth}\centering
\begin{tikzpicture}[>=Latex,scale=.9]

\draw[thick] (-1.8,0)--(4.1,0);\draw[thick] (0,1.3) circle(1.3);
\coordinate (A) at (3.2,0);\coordinate (T) at (0,2.6);\coordinate (O) at (0,1.3);
\draw[kguide] (0,0)--(T);
\draw[kcurve,dashed,domain=0:3.2,samples=100] plot(\x,{2.6+.75*\x-(5/10.24)*\x*\x});
\pgfmathsetmacro{\launchangle}{180-atan(2.375)}
\draw[vector] (A)--++(\launchangle:1.3) node[above right]{$v_0$};
\draw[vector] (T)--++(-.8,-.6) node[above left]{$v_T$};
\draw[<->] (0,-.55)--(3.2,-.55) node[midway,below]{$d$};
\draw (O)--(-1.3,1.3) node[midway,above]{$R$};
\draw (A)++(\launchangle:.75) arc(\launchangle:180:.75);\node at (2.2,.38){$\theta$};
\node[point,label=below:$O$] at (O){};\node[point,label=above left:$T$] at (T){};
\node[point,label=below right:$A$] at (A){};

\end{tikzpicture}
\figcap{题图：A 为抛出点，T 为球顶；虚线仅示意一条可行轨迹，并非唯一轨迹。}\label{question-figure-2-45}
\end{minipage}\end{center}
}

% Source example 2.47; question-only view.
\newcommand{\KQuestionBV}{%
%
\begin{center}\RadialCollisionBase{\vect v_B}
\figcap{题图：A 沿半径为 $R$ 的圆周运动，B 沿 OP 匀速运动。}\label{question-figure-2-47}\end{center}
}

% Source example 2.48; question-only view.
\newcommand{\KQuestionBW}{%
%
\begin{center}\RadialCollisionBase{\vect a_B}
\figcap{题图：A 沿圆周匀速运动，B 从 O 由静止沿 OP 匀加速运动。}\label{question-figure-2-48}\end{center}
}

% Source example 2.49; question-only view.
\newcommand{\KQuestionBX}{%
%
\begin{center}\begin{tikzpicture}[scale=.8,>=Latex]
  \draw[very thick] (0,0) circle (2.25);
  \coordinate (O) at (0,0);\coordinate (P) at (90:2.25);\coordinate (Q) at (135:2.25);
  \draw[very thick] (O)--(P);\draw[kguide] (O)--(Q);
  \draw[DeepBlue,very thick,->] (125:2.5) arc[start angle=125,end angle=175,radius=2.5]
    node[midway,clearlabel,left]{$\omega_1$};
  \draw[Blue,very thick,->] (-2.5,-.55) arc[start angle=192,end angle=230,radius=2.5]
    node[midway,clearlabel,left]{$\omega_2$};
  \draw pic[draw=TextGray,->,angle radius=8mm,"$\theta$"{clearlabel}] {angle=P--O--Q};
  \node[point] at (Q) {};\node[left] at ($(Q)+(-.7,.32)$){$Q$};
  \node[point,label={[clearlabel]above:$P$}] at (P) {};
  \node[point,label={[clearlabel]below right:$O$}] at (O) {};
\end{tikzpicture}\figcap{题图：两质点从正上方 P 同时出发，逆时针运动；$\theta$ 标记再次相遇的位置。}\label{question-figure-2-49}\end{center}
}

% Source example 2.50; question-only view.
\newcommand{\KQuestionBY}{%
%
\begin{center}\begin{tikzpicture}[scale=.8,>=Latex]
  \draw[very thick] (0,0) circle (1.35);\draw[very thick] (0,0) circle (2.35);
  \coordinate (O) at (0,0);\coordinate (A) at (28:1.35);\coordinate (B) at (28:2.35);
  \draw[kguide] (O)--(B);
  \draw[vector] (A)--++(118:1.05) node[midway,clearlabel,right]{$\vect v_A$};
  \draw[vector] (B)--++(118:.82);
  \node at (2.12,2.17){$\vect v_B$};
  \draw[TextGray] (O)--(145:1.35) node[midway,clearlabel,above=3pt]{$R_1$};
  \draw[TextGray] (O)--(225:2.35);
  \node at (-1.75,-.65){$R_2$};
  \node[point] at (O) {};\node at (.25,-.30){$O$};
  \node[point,label={[clearlabel]right:$A$}] at (A) {};
  \node[point,label={[clearlabel]below right:$B$}] at (B) {};
\end{tikzpicture}\figcap{题图：A、B 分别在半径为 $R_1$、$R_2$ 的同心圆上同向运动。}\label{question-figure-2-50}\end{center}
}

% Source example 2.51; question-only view.
\newcommand{\KQuestionBZ}{%
%
\begin{center}\begin{minipage}{.96\linewidth}\centering
\begin{tikzpicture}[scale=1,>=Latex]
 \draw[TextGray,dashed] (3,0) arc[start angle=0,end angle=180,x radius=3,y radius=.75];
 \draw[thick] (-3,0) arc[start angle=180,end angle=360,x radius=3,y radius=.75];
 \coordinate (O) at (0,0);\coordinate (A) at (3,0);\coordinate (S) at (0,3);
 \draw[kguide] (O)--(A) (O)--(S);
 \draw[vector] (S)--(1.25,3) node[midway,above]{$v_0$};
 \draw[kcurve,dashed,->,domain=0:3,samples=90] plot(\x,{3-\x*\x/3});
 \draw[Blue,->,thick] (3.3,0) arc[start angle=0,end angle=52,x radius=3.3,y radius=.83];
 \node[above right] at (2.6,.7){$\omega$};
 \draw[<->,TextGray] (-.45,0)--(-.45,3) node[midway,left]{$H$};
 \node[below] at (1.5,0){$R$};
 \node[point,label=below left:$O$] at (O){};
 \node[point,label=right:$A$] at (A){};
 \node[point,label=above left:{$S\;(B\text{的初始位置})$}] at (S){};
\end{tikzpicture}
\figcap{题图：水平圆心 O 上方为 S；B 从 S 水平抛出，A 沿圆周运动。}\label{question-figure-2-51}
\end{minipage}\end{center}
}

% Source example 2.52; question-only view.
\newcommand{\KQuestionCA}{%
\begin{center}\begin{minipage}{.97\linewidth}\centering
\begin{tikzpicture}[>=Latex,scale=.9]

\draw[thick] (0,0) circle(1.5) (0,0) circle(2.6);
\coordinate (O) at (0,0);\coordinate (A) at (75:1.5);\coordinate (B) at (27:2.6);
\draw[kguide] (O)--(27:2.6) (O)--(A);\draw[thick] (A)--(B);
\pgfmathsetmacro{\lineangle}{atan2(1.5*sin(75)-2.6*sin(27),1.5*cos(75)-2.6*cos(27))}
\draw (B)++(\lineangle:.6) arc(\lineangle:207:.6);
\node at ($(B)+(190:.91)$){$\theta$};
\draw[->,DeepBlue] (80:1.78) arc(80:122:1.78) node[above left]{$\omega_1$};
\draw[->,DeepBlue] (5:2.87) arc(5:39:2.87) node[above right]{$\omega_2$};
\draw (O)--(150:1.5);\node at (-1.72,1.05){$R_1$};
\draw (O)--(230:2.6) node[pos=.7,below right]{$R_2$};
\node[point,label=below:$O$] at (O){};\node[point,label=above:$A$] at (A){};\node[point,label=right:$B$] at (B){};

\end{tikzpicture}
\figcap{题图：B 观察 A，视线 BA 与指向圆心的 BO 夹角为 $\theta$；图中未预设最大夹角位置。}\label{question-figure-2-52}
\end{minipage}\end{center}
}

% Source example 2.57; question-only view.
\newcommand{\KQuestionCF}{%
%
\begin{center}\begin{tikzpicture}[scale=.72,>=Latex]
 \draw[very thick] (-2.4,-2.4)--(-2.4,0) arc[start angle=180,end angle=0,radius=2.4]--(2.4,-2.4);
 \draw[very thick] (-1.4,-2.4)--(-1.4,0) arc[start angle=180,end angle=0,radius=1.4]--(1.4,-2.4);
 \draw[->,DeepBlue,very thick] (-1.4,-1.9)--(-1.4,-.85);
 \draw[->,Blue,very thick] (-2.4,-1.9)--(-2.4,-.85);
 \draw[<->,TextGray] (0,0)--(0,1.4) node[midway,clearlabel,right]{$R_1$};
 \draw[<->,TextGray] (0,0)--(135:2.4) node[pos=.85,clearlabel,above right=3pt]{$R_2$};
 \node[clearlabel,below] at (0,-2.75){内、外路线均为半圆弯道};
\end{tikzpicture}\figcap{题图：内、外半圆弯道半径分别为 $R_1,R_2$。}\label{question-figure-2-57}\end{center}
}

% Source example 2.59; question-only view.
\newcommand{\KQuestionCH}{%
%
\begin{center}\begin{minipage}{.96\linewidth}\centering
\begin{tikzpicture}[scale=.95,>=Latex]
 \draw[thick] (0,0) circle (2.0);\coordinate (P) at (35:2);
 \coordinate (V) at ($(P)+(125:1.4)$);\coordinate (A) at ($(P)+(170:1.4)$);
 \draw[kguide] (0,0)--(P) node[pos=.4,below right]{$R$};
 \draw[vector] (P)--(V) node[above right]{$\vect v$};
 \draw[vector] (P)--(A) node[above left]{$\vect a$};
 \draw pic[draw=TextGray,angle radius=7mm] {angle=V--P--A};
 % The circumference crosses this narrow wedge. Show its label separately,
 % retaining the physical vectors and the main angle arc without a mask.
 \begin{scope}[shift={(4.7,1.3)}]
 \draw[vector] (0,0)--(125:1.6) node[above right]{$\vect v$};
 \draw[vector] (0,0)--(170:1.6) node[above left]{$\vect a$};
 \draw[TextGray] (125:.6) arc(125:170:.6);
 \node at (147.5:.95){$\theta$};
 \node[below] at (-.5,-.4){夹角放大};
 \end{scope}
 \node[point,label=below left:$O$] at (0,0){};\node[point] at (P){};
\end{tikzpicture}
\figcap{题图：速度 $\vect v$ 与加速度 $\vect a$ 的夹角为 $\theta$。}\label{question-figure-2-59}
\end{minipage}\end{center}
}

% Source example 2.60; question-only view.
\newcommand{\KQuestionCI}{%
\begin{center}\begin{minipage}{.97\linewidth}\centering
\begin{tikzpicture}[>=Latex,scale=.9]

\draw[kaxis] (-3,0)--(3.2,0) node[right]{$x$};
\draw[kaxis] (0,-1.8)--(0,2) node[above]{$y$};\node[below left] at (0,0){$O$};

\end{tikzpicture}
\figcap{题图：根据给定参数方程在空坐标系中画出轨迹。}\label{question-figure-2-60}
\end{minipage}\end{center}
}

% Source example 2.61; question-only view.
\newcommand{\KQuestionCJ}{%
\begin{center}\begin{minipage}{.97\linewidth}\centering
\begin{tikzpicture}[>=Latex,scale=.9]

\draw[kaxis] (-3,0)--(3.2,0) node[right]{$x$};
\draw[kaxis] (0,-1.8)--(0,2) node[above]{$y$};\node[below left] at (0,0){$O$};

\end{tikzpicture}
\figcap{题图：根据给定参数方程在空坐标系中画出轨迹。}\label{question-figure-2-61}
\end{minipage}\end{center}
}

% Source example 2.62; question-only view.
\newcommand{\KQuestionCK}{%
%
\begin{center}\begin{tikzpicture}[scale=.72,>=Latex]
 \draw[very thick] (-.4,0)--(10,0);
 \draw[kaxis] (0,0)--(10.2,0) node[right,clearlabel]{$x$};\draw[kaxis] (0,0)--(0,3.4) node[above,clearlabel]{$y$};
 \draw[kcurve,domain=0:6.283,samples=120] plot ({1.35*(\x-sin(\x r))},{1.35*(1-cos(\x r))});
 \draw[very thick] (0,1.35) circle (1.35);\node[point,label={[clearlabel]below left:$P$}] at (0,0) {};
 \draw[vector] (0,1.35)--(1.1,1.35) node[pos=.22,clearlabel,above=3pt]{$v$};
\end{tikzpicture}\figcap{题图：圆沿水平地面无滑动滚动，P 为圆周上的标记点，初始位于接触处。}\label{question-figure-2-62}\end{center}
}

% Source example 2.63; question-only view.
\newcommand{\KQuestionCL}{%
\begin{center}\begin{minipage}{.97\linewidth}\centering
\begin{tikzpicture}[>=Latex,scale=.85,x={(1cm,0cm)},y={(-.50cm,-.28cm)},z={(0cm,.80cm)}]

\draw[kaxis] (0,0,0)--(3,0,0) node[right]{$x$};
\draw[kaxis] (0,0,0)--(0,2.8,0) node[below left]{$y$};
\draw[kaxis] (0,0,0)--(0,0,3.5) node[above]{$z$};
\node[below right] at (0,0,0){$O$};

\end{tikzpicture}
\figcap{题图：在右手直角坐标系中，根据三个坐标随时间的变化画出轨迹。}\label{question-figure-2-63}
\end{minipage}\end{center}
}

% Source example 2.64; question-only view.
\newcommand{\KQuestionCM}{%
%
\begin{center}\begin{tikzpicture}[scale=.78,>=Latex]
  \draw[kaxis] (-3.1,0)--(3.25,0) node[right,clearlabel]{$x$};
  \draw[kaxis] (0,-3.0)--(0,3.15) node[above,clearlabel]{$y$};
  \draw[very thick] (0,0) circle (2.25);
  \coordinate (O) at (0,0);\coordinate (P) at (38:2.25);
  \draw[very thick] (O)--(P) node[midway,clearlabel,below right]{$R$};
  \draw[kguide] (P)--(P|-O) node[clearlabel,below]{$x$};
  \draw[kguide] (P)--(P-|O) node[clearlabel,left]{$y$};
  \node[point,label={[clearlabel]below left:$O$}] at (O) {};
  \node[point,label={[clearlabel]above right:$P$}] at (P) {};
\end{tikzpicture}\figcap{题图：圆周上 P 点在两个坐标轴上的投影分别为 $x,y$。}\label{question-figure-2-64}\end{center}
}

% Source example 2.65; question-only view.
\newcommand{\KQuestionCN}{%
%
\begin{center}\begin{tikzpicture}[scale=.8,>=Latex]
  \draw[very thick] (-3.7,0)--(3.7,0);
  \draw[very thick] (0,1.55) circle (1.55);
  \coordinate (O) at (0,1.55);\coordinate (P) at (0,0);
  \draw[vector] (O)--++(1.45,0) node[midway,clearlabel,above]{$\vect v_O$};
  \draw[DeepBlue,very thick,->] ({1.85*cos(30)},{1.55+1.85*sin(30)}) arc[start angle=30,end angle=-20,radius=1.85];
  \node[clearlabel,right] at (2.05,1.15){$\omega$};
  \draw[very thick] (O)--(P) node[midway,clearlabel,right]{$R$};
  \node[point] at (P) {};\node[below=4pt] at (P){$P$};
  \node[point,label={[clearlabel]above left:$O$}] at (O) {};
\end{tikzpicture}\figcap{题图：圆无滑动向右滚动，P 为此刻与地面接触的圆周点。}\label{question-figure-2-65}\end{center}
}

% Source example 2.66; question-only view.
\newcommand{\KQuestionCO}{%
\begin{center}\begin{minipage}{.97\linewidth}\centering
\begin{tikzpicture}[>=Latex,scale=.9]

\draw[kaxis] (0,0)--(6,0) node[right]{$r$};\node[below left] at (0,0){$O$};
\path[use as bounding box] (-.5,-1.35) rectangle (6.4,1.35);

\end{tikzpicture}
\figcap{题图：以 O 为极点、向右为极轴，描出给定极坐标运动。}\label{question-figure-2-66}
\end{minipage}\end{center}
}

% Source example 2.68; question-only view.
\newcommand{\KQuestionCQ}{%
\begin{center}\begin{minipage}{.97\linewidth}\centering
\begin{tikzpicture}[>=Latex,scale=.9]

\draw[kaxis] (-2.2,0)--(2.3,0) node[right]{$x$};
\draw[kaxis] (0,-.7)--(0,2.7) node[above]{$y$};
\draw[kcurve,domain=-1.8:1.8,samples=100] plot(\x,{.65*\x*\x});
\node[below left] at (0,0){$O$};\node[right] at (1.8,2.1){$y=ax^2$};

\end{tikzpicture}
\figcap{题图：确定抛物线在原点处的曲率半径；这里示意 $a>0$。}\label{question-figure-2-68}
\end{minipage}\end{center}
}

% Source example 2.69; question-only view.
\newcommand{\KQuestionCR}{%
\begin{center}\begin{minipage}{.97\linewidth}\centering
\begin{tikzpicture}[>=Latex,scale=1.15]
\draw[kcurve] \FoldedCurvePath;\FoldedCurvePoints
\end{tikzpicture}
\figcap{题图：比较 $P_1,P_2,P_3,P_4$ 的局部弯曲，注意 $P_3$ 后向左下回折。}\label{question-figure-2-69}
\end{minipage}\end{center}
}

% Source example 2.70; question-only view.
\newcommand{\KQuestionCS}{%
\begin{center}\begin{minipage}{.97\linewidth}\centering
\begin{tikzpicture}[>=Latex,scale=.9]

\draw[thick] (0,0) ellipse(3 and 1.6);\draw (0,0)--(3,0) node[midway,below]{$a$};
\draw (0,0)--(0,1.6) node[midway,left]{$b$};\node[below left] at (0,0){$O$};

\end{tikzpicture}
\figcap{题图：椭圆的半长轴为 $a$、半短轴为 $b$。}\label{question-figure-2-70}
\end{minipage}\end{center}
}

% Source example 2.73; question-only view.
\newcommand{\KQuestionCV}{%
%
\begin{center}\begin{tikzpicture}[scale=.78,>=Latex]
 \fill[Blue!7] (-3,-1.45) rectangle (3,1.45);
 \draw[very thick] (-3,-1.45)--(3,-1.45) (-3,1.45)--(3,1.45);
 \coordinate (A) at (0,-1.45);\coordinate (B) at (0,1.45);
 \node[point,label={[clearlabel]below:$A$}] at (A) {};
 \node[point,label={[clearlabel]above:$B$}] at (B) {};
 \draw[vector] (-2.35,.15)--(-1.05,.15) node[midway,clearlabel,above]{$\vect u$};
 \draw[vector] (A)--++(145:1.25) node[near end,clearlabel,below left]{$\vect v$};
 \draw[<->,TextGray] (2.45,-1.45)--(2.45,1.45) node[midway,clearlabel,right]{$L$};
\end{tikzpicture}\figcap{题图：两岸正对的 A、B 距离为 $L$，水流速度向右，船速指船相对水的速度。}\label{question-figure-2-73}\end{center}
}

% Source example 2.74; question-only view.
\newcommand{\KQuestionCW}{%
%
\begin{center}\begin{tikzpicture}[scale=.82,>=Latex]
 \draw[very thick] (-3,0)--(3,0);
 \coordinate (A) at (-2.1,0);\coordinate (B) at (2.1,0);
 \node[point,label={[clearlabel]below:$A$}] at (A) {};
 \node[point,label={[clearlabel]below:$B$}] at (B) {};
 \draw[vector] (A)--++(1.55,0) node[midway,clearlabel,above]{$v$};
 \draw[vector] (-.75,1.0)--(.75,1.0) node[midway,clearlabel,above]{$S'\text{ 的速度 }u$};
\end{tikzpicture}\figcap{题图：地面系中 A 向右运动、B 静止；参考系 $S\prime$ 以速度 $u$ 向右平动。}\label{question-figure-2-74}\end{center}
}

% Source example 2.82; question-only view.
\newcommand{\KQuestionDE}{%
\begin{center}\begin{minipage}{.97\linewidth}\centering
\begin{tikzpicture}[>=Latex,scale=.9]

\draw[thick] (-2,0)--(3,0) (-2,2.7)--(3,2.7);
\node[point,label=below:$A$] at (0,0){};
\draw[vector] (-1.7,1.5)--(-.3,1.5) node[midway,above]{$v_2$};
\draw[vector] (0,0)--(110:1.4) node[pos=.5,right=7pt]{$v_1$};
\draw[<->] (2.4,0)--(2.4,2.7) node[midway,left]{$d$};

\end{tikzpicture}
\figcap{题图：水流平行于河岸，小船对水速度的大小为 $v_1$，船头方向可以选择。}\label{question-figure-2-82}
\end{minipage}\end{center}
}

% Source example 2.85; question-only view.
\newcommand{\KQuestionDH}{%
\begin{center}\begin{minipage}{.97\linewidth}\centering
\begin{tikzpicture}[>=Latex,scale=.9]

\coordinate (A) at (4,0);\coordinate (B) at (-2,2.4);\coordinate (C) at (2.4,2.4);
\draw[thick] (A)--(B) node[midway,below left]{$L$};
\draw[kguide] (B)--(5.1,2.4) (A)--(4,3.2) (A)--(C);
\draw[vector] (B)--(-.7,2.4) node[midway,above]{$v_1$};
\draw[vector] (A)--($(A)!.43!(C)$) node[left]{$v_2$};
\draw (B)++(-21.801:.85) arc(-21.801:0:.85);\node at (-.7,2.1){$\alpha$};
\draw (A)++(123.690:.65) arc(123.690:158.199:.65);\node at (3.23,.69){$\beta$};
\foreach \p/\pos in {A/below right,B/above left,C/above}\node[point,label=\pos:$\p$] at (\p){};

\end{tikzpicture}
\figcap{题图：A 为潜艇、B 为敌船初始位置，C 表示可能的相遇点；发射方向待定。}\label{question-figure-2-85}
\end{minipage}\end{center}
}

% Source example 2.87; question-only view.
\newcommand{\KQuestionDJ}{%
\begin{center}\begin{minipage}{.97\linewidth}\centering
\begin{tikzpicture}[>=Latex,scale=.9]

\draw[kguide] (0,0)--(4.8,0);
\foreach \x/\lab in {0/A,1.8/B,3.6/C}{\fill (\x,0) circle(1.5pt);\node[below] at (\x,0){$\lab$};}
\draw[vector] (0,0)--(.6,0);\node[above] at (.3,0){$u$};
\draw[vector] (1.8,0)--(1.8,1.2) node[above]{$2u$};
\draw[vector] (3.6,0)--(4.8,1.2) node[above]{$2\sqrt2u$};
\draw[<->] (0,-.65)--(1.8,-.65) node[midway,below]{$a$};
\draw[<->] (1.8,-.65)--(3.6,-.65) node[midway,below]{$a$};
\end{tikzpicture}
\figcap{题图：三船初始位置及对地速度，图中 A、B、C 等间距。}\label{question-figure-2-87}
\end{minipage}\end{center}
}

% Source example 2.89; question-only view.
\newcommand{\KQuestionDL}{%
\begin{center}\begin{minipage}{.97\linewidth}\centering
\begin{tikzpicture}[>=Latex,scale=.9]

\coordinate (O) at (0,0);\coordinate (A) at (0,3.6);\coordinate (B) at (4,0);\coordinate (P) at (1.52,2.232);
\draw[thick] (0,4.1)--(O)--(4.9,0);\draw[kcurve] (A)--(B);
\foreach \p/\pos in {A/left,B/below,P/above right}{\fill (\p) circle(1.5pt);\node[\pos] at (\p){$\p$};}
\draw[vector] (A)--++(0,-.8) node[midway,left]{$v_A$};\draw[vector] (B)--++(.8,0) node[above]{$v_B$};
\draw (0,2.9) arc(-90:-41.987:.7);\node at (.42,2.7){$\theta$};
\node[above right] at (.75,3.25){$aL$};\node[above right] at (2.7,1.4){$(1-a)L$};
\end{tikzpicture}
\figcap{题图：杆的 A、B 端接触墙、地；$\theta$ 从竖直墙量起，$AP=aL$。}\label{question-figure-2-89}
\end{minipage}\end{center}
}

% Source example 2.90; question-only view.
\newcommand{\KQuestionDM}{%
%
\begin{center}\begin{minipage}{\linewidth}\centering\begin{tikzpicture}[>=Latex,scale=.93]
\coordinate (O) at (0,2.5);\coordinate (B) at (4,0);
\draw[thick] (-3.8,2.3)--(-.15,2.3)--(-.15,-.4);
\fill[Blue!5] (0,-.7) rectangle (5.3,0);\draw[thin,Blue] (0,0)--(5.3,0);
\draw[thick] (-3.2,2.5) rectangle (-2.2,3);\draw (-3,2.4) circle(.1) (-2.4,2.4) circle(.1);
\draw[kcurve] (-2.2,2.5)--(O)--(B);
\fill (O) circle(1.8pt);\node[above right] at (O){$O$};
\draw[thick] (B)--(5.1,0)--(4.9,-.4)--(4.2,-.4)--cycle;
\fill (B) circle(1.5pt);\node[above right] at (B){$B$};
\draw[vector] (-2.3,3.45)--(-3.4,3.45) node[midway,above]{$v$};
\draw[vector] (4.9,-1)--(3.9,-1) node[midway,below]{$u$};
\node[above] at (-1.2,2.5){$l_c$};\node[above right] at (1.9,1.32){$l$};
\draw[<->] (-4.25,0)--(-4.25,2.5) node[midway,left]{$h$};
\draw[kguide] (-4.25,0)--(0,0) (-4.25,2.5)--(-3.5,2.5);
\draw (B)++(148:.9) arc(148:180:.9);\node at (2.82,.31){$\theta$};
\draw[<->] (0,-.85)--(4,-.85) node[midway,below]{$x$};
\end{tikzpicture}\figcap{题图：汽车经岸上定滑轮拉船；绳不可伸长，$h$ 为滑轮相对水面的高度。}\label{question-figure-2-90}\end{minipage}\end{center}
}

% Source example 2.91; question-only view.
\newcommand{\KQuestionDN}{%
\begin{center}\begin{minipage}{.97\linewidth}\centering
\begin{tikzpicture}[>=Latex,scale=.9]

\draw[thick] (0,0)--(4,0)--(0,2)--cycle;
\coordinate (P) at (2.2,.9);\draw[very thick] (P)--(2.2,3.1);
\foreach \x in {2.06,2.34}{\draw[thick] (\x,2)--(\x,2.6);}
\draw (1.6,2)--(2.06,2) (2.34,2)--(2.8,2);
\fill (P) circle(1.5pt);\node[above right=3pt] at (P){$P$};
\draw[vector] (.6,-.5)--(1.7,-.5) node[midway,below]{$v_0$};
\draw[vector] (2.7,2.65)--(2.7,3.5) node[right]{$v$};
\draw (3.2,0) arc(180:153.435:.8);\node at (2.95,.24){$\theta$};
\end{tikzpicture}
\figcap{题图：斜面向右运动，竖杆位于两侧固定导向之间，杆底 P 与斜面保持接触。}\label{question-figure-2-91}
\end{minipage}\end{center}
}

% Source example 2.92; question-only view.
\newcommand{\KQuestionDO}{%
\begin{center}\begin{minipage}{.97\linewidth}\centering
\begin{tikzpicture}[>=Latex,scale=.9]

\coordinate (C) at (0,0);\coordinate (P) at (1.2,1.6);
\draw[thick] (-2,0)--(2,0) arc(0:180:2);
\draw[kguide] (C)--(0,2.2) (C)--(P) (P)--(1.2,0);
\draw[very thick] (P)--(1.2,3.2);
\foreach \x in {1.06,1.34}{\draw[thick] (\x,2.4)--(\x,3);}
\draw (.6,2.4)--(1.06,2.4) (1.34,2.4)--(1.8,2.4);
\fill (P) circle(1.5pt) (C) circle(1.5pt);
\node[below left] at (C){$C=(X,0)$};\node[above right=3pt] at (P){$P$};
\node at (.7,.46){$R$};\node[below] at (1.2,0){$x_0$};
\draw (0,.75) arc(90:53.1301:.75);\node at (.29,1.13){$\theta$};
\draw[vector] (.5,-.8)--(1.6,-.8) node[midway,below]{$v_0$};
\end{tikzpicture}
\figcap{题图：半圆柱向右运动；竖杆横坐标固定，P 为接触点，$\theta$ 从竖直方向量起。}\label{question-figure-2-92}
\end{minipage}\end{center}
}

% Source example 2.93; question-only view.
\newcommand{\KQuestionDP}{%
%
\begin{center}\begin{minipage}{\linewidth}\centering\begin{tikzpicture}[>=Latex,scale=1]
\coordinate (O) at (0,0);\coordinate (B) at (0,3.4);\coordinate (A) at (2.3,1.2);\coordinate (C) at (-2.2,3.4);
\draw[very thick] (O)--(A) node[midway,below=5pt]{$R$};
\draw[kcurve] (A)--(B)--(C)--(-2.2,.65);
\draw[thick] (-2.55,0) rectangle (-1.85,.65);\node at (-2.2,.32){$M$};
\draw[kguide] (O)--(B);
\foreach \p/\pos in {O/below,B/above,C/above,A/right}{\fill (\p) circle(1.7pt);\node[\pos] at (\p){$\p$};}
\draw[<->] (-.65,0)--(-.65,3.4) node[midway,left]{$H$};
\draw (B)++(-90:.8) arc(-90:-43.727:.8);\node at (.42,2.4){$\alpha$};
\draw[->,DeepBlue] (35:.75) arc(35:120:.75);\node[above] at (.3,1.2){$\omega$};
\node[left] at (.83,2.0){$l$};
\end{tikzpicture}\figcap{题图：杆 OA 绕 O 转动，绳从 A 经固定转向点 B、C 连到物块 M。}\label{question-figure-2-93}\end{minipage}\end{center}
}

% Source example 2.94; question-only view.
\newcommand{\KQuestionDQ}{%
\begin{center}\begin{minipage}{.97\linewidth}\centering
\begin{tikzpicture}[>=Latex,scale=.9]

\draw[thick] (0,0)--(4,0);\node[point,label=left:$A$] at (0,0){};\node[point,label=right:$B$] at (4,0){};
\draw[vector] (0,0)--(0,1.6) node[above]{$v_A$};\draw[vector] (4,0)--(4,1.6) node[above]{$v_B$};
\draw[<->] (0,-.55)--(4,-.55) node[midway,below]{$l$};

\end{tikzpicture}
\figcap{题图：A、B 用轻杆连接；图示规定竖直向上为端点速度的正方向，各问按给定符号取向。}\label{question-figure-2-94}
\end{minipage}\end{center}
}

% Source example 2.95; question-only view.
\newcommand{\KQuestionDR}{%
%
\begin{center}\begin{minipage}{\linewidth}\centering\begin{tikzpicture}[>=Latex,scale=1]
\draw[thick] (-2,0)--(4,0);\draw[thick] (0,1) circle(1);
\fill (0,1) circle(1.5pt) (0,0) circle(1.5pt);\node[below left] at (0,1){$C$};\node[below left] at (0,0){$P$};
\draw (0,1)--(-.707,1.707);\node[left] at (-.58,1.28){$R$};
\draw[vector] (0,1)--(2.6,1);\node[above] at (2,1){$v_0$};
\draw[->,DeepBlue,thick] (40:1.5)++(0,1) arc(40:145:1.5);
\node[above] at (0,2.65){$\omega_0>0$（逆时针）};
\end{tikzpicture}\figcap{题图：球初始向右平动并逆时针转动，地面为水平面。}\label{question-figure-2-95}\end{minipage}\end{center}
}

% Source example 2.96; question-only view.
\newcommand{\KQuestionDS}{%
%
\begin{center}\begin{minipage}{\linewidth}\centering\begin{tikzpicture}[>=Latex]
\begin{scope}[x=2.1cm,y=2.1cm]
\draw[thick] (5,-.1)--(5,1.1);\node[right] at (5,.7){木靶};
\draw[kguide,domain=0:1,samples=100,variable=\t] plot({4.9*\t},{4.9*tan(30)*\t*(1-\t)});
\foreach \t/\n in {0/1,.25/2,.5/3,.75/4,1/5}{
\begin{scope}[shift={({4.9*\t},{4.9*tan(30)*\t*(1-\t)})},rotate={60+300*\t}]
\KnifeOutline\fill (0,0) circle(.7pt);
\end{scope}
\node[above] at ({4.9*\t},{4.9*tan(30)*\t*(1-\t)+.17}){\n};}
\draw[<->] (0,-.35)--(5,-.35) node[midway,below]{$s=5.0\,\mathrm m$};
\node at (2.5,-.8){(a) 质心抛物线与五个等时姿态};
\end{scope}
\begin{scope}[xshift=2.1cm,yshift=-3.8cm]
\begin{scope}[scale=15]\KnifeOutline\end{scope}
\fill (0,0) circle(1.5pt);\node[below] at (0,-.25){$C$};\node[right] at (1.5,0){刀尖};
\draw[<->] (-2.25,-.65)--(1.5,-.65) node[midway,below]{$L=25\,\mathrm{cm}$};
\draw[<->] (0,.65)--(1.5,.65) node[midway,above]{$l=10\,\mathrm{cm}$};
\node at (0,-1.45){(b) 长度定义（局部放大）};
\end{scope}
\begin{scope}[xshift=8cm,yshift=-3.8cm]
\draw[kaxis] (-.3,0)--(2.5,0) node[right]{$x$};
\begin{scope}[rotate=60,scale=10]\KnifeOutline\end{scope}
\draw[vector] (0,0)--(30:2.1) node[above right]{$\vect v_C$};
\draw (.65,0) arc(0:30:.65);\node at (15:.95){$30^\circ$};
\node[right] at (.25,-.45){$C$};\draw[thin] (.25,-.45)--(0,0);
\node[above] at (60:2.05){$\theta=60^\circ$};
\node at (.3,-2.45){(c) 初姿态与速度方向};
\end{scope}
\end{tikzpicture}\figcap{题图：飞刀的飞行姿态、长度及初始方向；木靶与初始质心等高。}\label{question-figure-2-96}\end{minipage}\end{center}
}

% Source example 2.103; question-only view.
\newcommand{\KQuestionDZ}{%
\begin{center}\begin{minipage}{.97\linewidth}\centering
\begin{tikzpicture}[>=Latex,scale=.9]

\coordinate (E) at (0,0);\coordinate (P) at (0,3.3);
\draw[thick] (E)--(P) node[midway,right]{$L$};
\draw (E) circle(.25);\node[below] at (0,-.3){地球 E};
\draw[vector] (P)--++(-135:1.5) node[below left]{$v$};
\draw (P)++(-135:.7) arc(-135:-90:.7);\node at (-.4,2.34){$\theta$};
\node[point,label=above:$P$] at (P){};

\end{tikzpicture}
\figcap{题图：天体初始距地球 $L$，运动方向与指向地球的视线成角 $\theta$。}\label{question-figure-2-103}
\end{minipage}\end{center}
}

% Source example 2.105; question-only view.
\newcommand{\KQuestionEB}{%
%
\begin{center}\begin{minipage}{\linewidth}\centering\begin{tikzpicture}[>=Latex,scale=.9]
\foreach \t/\n/\dx/\dy in {-90/1/0/0,-60/2/3.5/0,-36.37/3/7/0,0/4/10.5/0,36.37/5/1.75/-4,60/6/5.25/-4,90/7/8.75/-4}{
\begin{scope}[shift={(\dx,\dy)}]
\draw[TextGray] (0,0) circle(.8);\draw[thick] (0,0) circle({.8*1.524});
\fill (0,0) circle(1.3pt);\node[above] at (0,0){$S$};
\pgfmathsetmacro{\ea}{.985609*\t-90}
\pgfmathsetmacro{\ma}{.985609*\t*pow(1.524,-1.5)-90}
\pgfmathsetmacro{\ex}{.8*cos(\ea)}\pgfmathsetmacro{\ey}{.8*sin(\ea)}
\pgfmathsetmacro{\mx}{.8*1.524*cos(\ma)}\pgfmathsetmacro{\my}{.8*1.524*sin(\ma)}
\draw[DeepBlue,dashed,thin,->] (\ex,\ey)--({\ex+1.7*(\mx-\ex)},{\ey+1.7*(\my-\ey)});
\fill (\ex,\ey) circle(1.5pt) (\mx,\my) circle(1.5pt);
\node[font=\scriptsize] at (\ea:.5){$E_{\n}$};
\ifdim\t pt<0pt \node[font=\scriptsize,left] at ({\mx-.13},{\my-.16}){$M_{\n}$};
\else \node[font=\scriptsize,right] at ({\mx+.13},{\my-.16}){$M_{\n}$};\fi
\node[font=\footnotesize] at (0,1.7){\n：$t=\t\,\mathrm d$};
\end{scope}}
\node at (5.25,-6){虚线由地球经过火星，指向背景星空；各小图比例相同};
\end{tikzpicture}\figcap{题图：七个时刻的日、地、火位置与观测视线；E 为地球、M 为火星、S 为太阳。}\label{question-figure-2-105}\end{minipage}\end{center}
%
\begin{center}\begin{minipage}{\linewidth}\centering\begin{tikzpicture}[>=Latex]
\draw[kaxis] (-4.4,-2.1)--(4.6,-2.1) node[right]{视经度增加};
\draw[kaxis] (-4.4,-2.1)--(-4.4,3.05) node[above]{视纬度增加};
\MarsInclinedPath
\node[above] at (0,3.25){三维圆轨道示例：轨道倾角 $1.85^\circ$};
\node at (0,-2.8){箭头表示观测的时间先后};
\end{tikzpicture}\figcap{同一现象在背景星空上的视位置回环示意。}\end{minipage}\end{center}
}

